\documentclass[10pt,a4paper]{article}

% ----------------- Paquetes -------------------%
\usepackage[T1]{fontenc}
\usepackage[utf8]{inputenc}
\usepackage[a4paper,margin=2.5cm]{geometry}
\usepackage{mathptmx}             
\usepackage{changepage} 
\usepackage{setspace}
\usepackage[spanish]{babel}
\usepackage[labelfont=bf,labelsep=space]{caption}
\usepackage{amssymb}
\usepackage{amsmath}
\usepackage{ebgaramond}
\usepackage{placeins}
\usepackage{etoolbox}
\usepackage{setspace}
\usepackage{parskip} 
\usepackage{tcolorbox}
\usepackage{needspace}
\usepackage{xparse}
\usepackage{ocgx2}
\usepackage{hyperref}
\usepackage{hypcap}
\usepackage{multirow}
\usepackage{soul}\sethlcolor{yellow}% resaltado amarillo: \hl{texto} (Panel TCEE, ⌃H)


%%%%%%%%%%%%%%%%%%%%%%%%%%%%%%%%%%%%%%%%%%%%%%%%%%%%%%%%%%%%%%%%
% --- Fija primero tus valores globales "buenos" ---
\setstretch{1.2}                 % Interlineado global
\setlength{\parskip}{0.7\baselineskip}
\setlength{\parindent}{0pt}
\newlength{\GlobalParskip}
\setlength{\GlobalParskip}{\parskip}

\newcounter{eqnum}[subsection]
\renewcommand{\theeqnum}{\arabic{eqnum}}
\newcounter{alphaitem}
\newcounter{numitem}

%% ------------------- Recuadros----------------------%%
\newcommand{\puntosclave}[1]{%
  \begin{tcolorbox}[
    colframe=black,
    title={Puntos clave},
    before upper={\setlength{\parindent}{0.5cm}\setlength{\parskip}{0.5\baselineskip}}
  ]
  #1
  \end{tcolorbox}
}

\newcommand{\modificaciones}[1]{
  \begin{tcolorbox}[
    colback=white,
    colframe=red,
    title={Modificaciones a realizar},
    before upper={\setlength{\parindent}{0.5cm}\setlength{\parskip}{0.5\baselineskip}}
  ]
  #1
  \end{tcolorbox}
}

%% ------------------- Preguntas TEST----------------------%%
% Opciones: radio (single-choice)
\NewDocumentCommand{\OptRadio}{m m m}{%
  \noindent\CheckBox[radio,name=#1,value=#2,width=1.2ex,height=1.2ex]{}%
  \;\textbf{#2.}~#3\par
}

% Opciones: checkbox (multi-choice)
\NewDocumentCommand{\OptCheck}{m m}{%
  \noindent\CheckBox[name=#1,width=1.2ex,height=1.2ex,bordercolor={0 0 0}]{}%
  \;#2\par
}

% Pregunta single-choice (radio)
% Args: 1=id, 2=enunciado, 3=opciones, 4=solución+justificación, 5=imagen (o {}), 6=origen
\NewDocumentCommand{\PreguntaTestSingle}{m m m m m m}{%
  \par\medskip
  \noindent\textbf{#1.}~#2\par\smallskip

    % Imagen (si no es {})
    \IfBlankTF{#5}{}{%
      \begin{center}
        \includegraphics[width=0.75\linewidth]{#5}
      \end{center}
    }

    \begin{Form}
      #3
    \end{Form}

    \par\smallskip
    \switchocg{sol-#1}{\fbox{Mostrar / ocultar solución}}%
    \hfill{\footnotesize #6}\par\smallskip

    \begin{ocg}{Solución}{sol-#1}{0}
      \noindent #4
    \end{ocg}
  \par\medskip
}

% Pregunta multi-choice (checkboxes)
\NewDocumentCommand{\PreguntaTestMulti}{m m m m m m}{%
  \par\medskip
  \noindent\textbf{#1.}~#2\par\smallskip

    % Imagen (si no es {})
    \IfBlankTF{#5}{}{%
      \begin{center}
        \includegraphics[width=0.75\linewidth]{#5}
      \end{center}
    }
    \begin{Form}
      #3
    \end{Form}

    \par\smallskip
    \switchocg{sol-#1}{\fbox{Mostrar / ocultar solución}}%
    \hfill{\footnotesize #6}\par\smallskip

    \begin{ocg}{Solución}{sol-#1}{0}
      \noindent #4
    \end{ocg}
  \par\medskip
}

%% ------------------- Listas----------------------%%
    % --- Lista alfabética ---
    \newenvironment{la}
      {\begin{list}{\alph{alphaitem})}{%
          \usecounter{alphaitem}%
          \setlength{\leftmargin}{1cm}%
          \setlength{\parskip}{3pt}% 
          \setlength{\itemsep}{0pt}% ← espacio entre ítems
          \setlength{\parsep}{0pt}%  ← párrafos dentro de un ítem
      }}
      {\end{list}}
      
    % --- Lista numérica ---
    \newenvironment{lnum}
      {\begin{list}{\arabic{numitem}.}{%
          \usecounter{numitem}%
          \setlength{\leftmargin}{1cm}%
          \setlength{\parskip}{3pt}% 
          \setlength{\itemsep}{0pt}% ← espacio entre ítems
          \setlength{\parsep}{0pt}%  ← párrafos dentro de un ítem
      }}
      {\end{list}}
      
% ------------------- Imágenes----------------------%
    \graphicspath{{Img/}}% Rutas por defecto 
    % --- Imagen adaptativa ---
    % Registros de dimensión definidos una sola vez
    \newdimen\imgcap
    \newdimen\imgmin
    \newdimen\imgmax
    \newdimen\imgremain
    \newdimen\imgh
    \newdimen\imgneed
    
    \newcommand{\imagenfit}[3]{%
      \addvspace{10pt}% separación antes
      \par\begingroup
        % Parámetros de composición (asignaciones, no \newdimen)
        \imgcap=1\baselineskip            % alto estimado de título+fuente
        \imgmin=0.15\textheight
        \imgmax=0.15\textheight
        \imgremain=\pagegoal
        \ifdim\pagegoal=\maxdimen \imgremain=0pt\fi
        \advance\imgremain by -\pagetotal
        \advance\imgremain by -\imgcap
        % Decisión de altura destino
        \ifdim\imgremain<\imgmin
          \newpage
          \imgh=0.20\textheight
        \else
          \imgh=\imgremain
          \ifdim\imgh>\imgmax \imgh=\imgmax \fi
        \fi
        % Reservar espacio para evitar cortes del bloque completo
        \imgneed=\imgh \advance\imgneed by \imgcap
        \Needspace{\imgneed}
        % Cuerpo
        \noindent\begin{minipage}{\textwidth}
          \centering
          \captionof{figure}{\textemdash\ #2}\par\vspace{0.3em}% título arriba
          \includegraphics[width=\linewidth,height=\imgh,keepaspectratio]{\detokenize{#1}}\par
          \vspace{0.3em}\small\textit{Fuente: }#3% fuente abajo
        \end{minipage}%
      \endgroup
      \par\addvspace{10pt}%
    }

%------------ Bloque de RECUADRO ESPECIAL -------------%
    \newenvironment{specialblock}[1]{%
      \begin{quote} % crea sangría a izquierda y derecha
      \small % texto más pequeño
      \noindent\textbf{#1}\\[-0.3em] % título en negrita
      \rule{\linewidth}{0.4pt}\\[0.6em]}{
      \end{quote}}
      
%-------------------------- Portada -----------------------%
\newcommand{\oldtitle}[1]{\def\OldTitle{#1}}
\newcommand{\changerationale}[1]{\def\ChangeWhy{#1}}

\newcommand{\changebox}{%
\begin{tcolorbox}[title={Historial de cambios del tema}]
\textbf{Título anterior:} \OldTitle\par
\textbf{Motivo del cambio:} \ChangeWhy
\end{tcolorbox}
}

%-------------------------- Author Font -----------------------%
    \newcommand{\authorfont}[1]{{\fontfamily{EBGaramond-TLF}\selectfont #1}}

%-------------------------- Anexos -----------------------%
    \makeatletter
    \newcounter{anexo}
    \renewcommand{\theanexo}{\Roman{anexo}}
    \newcommand{\anexo}[1]{%
      \clearpage
      \phantomsection
      \refstepcounter{anexo}%
      \noindent\textbf{Anexo \theanexo.- }#1\par\medskip
      \addcontentsline{stc}{section}{Anexo \theanexo.- #1}%
    }
    \makeatother

    %-------------------------- Lista de figuras (personal) -----------------------%
\makeatletter
\newcommand{\MyListOfFigures}{%
  \clearpage
  \phantomsection
  {\centering\bfseries\Large Índice de figuras\par}%
  \medskip
  % Entrada en nuestro índice de contenido (stc)
  \addcontentsline{stc}{section}{Índice de figuras}%
  % Recuperación del listado (sin cabecera propia)
  \begingroup
    \parindent=0pt\parskip=2pt
    \@starttoc{lof}%
  \endgroup
}
\makeatother

%-------------------------- Bibliografía -----------------------%
\makeatletter
% Contador para las referencias
\newcounter{myref}
% Cabecera de la bibliografía, entra en el índice 'stc'
\newcommand{\MyBibliography}{%
  \clearpage
  \phantomsection
  {\centering\bfseries\Large Bibliografía y referencias\par}%
  \medskip
  \addcontentsline{stc}{section}{Bibliografía y referencias}%
  \setcounter{myref}{0}%
}
\newcommand{\MyBibItem}[4]{%
  \refstepcounter{myref}%
  \noindent[\themyref]\ #1.\ \emph{#2}. #3. #4\par
}
\makeatother

%--------- Establecer cuatro niveles de índice --------%
    \setcounter{tocdepth}{4}   
    \setcounter{secnumdepth}{4}% numera hasta \paragraph

%%%%%%%%%%%%%%%%%%%%%%%%%%%%%%%%

\makeatletter

    %--------- Interlineado Global --------%
    \edef\GlobalBaseLineStretch{\baselinestretch}
    
    %--------- Espaciado de seccinones --------%
    \renewcommand\section{\@startsection{section}
      {1}{\z@}
      {2.5ex \@plus 1ex \@minus .2ex} % espacio antes
      {1.5ex \@plus .2ex}             % espacio después
      {\normalfont\Large\bfseries}    % formato del título
    }
    \renewcommand\subsection{\@startsection{subsection}{2}{\z@}
      {3ex \@plus 0.8ex \@minus .2ex}
      {1.5ex \@plus .2ex}
      {\normalfont\large\bfseries}
    }
    \renewcommand\subsubsection{\@startsection{subsubsection}{3}{\z@}
      {3ex \@plus 0.5ex \@minus .2ex}
      {1ex \@plus .2ex}
      {\normalfont\normalsize\bfseries}
    }
    \renewcommand\paragraph{\@startsection{paragraph}{4}{\z@}%
    {-2ex \@plus -0.5ex \@minus -.2ex}% espacio antes
    {0.8ex \@plus .2ex}% espacio después
    {\normalfont\normalsize\itshape}}% ← cursiva en lugar de negrita

    %--------- Ecuaciones --------%   
    % Variables globales para almacenar títulos y notas
    \gdef\leq@current@section{}%
    \gdef\leq@current@subsection{}%
    \gdef\leq@last@printed@section{}%
    \gdef\leq@last@printed@subsection{}%
    
    % Comando para añadir nota (invisible en el cuerpo)
    \newcommand{\eqnote}[1]{%
      \addcontentsline{leq}{leqnote}{#1}%
    }
    
    % Comando para ecuaciones
    \newcommand{\eqblock}[2]{%
      % Verificar si necesitamos imprimir el título de sección
      \ifx\leq@current@section\leq@last@printed@section\else
        \addcontentsline{leq}{leqsection}{\leq@current@section}%
        \global\let\leq@last@printed@section\leq@current@section
        \gdef\leq@last@printed@subsection{}% Reset subsección al cambiar sección
      \fi
      % Verificar si necesitamos imprimir el título de subsección
      \ifx\leq@current@subsection\leq@last@printed@subsection\else
        \ifx\leq@current@subsection\@empty\else
          \addcontentsline{leq}{leqsubsection}{\leq@current@subsection}%
          \global\let\leq@last@printed@subsection\leq@current@subsection
        \fi
      \fi
      % Ecuación
      \refstepcounter{eqnum}%
      \begin{equation*}
        #1\tag{E.\theeqnum}
      \end{equation*}%
      \addcontentsline{leq}{leqt}{[E.\theeqnum] -- #2}%
      \addcontentsline{leq}{leqm}{#1}%
    }
    
    %%% ----- Índice de ecuaciones ----------%%%
    % Tamaño para []
    \newcommand{\leqIndexFont}{\normalsize}
    \newcommand{\leqTitleFont}{\tiny}
    \newcommand{\leqNoteFont}{\tiny}
    % Cabecera de sección 
    \newcommand*\l@leqsection[2]{%
      \addvspace{25pt}% Mayor separación antes de nueva sección
      \noindent\leqIndexFont
      \makebox[\linewidth]{\rule{.38\linewidth}{0.4pt}\ \ \textbf{  \MakeUppercase{#1}}\ \ \rule{.38\linewidth}{0.4pt}}%
      \par\vspace{-10pt}
    }
    % Cabecera de subsección 
    \newcommand*\l@leqsubsection[2]{%
      \par\addvspace{15pt}%
      \noindent\leqIndexFont #1
      \par\vspace{-7pt}
    }
    % Título de ecuación 
    \newcommand*\l@leqt[2]{%
    \par\addvspace{10pt}%
      \par\noindent\leqTitleFont #1\par
    }
    % Miniatura de ecuación
    \newcommand*\l@leqm[2]{%
      \vspace{0.3\baselineskip}%
      \par\scriptsize\noindent\hspace*{1.5em}\(#1\)\par%
    }
    % Nota de ecuación 
    \newcommand*\l@leqnote[2]{%
      \vspace{0.2\baselineskip}%
      \par\noindent\leqNoteFont\hspace*{1.5em}\textbf{Nota.--} #1\par\vspace{0.5\baselineskip}%
    }
    % Generación del índice
    \newcommand{\mylistofequations}{%
      \clearpage
      \phantomsection
      % Título visual de la lista (ajústalo a tu gusto):
      {\centering\bfseries\Large Índice de ecuaciones\par}
      % Entrada en nuestro índice 'stc':
      \addcontentsline{stc}{section}{Índice de ecuaciones}%
      % Llamada a tu lista real (usa el nombre exacto de tu macro):
      \@starttoc{leq}%
    }
    
    %--------- Captura de títulos de sección --------%
    \let\leq@original@section\section
    \let\leq@original@subsection\subsection
    % Flag para saber si estamos en una sección asterisco
    \newif\ifLeqStarSection
    \LeqStarSectionfalse
    
    % Redefinir \section CON CONTROL DE ASTERISCO
    \renewcommand{\section}{%
      \@ifstar{\leq@section@star}{\@dblarg\leq@section@nostar}%
    }
    \def\leq@section@star#1{%
      \global\LeqStarSectiontrue
      \gdef\leq@current@section{#1}%
      \gdef\leq@current@subsection{}%
      \leq@original@section*{#1}%
      \global\LeqStarSectionfalse
    }
    \def\leq@section@nostar[#1]#2{%
      \global\LeqStarSectionfalse
      \gdef\leq@current@section{#2}%
      \gdef\leq@current@subsection{}%
      \leq@original@section[#1]{#2}%
    }
    % Redefinir \subsection
    \renewcommand{\subsection}{%
      \@ifstar{\leq@subsection@star}{\@dblarg\leq@subsection@nostar}%
    }
    \def\leq@subsection@star#1{%
      \global\LeqStarSectiontrue
      \gdef\leq@current@subsection{#1}%
      \leq@original@subsection*{#1}%
      \global\LeqStarSectionfalse
    }
    \def\leq@subsection@nostar[#1]#2{%
      \global\LeqStarSectionfalse
      \gdef\leq@current@subsection{#2}%
      \leq@original@subsection[#1]{#2}%
    }

    % ----- Restauración de valores Globales de estilo ----------%
    % Flags
    \newif\ifInFootnote
    \InFootnotefalse
    \pretocmd{\@footnotetext}{\InFootnotetrue}{}{}
    \apptocmd{\@footnotetext}{\InFootnotefalse}{}{}
    % Profundidad de listas (soporta anidadas)
    \newcount\MyListDepth \MyListDepth=0
    \AtBeginEnvironment{la}{\advance\MyListDepth by 1}
    \AfterEndEnvironment{la}{\advance\MyListDepth by -1}
    \AtBeginEnvironment{lnum}{\advance\MyListDepth by 1}
    \AfterEndEnvironment{lnum}{\advance\MyListDepth by -1}
    \AtBeginEnvironment{itemize}{\advance\MyListDepth by 1}
    \AfterEndEnvironment{itemize}{\advance\MyListDepth by -1}
    \AtBeginEnvironment{enumerate}{\advance\MyListDepth by 1}
    \AfterEndEnvironment{enumerate}{\advance\MyListDepth by -1}
    % Restauración diferida: hazlo al INICIO del próximo párrafo real
    \newif\if@pendingRestore
    \@pendingRestorefalse
    \newcommand{\RequestRestore}{\global\@pendingRestoretrue}
    \everypar{%
      \if@pendingRestore
        \global\@pendingRestorefalse
        \ifInFootnote\else
          \ifnum\MyListDepth=0
            \setlength{\parskip}{\GlobalParskip}%
            \setstretch{\GlobalBaseLineStretch}%
          \fi
        \fi
      \fi
    }
    % Al final de bloques:
    \AfterEndEnvironment{equation}{\RequestRestore}
    \AfterEndEnvironment{equation*}{\RequestRestore}
    \AfterEndEnvironment{la}{\RequestRestore}
    \AfterEndEnvironment{lnum}{\RequestRestore}
    % Al final de macros:
    \apptocmd{\eqblock}{\RequestRestore}{}{}
    \apptocmd{\imagenfit}{\RequestRestore}{}{}
    
\makeatother

% ===================== ToC propio con 4 niveles (único bloque) =====================
\makeatletter

% --- Escritores: añadimos una línea al .stc en cada cabecera numerada ---
% Guardamos las versiones actuales para llamar al comportamiento normal
\let\MyTOC@orig@section\section
\let\MyTOC@orig@subsection\subsection
\let\MyTOC@orig@subsubsection\subsubsection
\let\MyTOC@orig@paragraph\paragraph

% SECTION
\renewcommand{\section}{%
  \@ifstar{\MyTOC@section@star}{\@dblarg\MyTOC@section@nostar}%
}
\def\MyTOC@section@star#1{%
  \MyTOC@orig@section*{#1}% sin entrada al .stc
}
\def\MyTOC@section@nostar[#1]#2{%
  \MyTOC@orig@section[{#1}]{#2}% comportamiento normal (escribe al .toc)
  % Escribimos también nuestra línea al .stc (texto con \numberline)
  \addcontentsline{stc}{section}{\protect\numberline{\thesection}#2}%
}

% SUBSECTION
\renewcommand{\subsection}{%
  \@ifstar{\MyTOC@subsection@star}{\@dblarg\MyTOC@subsection@nostar}%
}
\def\MyTOC@subsection@star#1{%
  \MyTOC@orig@subsection*{#1}%
}
\def\MyTOC@subsection@nostar[#1]#2{%
  \MyTOC@orig@subsection[{#1}]{#2}%
  \addcontentsline{stc}{subsection}{\protect\numberline{\thesubsection}#2}%
}

% SUBSUBSECTION
\renewcommand{\subsubsection}{%
  \@ifstar{\MyTOC@subsubsection@star}{\@dblarg\MyTOC@subsubsection@nostar}%
}
\def\MyTOC@subsubsection@star#1{%
  \MyTOC@orig@subsubsection*{#1}%
}
\def\MyTOC@subsubsection@nostar[#1]#2{%
  \MyTOC@orig@subsubsection[{#1}]{#2}%
  \addcontentsline{stc}{subsubsection}{\protect\numberline{\thesubsubsection}#2}%
}

% PARAGRAPH
\renewcommand{\paragraph}{%
  \@ifstar{\MyTOC@paragraph@star}{\@dblarg\MyTOC@paragraph@nostar}%
}
\def\MyTOC@paragraph@star#1{%
  \MyTOC@orig@paragraph*{#1}%
}
\def\MyTOC@paragraph@nostar[#1]#2{%
  \MyTOC@orig@paragraph[{#1}]{#2}%
  % Si no numeras los \paragraph, cambia \theparagraph por texto plano sin \numberline
  \addcontentsline{stc}{paragraph}{\protect\numberline{\theparagraph}#2}%
}

% --- Impresor del índice propio (lee el .stc) ---
\newcommand{\MyTOC}{%
  \par\bigskip
  {\centering\bfseries\Large Índice de contenido\par}%
  \medskip
  \begingroup
    \parindent=0pt\parskip=2pt
    \@starttoc{stc}%
  \endgroup
}

\makeatother
% ===================== fin ToC propio =====================


%%%%%%%%%%%%%%


% --- Datos del tema ---
\title{Sistema económico internacional (I)\\
\large Desde el siglo XIX hasta la ruptura del sistema de Bretton-Woods}
\author{Marcelo Ortega}
\date{Marzo 2026}
\newcommand{\TemaCode}{3B20}

\oldtitle{
    B.21. El sistema económico desde el siglo XIX hasta la ruptura del sistema de Bretton-Woods
}
\changerationale{
    Sin cambios.
}

\begin{document}


\maketitle
\changebox

\newpage

% --- Página de resumen y modificaciones ---
\puntosclave{ %% Escribir puntos clave Apuntes CECO
     Se ha decidido poner mayor énfasis en la evolución del sistema monetario internacional (SMI). Esto se debe a la influencia que ejerce sobre las demás relaciones de tipo económico que se establecen entre diferentes países. Sin embargo, ha de equilibrarse el contenido de la exposición de tal manera que se traten igualmente los principales desarrollos que tuvieron lugar en relación con: (i) el comercio internacional; (ii) los movimientosde factores productivos; (iii) la situación económica general; y, (iv) la Política Económica.

     En cuanto al inicio del período histórico que se estudia, tradicionalmente se ha considerado que debe situarse en torno a mediados del siglo XIX, época en la que comienza la conocida como Primera Globalización y en la que la industrialización de las principales economías desarrolladas adquiere una mayor entidad (segunda revolución industrial).

     Como tema de Historia económica, una exposición correcta consistirá en: (i) una cobertura completa y rigurosa de los desarrollos históricos de contenido económico; ampliado por, (ii) una lectura analítica puramente económica (incentivos, fallos del sistema, determinantes del crecimiento y de los flujos comerciales y de factores, etc) de los fenómenos y acontecimientos que se recogen.
    }
\vspace{1cm}

\modificaciones{
    
    }

\newpage
\MyTOC
\newpage

\mylistofequations
\clearpage

\MyListOfFigures
\clearpage


\pagenumbering{arabic}
\setcounter{page}{1}


\section{Introducción}



\subsection{Contextualización}


\subsubsection{Relevancia}




\subsubsection{Problemática}



\subsection{Estructura de la exposición}



\clearpage
\section{Primera Globalización: s.XIX-1913}
\puntosclave{
    La época comprendida entre mediados del siglo XIX y el estallido de la I Guerra Mundial estuvo marcada por elevados crecimientos de la población, la economía mundial y el comercio internacional. 
    
    En el plano comercial, factores como la revolución de los transportes y comunicaciones y unas políticas comerciales tendentes, en general, al libre comercio, fueron decisivos para la expansión de la red de comercio internacional.
    
    Durante este período también se produjeron flujos de inversión y migratorios internacionales sin precedentes, lo cual suponía una condición necesaria para el crecimiento económico de muchas regiones.
    
    En el ámbito monetario, la adopción del patrón oro universal, aunque no exenta de fallos, proporcionó una cierta estabilidad interna y externa y facilitó el desarrollo del comercio internacional.
}


Durante los siglos XVII y XVIII, el comercio mundial se había expandido a un ritmo constante de alrededor del 1\% anual, superando el aumento de los ingresos mundiales, aunque no por mucho. En algún momento del comienzo del siglo XIX, el \textbf{comercio} mundial empezó a crecer a pasos agigantados, registrando una tasa sin precedentes de casi el 4\% anual durante todo el siglo. 

Los costes de transacción que impiden el comercio a larga distancia empezaron a caer vertiginosamente. Los flujos de capital experimentaron un gran auge y las economías mundiales fueron objeto de una \textbf{integración financiera} sin precedentes. Fue también esta una época de grandes \textbf{movimientos de población} entre continentes, en la que europeos de clase obrera emigraron en masa a América y a otros territorios de reciente colonización. 

Por estas razones, la mayoría de los historiadores económicos consideran que el largo siglo que termina en 1914 fue la primera fase de la globalización. De hecho, desde muchos puntos de vista, la economía mundial no ha superado hasta hace poco los niveles de globalización de 1913 en cuanto a comercio y finanzas. En términos de volumen de emigración, aún no lo ha alcanzado.

¿Qué hizo posible esta fase de globalización? Las interpretaciones más habituales identifican tres cambios importantes en este periodo:
\begin{lnum}
    \item \textbf{Segunda revolución industrial}. Primero, nuevas tecnologías en forma de barcos de vapor, ferrocarriles, canales y telégrafos revolucionaron el transporte y las comunicaciones internacionales y redujeron en gran medida los costes comerciales desde los primeros años del siglo XIX;
    \item \textbf{Libre comercio: consenso teórico}. Segundo, la narrativa económica cambió a medida que las ideas de los economistas del libre mercado, como SMITH y RICARDO, fueron tomando fuerza. Esto llevó a los gobiernos de las principales economías del mundo a relajar sustancialmente las restricciones sobre el comercio en forma de aranceles y contingentes a la importación;
    \item \textbf{Sistema monetario internacional}. Finalmente, a partir de la década de 1870, la adopción generalizada del patrón oro permitió que el capital se moviera internacionalmente sin temor a cambios arbitrarios en los valores de las monedas u otros altibajos financieros. 
\end{lnum}

Sin embargo, estos argumentos no lo explican todo. 

\textbf{¿Qué queremos decir?} Como sabemos, los costes de transacción en la economía mundial van más allá del transporte, los aranceles y la inestabilidad monetaria. Por ello, a lo largo de esta sección presentaremos la estructura que seguirán las sucesivas con el fin de revestir estos tres argumentos básicos de un contexto geopolítico que permita entenderlos e interpretarlos de manera apropiada. Como veremos, describir las relaciones económicas internacionales en un periodo concreto de la historia requiere: (i) contextualizarlas; (ii) identificar los cambios; y, (iii) saberlos proyectar sobre estas tres variables de interés económico: tecnología, teoría y manifestación práctica.

\subsection{¿Cuál es el contexto geopolítico?: Imperialismo}
Empecemos, por ello, entendiendo cómo era el mundo, desde una perspectiva política, durante la primera globalización.

Tanto si era de tipo formal como de tipo informal, el imperialismo fue un mecanismo que sirvió para imponer reglas favorables al comercio, de tal manera que los gobiernos de los países avanzados se convirtieron en los árbitros que hacían cumplir dichas reglas. Las políticas imperialistas desplegaban el poder político y militar de los países poderosos para hacer que el resto del mundo se plegara a sus intereses, siempre que fuera posible. De ahí que proporcionaran un apoyo básico para impedir que la globalización descarrilara en las zonas periféricas de la economía mundial (Latinoamérica, Asia y Oriente Próximo)⁠ de manera que estas regiones fueran seguras para el comercio y las finanzas internacionales. En este capítulo se explica cómo surgió la versión siglo XIX de la globalización y de qué modo las políticas internas de los países lograron desbaratarla. Empezamos nuestro recorrido con las políticas comerciales y seguiremos con el patrón oro.



Una vez comprendemos el contexto geopolítico, podemos pasar a analizar los aspectos económicos que caractrizaron la primera globalización.

\subsection{Comercio internacional: la victoria (limitada) del libre comercio}
Las creencias favorables a la libertad de comercio fueron expandiéndose durante todo el siglo XIX gracias a los esfuerzos de economistas como RICARDO y MILL, que ampliaron las observaciones de SMITH para mostrar que el comercio sin restricciones es bueno para todos los países que participan de él. Esto condujo progresivamente a la convergencia en los sistemas de creencias de quienes tomaban las decisiones económicas clave en ese periodo, permitiendo abandonar las ideas mercantilistas y adoptar el liberalismo económico y las reglas del patrón oro, que unieron a los gobernantes de distintas naciones en su propósito de minimizar los costes de transacción. 

No obstante, estas ideas eran elegantes, potentes y podían describirse con precisión lógica pero, como veremos, su influencia fue diferente de país a país y en momentos distintos del tiempo. 

\subsubsection{Economías avanzadas: ¿libre comercio?}
La fecha crucial en la historia de los aranceles del siglo XIX es 1846, el año en que Gran Bretaña suprimió los aranceles sobre las importaciones de cereales que había instaurado durante las Guerras Napoleónicas.\footnote{%
    Estas llamadas Corn Laws fueron el centro de sonoros combates políticos en la Gran Bretaña de principios del siglo XIX, puesto que enfrentaron los intereses rurales con los intereses urbanos. (Aquí, corn se refiere a cereales, y los aranceles en cuestión se aplicaban a todas las importaciones de alimentos y cereales.) Los terratenientes querían aranceles altos que mantuvieran elevados los precios de los alimentos e hicieran subir sus ingresos. Los fabricantes, cada vez más poderosos debido a la Revolucion Industrial que se extendía por Londres, Manchester y otras ciudades, querían suprimir los aranceles para reducir el coste de la vida. Esa reducción, como argumentaba MARX, entre otros, permitiría a los capitalistas pagar salarios aún más bajos a sus trabajadores. La pugna por las Corn Laws ilustra un tema al que tendremos muchas ocasiones de volver: debido a que las políticas comerciales tienen consecuencias importantes en la distribución de los ingresos, se ven envueltas en luchas políticas de gran envergadura. Los economistas pueden criticar la artificialidad de los costes de en transacción que crean las barreras comerciales impuestas por los gobiernos, pero este argumento no siempre prevalece si existen fuertes intereses políticos o argumentos económicos de peso que vayan en dirección contraria.

Este debate movilizó a la sociedad y la política británicas, y surgieron bandos a favor y en contra de las Corn Laws, comprometidos en lo que parecía ser una lucha feroz por unos cuantos impuestos sobre las importaciones pero que, en realidad, era un enfrentamiento para decidir quién iba a gobernar en Gran Bretaña y a prosperar en los años venideros. La conocida revista The Economist es producto de esa época, fundada por opositores a las Corn Laws para difundir y popularizar opiniones en favor de la libertad de comercio, un papel que continúa desempeñando en la actualidad. Al final, ganaron los intereses de los empresarios emergentes: tenían de su parte tanto los argumentos intelectuales como a las fuerzas de la Revolución Industrial.
} Muchos percibieron la reforma como un éxito político y económico; los comentaristas económicos del continente europeo señalaban con admiración el gran aumento del comercio y de los ingresos en Gran Bretaña desde la derogación (aunque, naturalmente, era la Revolución Industrial la que en realidad se merecía la mayor parte del mérito). 

En Gran Bretaña, la ideología del libre comercio no solo dominaba el discurso público en un grado mucho mayor, sino que proteccionismo llegó a convertirse en un término peyorativo con el que se intimidaba a los oponentes. Además, la fuerte posición comercial británica en el sector industrial hacía de la arancelaria una política más bien redundante y carente de sentido. Cuando el primer ministro William Gladstone se burló de quienes buscaban venganza en la Política Comercial, se refería a ambos factores. ¿Qué diablos significa \textit{comercio justo}?, preguntó. \textit{Pues bien, señores}, respondió empleando un argumento que a partir de entonces repetirían incansablemente quienes estaban a favor del libre comercio, \textit{he de decir que tiene un sospechoso parecido a nuestra vieja amiga Protección}. Era ni más ni menos que el proteccionismo intentaba resultar más atractivo adoptando un nombre nuevo. Gladstone señalaba que el Reino Unido iba a ganar poco de adoptar este tipo de represalias, puesto que las importaciones de bienes manufacturados eran mucho menores que las exportaciones. Los aranceles que se gravaran en unas bases tan diminutas apenas iban a tener efecto sancionador. El enorme excedente comercial de Gran Bretaña, declaró Gladstone, hacía que el libre comercio fuera la mejor política para el país: el dominio británico en la industria ayudó a evitar una oleada proteccionista. 

Sin embargo, el aparente éxito británico no facilitó necesariamente la liberalización del comercio otros países.\footnote{%
    Como el emperador Napoleón III comentó a R. COBDEN, miembro del Parlamento británico y proselitista del libre comercio: \textit{Me atrae y me gusta la idea de realizar una obra similar en mi país; pero es muy difícil hacer reformas en Francia; en Francia hacemos revoluciones, no reformas}.
}

\paragraph{Transaccionalidad de la Política Comercial: Tratado COBDEN-CHEVALIER (1860)}
Hizo falta otro argumento que sellaría el destino de las negociaciones comerciales.

Poco después de la abolición, y tras abundantes alagos hacia la hazaña británica, comenzó a surgir un recurso político al que los defensores del libre comercio han acudido desde entonces: reducir las barreras comerciales a cambio de que otro país haga lo mismo, y luego presentar la liberalización a la oposición como una \textbf{concesión} necesaria para conseguir que la otra parte abra sus mercados. Emerge así un término poco utilizado pero frecuentemente utilizado en la negociación comercial: la transaccionalidad de la Política Comercial. (El argumento teórico está bien, pero queremos algo a cambio).

Uno de los más célebres ejemplos de este exitoso nuevo argumento fue el Tratado COBDEN-CHEVALIER de 1860, que comprometía a Gran Bretaña a reducir sus impuestos sobre las bebidas alcohólicas francesas a cambio de que Francia redujera sus aranceles sobre los bienes manufacturados británicos. Además, trajo consigo una innovación importante importante en materia comercial: \textbf{la cláusula de nación más favorecida} (NMF). 

\paragraph{Europa continental: Un equilibrio inestable y frágil}
La red de tratados comerciales a la que este recurso retórico dió comienco, se convirtió en un importante instrumento de reducción arancelaria en toda Europa en las décadas de 1860 y 1870. 

A mediados de los años 1870, la mayoría de las prohibiciones al comercio habían desaparecido y la media de los aranceles sobre productos manufacturados permaneció en números bajos de una sola cifra en Gran Bretaña, Alemania, los Países Bajos, Suecia y Suiza, y un poco por encima de la decena en Francia e Italia, muy por debajo de los niveles anteriores, que eran múltiplos de estos índices.\footnote{%
    Las razones individuales por las cuales los países adoptaron medidas orientadas al libre comercio fueron muy variadas. En Países Bajos, la presión ejercida por agentes interesados en la exportación, como banqueros, mercaderes y dueños de navíos jugó un papel crucial. En Francia, por el contrario, el interés del sector industrial en lograr materias primas baratas y bienes de equipo fue determinante. En Alemania, la tendencia hacia políticas de libre comercio fue defendida por la aristocracia prusiana que contralaba la exportación de productos agrícolas.
}

Sin embargo, las dudas no tardaron en llegar. Un detonante importante, como suele ocurrir, fue una prolongada recesión económica que había empezado en la década de 1870, golpeando con especial dureza a los agricultores y ganaderos. La revolución en el transporte y los recortes arancelarios dieron como resultado un influjo de granos procedentes del Nuevo Mundo y un abrupto descenso de los precios. Por todo el continente, los agricultores pedían protección, a menudo haciendo causa común con industriales que se tambaleaban por la competencia de los más avanzados productores británicos (y cada vez más por la de los exportadores estadounidenses también).\footnote{%
    En la Alemania de BISMARCK, esto llevó al famoso \textbf{matrimonio del hierro y el centeno}, una coalición entre agricultores e industriales que dio lugar a un brusco aumento de los aranceles desde finales de los años 1870. BISMARCK, un político siempre inteligente y astuto, racionalizó la nueva política quejándose de que Alemania se había convertido en el vertedero del exceso de producción de otros países. Francia y otras potencias de la Europa continental le siguieron, aumentando sus propios aranceles, y la tendencia general hacia el endurecimiento de las restricciones al comercio continuó hasta el comienzo de la Primera Guerra Mundial.
} Hacia 1913, la media de los aranceles a la importación en la Europa continental sobre productos manufacturados se había duplicado y era de casi el 20\%. Así, muchas de las potencias de la Europa continental, volvieron a niveles de protección más altos después de un par de décadas. 

En Europa, fue solo en Gran Bretaña donde hubo resistencia a la resbaladiza pendiente proteccionista en las décadas anteriores, el resto de políticas nacionales demostraron ser propicias al libre comercio únicamente cuando la supremacía económica garantizaba inmunidad relativa frente a la competencia de las importaciones. 

\paragraph{Estados Unidos: El argumento de la industria naciente}
De hecho, el libre comercio no triunfó, ni por un momento, en todas partes. A pesar de las diferencias en la composición de la vida política entre Estados Unidos y Gran Bretaña durante la primera parte del siglo XIX, los dos países tenían una cosa en común: las controversias sobre los aranceles estuvieron en el centro de la política nacional.\footnote{%
    Como reflexionaría posteriormente un exasperado legislador de Pensilvania, \textit{el hombre es un animal que hace discursos sobre aranceles}.
} En este caso, la Política Comercial se nutría directamente de la división política y social más importante del país, la existente entre el Norte y el Sur.\footnote{%
    El Sur, esclavista, estaba organizado en torno a una economía de exportación basada en el tabaco y el algodón. Los Estados libres del Norte dependían de una incipiente base manufacturera que iba por detrás de Gran Bretaña en productividad y luchaba por competir con productos importados que eran más baratos.
}

En este sentido, la Guerra Civil de 1861-1866 fue motivada tanto por la política comercial estadounidense como por la esclavitud. La prosperidad del Sur dependía del comercio internacional, mientras que el Norte quería protección frente a las importaciones, al menos hasta que pudiera alcanzar cierto nivel. Al comienzo de la guerra LINCOLN subió los aranceles y la protección al comercio aumentó todavía más después de la victoria del Norte. La media de los aranceles sobre productos manufacturados fue del 45\% en los diez años siguientes a 1866 y nunca bajaron mucho de ese nivel hasta la Primera Guerra Mundial. Desde cualquier punto de vista, Estados Unidos era un país muy proteccionista a finales del siglo XIX.\footnote{%
    Además, resulta interesante que esta situación representa un caso en el que claramente el libre comercio no estuvo al servicio de una causa política progresista. Como ha escrito el destacado politólogo KEOHANE: \textit{Seguir la lógica del mercado tuvo efectos trágicos a largo plazo. Elimpacto económico en el Sur del crecimiento sin diversificación ni industrialización ya fue bastante perjudicial. Mucho más graves fueron los resultados políticos y sociales de hacer que el algodón fuera el rey: la esclavitud se enquistó y la guerra civil se hizo cada vez más probable}. 

Al margen de las demás consecuencias económicas, el libre comercio en Estados Unidos en el siglo XIX habría reforzado y fortalecido más la esclavitud como institución política y social. Solo se pueden hacer conjeturas sobre el daño que esto habría hecho al desarrollo de las instituciones políticas del país, pero no parece que la imagen hubiera sido bonita. La lección está clara: según dónde esté un país en relación con la economía mundial y cómo se alineen las políticas comerciales con sus divisiones políticas y sociales, el libre comercio puede ser una fuerza progresista o retrógrada.
}

En definitiva, Gran Bretaña fue el centro neurálgico industrial del mundo a mediados del siglo XIX y además las políticas de liberalizaron favorecían los intereses de los empresarios y las clases medias. En el extremo opuesto, Estados Unidos era un rezagado industrial con una ventaja en costes en las actividades propias de las plantaciones con esclavos, donde las políticas comerciales liberales habrían favorecido los intereses agrarios represivos. Libre comercio y buena política: ni van siempre juntos, ni son muchas veces compatibles.

\subsubsection{Proyección internacional: Imperialismo}
¿Y qué pasaba en el resto del mundo? Muchas veces se aborda este tema argumentando la gran brecha tecnológica que había entre las exportaciones de las regiones más desarrolladas y las menos. Sin embargo, esta aproximación no tiene mucho sentido una vez hemos comprendido el contexto geopolítico de la época: \textbf{si entre los países avanzados el libre comercio dependía de un equilibrio inestable y frágil entre una ideología compartida y una constelación de intereses políticos interiores, en el resto del mundo se imponía sobre todo desde el exterior}. Por tanto, ni era libre ni era comercio, atendiendo a la interpretación que le da la Teoría Económica.

Un reciente estudio estadístico constata que dos países miembros de un mismo imperio, \textit{ceteris paribus}, tenían el doble de volumen de comercio entre ellos que con otros que no pertenecían al imperio. ¿La razón?: \textit{Los imperios aumentaron el comercio bajando los costes de transacción y estableciendo políticas que estimulaban el comercio dentro del imperio}.\footnote{%
    Los ejemplos concretos de costes de transacción reducidos que los autores del estudio logran cuantificar son: el uso de una lengua común, la presencia de una moneda común, el uso de la moneda en las colonias adquiridas recientemente y los acuerdos comerciales preferenciales.
} 

Sostiene FERGURSON, el historiador de Harvard, que: \textit{Ninguna organización de la historia ha hecho más para fomentar la libre circulación de bienes, capital y trabajo que el Imperio Británico en el siglo XIX y principios del siglo XX}. No es necesario aceptar la optimista interpretación de FERGUSON sobre el Imperio Británico para estar de acuerdo con su afirmación de que el imperialismo fue una fuerza tremendamente potente a favor de la globalización económica.\footnote{
En Asia, el imperialismo europeo garantizaba que se protegían los derechos de los extranjeros, se hacían cumplir los contratos, las disputas se resolvían según las reglas de países europeos, se recibía con los brazos abiertos a exportadores e inversores, las deudas se pagaban, se invertía en infraestructuras, se pacificaba a los nativos, se desbarataban las ambiciones de los incipientes nacionalismos, etc., neutralizando la larga lista de costes de transacción que podrían obstaculizar el comercio internacional.

Recordemos la sustitución de la Compañía Inglesa de las Indias Orientales por el Raj Británico cuando aquella demostró ser incapaz de controlar la insurgencia local, o el traslado de los poderes policiales de la Compañia de la Bahía de Hudson al Dominio de Canadá. 
} 

\paragraph{Dominio: formal e informal}
Sin embargo, sería un error pensar que esto emergía únicamente donde la metropoli empleaba su dominio imperial directo: \textbf{el imperialismo se daba tanto en la variante formal como en la informal}. Gran Bretaña, Francia y Estados Unidos (que acabó por entrar en el juego), así como las demás potencias, no siempre tenían que extender su dominio directo para doblegar a otras regiones a su voluntad. Bastaba con la amenaza militar y la presión política.

En un artículo clásico titulado \textit{The Imperialism of Free Trade}, GALLAGHER y ROBINSON muestran que existía un continuo entre la influencia informal y el dominio formal, este utilizado solamente como último recurso cuando la situación estaba demasiado descontrolada para lograr los efectos deseados a través de los gobernantes locales. Uno de los principales instrumentos de influencia informal fue el tratado comercial. Pero si resultaba que la población local no reverenciaba lo suficiente las ideas de SMITH y RICARDO, las cañoneras listas para el ataque siempre podían proporcionar la persuasión necesaria.\footnote{%
    De ese modo, Gran Bretaña firmó un tratado con la Turquía otomana en 1838 que obligaba al país a limitar los aranceles sobre las importaciones a un máximo del 5\% y a abolir las prohibiciones a la importación y los monopolios. Los británicos también lucharon en la llamada Guerra del Opio contra China en 1839-1842 para abrir ese país a las importaciones de opio y otros productos exportados desde el Imperio Británico. 

El contraalmirante Matthew C. Perry firmó un tratado con Japón en representación de Estados Unidos en 1854 para abrir ese país al tráfico marítimo y al comercio exterior. Este y otros tratados similares obligaron a poner techos a los impuestos sobre las importaciones (unilateralmente, por supuesto), restringieron la capacidad de los países menos poderosos de dirigir sus políticas comerciales de forma independiente, concedieron privilegios legales a los comerciantes extranjeros y reforzaron el acceso de estos a sus puertos.
}

\subsubsection{Implicaciones de Política Económica: ¿comercio y crecimiento?}
En definitiva, debemos comprender que, a pesar de la inequívoca explosión comercial (probablemente más que en cualquier otro momento de la historia, a excepción de estas últimas décadas), descansaba sobre pilares institucionales frágiles y poco funcionales y no tanto en el libre comercio, como se suele decir. Sin duda, la revolución en el transporte, el aumento de la productividad y los ingresos y la reducción de barreras arancelarias por debido a la convergencia de percepciones jugó un rol esencial en la globalización. Pero Gran Bretaña fue la única gran economía que mantuvo políticas comerciales abiertas durante un periodo más o menos prolongado. Estados Unidos puso aranceles muy altos sobre importaciones manufacturadas durante la Guerra Civil y los mantuvo altos durante todo el siglo. Las principales potencias de la Europa continental se convirtieron sin vacilación al libre comercio solamente durante un breve periodo en las décadas de 1860 y 1870. Y en el resto del mundo, simplemente se hizo y deshizo a voluntad.

\textbf{NOTA.-} Existe, desde una óptica liberal, un aspecto paradójico en este incremento de la protección a finales del siglo XIX en Europa con grandes Implicaciones de Política Económica. Como señala BAIROCH, no solo fue el volumen del comercio lo que aumentó rápidamente después de 1890, también lo hicieron los ingresos, especialmente en aquellos países que pusieron barreras comerciales. Más aún, esta experiencia no se circunscribe a Europa, pues se observa en el periodo que siguió a la Guerra Civil en Estados Unidos, lo que refuerza las dudas respecto a la existencia de una relación lineal entre políticas de libre comercio y crecimiento económico. 

\subsection{Libre movilidad de factores productivos}
Ahora que conocemos la primera globalización desde una óptica comercial, ¿qué podemos decir sobre la internalización de los factores productivos?

\subsubsection{Sistema Monetario Internacional: Patrón oro}
Empecemos con el capital. El sistema monetario internacional fue un sistema de tipos de cambio fijos que establecía un valor en oro para cada moneda, lo que le ha atribuido el nombre de \textbf{patrón oro}.

Su objetivo declarado sería la estabilización y credibilidad del sistema monetario. Por ello, no debe entenderse como una fórmula ad hoc para regular los tipos de cambio sino una proyección internacional de lo que ya se había establecido tiempo atrás en determinadas economías y había resultado rentable. Por tanto, las causas de este sistema no son unidireccionales: ni el auge en el comercio trajo consigo el patrón oro, ni el patrón oro trajo consigo un auge del comercio. Tanto el patrón oro como el comercio aumentaron de forma paralela durante la primera mitad del siglo XIX y, en aras de garantizar la estabilidad de precios en cada economía (cada vez más influida por las importaciones exteriores) se internacionalizó lo que había demostrado ser una mecanismo seguro y creíble: la convertibilidad fija. (Además de ser el sistema que el Reino Unido venía utilizando desde 1717, impuesta por NEWTON).

Lo que sí fue novedoso fueron los rasgos ténicos que caracterizaron su proyección internacional.

\paragraph{Rasgos técnicos: la libre convertibilidad}
El patrón oro se apoyaba en unas pocas reglas sencillas. Si cada moneda nacional tenía su paridad con el oro, el banco central de cada país permanecía atento para convertir las monedas nacionales en oro según estas paridades. Por ejemplo, si la libra esterlina se definía como 4 granos de oro puro y el dólar estadounidense como 1 gramo, cada emisor de moneda debía respaldar sus pasivos con oro y garantizar la convertibilidad. ¿Cómo construía esto el sistema monetario internacional? Si esta es la paridad fijada, los tipos de cambio de las divisas también estaban fijados irrevocablemente: una libra esterlina equivalía a 4 dólares (4/1 gramos de oro). Esto, sumado a la libre circulación de oro a nivel internacional asegura la estabilidad del sistema cambiario internacional. Parece tremendamente sencillo, y lo era. 

Bajo las reglas del patrón oro, los gobiernos no tenían posibilidad de recurrir a la Política Monetaria para alterar la situación del crédito interior, porque la oferta de dinero interior estaba determinada únicamente por la circulación de oro y capitales a través de sus fronteras nacionales. En consecuencia, los bancos centrales tenían poco que hacer, aparte de emitir o retirar moneda nacional a medida que fluctuaba la cantidad de oro en sus arcas. El sistema contaba con reglas claras, universales y no discrecionales: el régimen financiero minimizaba los costes de transacción a través de las fronteras nacionales y financieros e inversores no tenían que enfrentarse a sorpresas ni controles en la frontera.

Una vez conocemos las reglas, ¿cómo podemos garantizar su estabilidad? Existen dos visiones complementarias. 

\paragraph{Estabilidad: mecanismos de ajuste automático}
La primera son los mecanismos de ajuste automático que emergen en presencia de desequilibrios. Bajo estas reglas, los cambios en la oferta interior de dinero están estrechamente ligados a los movimientos de las reservas de oro, por tanto, cualquier desajuste debe compensarse por algún de los siguientes canales/vías:\footnote{%
    Las autoridades se comprometen a no esterilizar, es decir, se permiten que las variaciones en las reservas de oro se transformen en variaciones proporcionales en la cantidad de dinero en circulación.
}
\begin{la}
    \item \textbf{Vía precios}. Un país con un déficit en su balanza de pagos exterior pierde oro en sus relaciones comerciales, y experimenta una reducción en su oferta de dinero. Si el dinero y el crédito son escasos, aumentarán las tipos de interés, se canalizará mayor ahorro y disminuirán los precios nacionales. El deterioramiento de su relación real de intercambio producirá dos efectos compensatorios: (i) reducción de importaciones/déficit comercial gastos; y, (ii) una mayor competitividad comercial y exportaciones (entrada de capital extranjero); restableciendo el equilibrio de los pagos exteriores en su posición de equilibrio original;
    \item \textbf{Vía renta}. Poco después, KEYNES aportó una explicación alternativa del ajuste. Un déficit exterior provocará una salida de oro y reducción de la oferta monetaria. En este contexto, aumentarán los tipos de interés, se producirá una deflación y aumentará la competitividad. Además, este diferencial atraerá también capitales del exterior, financiando el desequilibrio. No obstante, también se reducirá la inversión afectando al nivel de renta nacional, reduciendo las importaciones. Esta interpretación confirma y desarrollo la capacidad de este sistema para corregir desequilibrios, también a costa de una reducción de la renta.
\end{la}

En la práctica, los bancos centrales tuvieron cierto espacio de maniobra y a veces se alejaron de estas reglas del juego. Por ejemplo, un país con un déficit comercial podía demorar o evitar una subida de las tasas de interés si había flujos de capital privado compensatorios que vinieran del otro lado de las fronteras. Además, uno de las cláusulas de escape más recurrentes eran los conocidos \textbf{puntos del oro}. Cuando un país superavitario decidía liquidar sus divisas frente al banco emisor debía tener presente los costes de transporte sobre la paridad establecida. Es decir, en cada momento el tipo de cambio se fija por la paridad, pero en la práctica, se fija por los activos y contragarantías (reservas de divisas) en posesión de éste frente a su acreedor comercial: si la demanda de divisa aumenta, su moneda nacional se va depreciando internamente (en el Banco), hasta llegar al punto en que el tipo de cambio sea igual al tipo de paridad más los costes de transporte. En este punto, le compensa pagar en oro y no descapitalizarse en divisas. Por tanto, el deudor nacional comprará oro (en vez de divisas) y lo exportará al país extranjero, donde el acreedor cambiará el oro por su moneda y quedará saldada la deuda.

\imagenfit{PuntosOro.png}{Puntos del oro}{ICEX-CECO. Tema 3.B.20}

El nuevo punto de equilibrio estará en $C$, donde parte de la demanda de divisas (RC) se cubre con oferta autónoma (RB) y otra parte (BC) se cubre con una pérdida de oro, en un movimiento acomodaticio de las reservas. Por tanto, el tipo de cambio podrá fluctuar dentro del intervalo $P_M-P_X$, dentro del cual no habrá entradas ni salidas de oro.

\paragraph{Estabilidad: Economía Política}
No obstante, por bien que suenen estos principios de estabilización automática en la práctica, la realidad es que el sistema se mantenía por el la credibilidad del compromiso. La mayoría de ajustes se producían por la disponibilidad de flujos de capital compensatoria y estos dependen de forma crítica, de las expectativas sobre el compromiso de los bancos centrales con la paridad del oro.\footnote{%
    Los mercados suponían que los gobiernos acabarían defendiendo las paridades aviento y marea. Así lo hacían porque este era el sistema de creencias que gobernaba el comportamiento de los bancos centrales en aquel momento. El mantenimiento del patrón oro tenía prioridad absoluta en la aplicación de la política monetaria, tanto porque se llegó a ver este sistema como el fundamento de la estabilidad monetaria como porque no existían objetivos que compitieran con él (tales como el pleno empleo o el crecimiento económico) en la gestión de la política monetaria.
}

Por lo tanto, cabe preguntarse qué hacía de este sistema algo tan creíble. Dos son las razones principales:
\begin{lnum}
    \item \textbf{Desconocimiento teórico}. Las ideas importan aquí igual que en cualquier otra parte. El hecho que políticas monetarias y fiscales activas pudieran suavizar los ciclos económicos o que una devaluación de la moneda pudiera ayudar a reducir los desequilibrios comerciales, eran ideas que aún estaban por venir o que, en el mejor de los casos, eran heréticas. No existía una creencia extendida o una concepción bien articulada de cómo los gobiernos podían estabilizar la demanda, la producción o el empleo a través de instrumentos de Política Fiscal o Monetaria. El debate, tanto político como académico era la Político Comercial; y,
    \item \textbf{¿Bancos Centrales?}. A diferencia de los responsables de las políticas comerciales, los bancos centrales consiguieron aislarse del tira y afloja de la política nacional y pudieron operar de forma autónoma.\footnote{%
        EICHENGREEN, uno de nuestros más perspicaces historiadores de la globalización financiera, da directamente en el clavo: la capacidad de los bancos centrales para mantener la libre circulación del capital y los tipos fijos de cambio de divisas frente a las crisis económicas descansaba sobre las cortapisas a la presión política que hubiera podido ejercerse sobre ellos para perseguir otros objetivos incompatibles con la defensa de la convertibilidad del oro.
    } ¿Por qué? La razón es que el término Banco Central induce al error. Los emisores autorizados de moneda de las primeras potencias (Gran Bretaña, Francia y muchas otras) estaban, en realidad, en manos privadas y no tenían otra función pública que la emisión de moneda.\footnote{%
        Estados Unidos no tuvo una institución pública que actuara como banco central hasta 1913, cuando el Acta de la Reserva Federal creó el Sistema de la Reserva Federal.

    Estos titulares operaban como miembros de un club, con mayor afinidad entre ellos que con sus conciudadanos procedentes de ámbitos ajenos a las finanzas. En palabras de EICHENGREEN: El patrón oro era una institución construida socialmente.
    } 
\end{lnum}

De hecho, la capacidad de resistencia de esta institución se sometió a una dura prueba en la década de 1870, cuando una escasez de oro provocó, como requerían las reglas del patrón oro, una escasez de crédito y una deflación de los precios, tanto en Europa como en Estados Unidos. El golpe más fuerte fue para los campesinos, para quienes los altos intereses y la caída de los precios fueron devastadores. Hubo un gran clamor para volver a un estándar bimetálico que permitiera a los gobiernos monetizar la plata e incrementar la oferta de dinero. La revuelta llegó a su apogeo en Estados Unidos, donde William Jennings Bryan, tres veces candidato demócrata a la presidencia, pronunció su famoso discurso \textit{No crucificaréis a la humanidad en una cruz de oro} en el Congreso Nacional Demócrata de 1896. No obstante, los bancos centrales siguieron firmes y se mantuvo el patrón oro.\footnote{%
    Al final, lo que probablemente lo salvó fue que la deflación de los precios llegó a su fin cuando el descubrimiento de oro en Sudáfrica, después de 1886, derivó en un aumento de su oferta.
}

\subsubsection{Movimientos de Capital: Integración (casi)perfecta}
Los sólidos fundamentos del sistema monetario internacional sin duda contribuyeron a la movilidad de los flujos de capital; pues, la globalización financiera experimentada por la economía mundial en las décadas anteriores a la Primera Guerra Mundial fue, como mínimo, extraordinaria. 

En un pasaje que se reproduce en casi todos los libros sobre la globalización, KEYNES recordaba con nostalgia, en 1919, el hecho de que los habitantes de Londres podían invertir libremente su dinero en cualquier parte del mundo sin ningún impedimento o preocupación que pudiera privarles del fruto de tal inversión. En este periodo, los mercados financieros mundiales operaban con unos costes de transacción mínimos. Las tasas de interés en Londres, Nueva York y los principales centros financieros de Europa se movían al unísono, \textbf{como si formaran parte de un único mercado}. 

El capital \textbf{fluía con libertad y en enormes cantidades} desde donde lo había en abundancia (Gran Bretaña, en particular) hacia donde era escaso (el Nuevo Mundo, en particular). Además, a diferencia de lo que ocurrió con el libre comercio, no se abandonaron el oro ni los flujos libres del capital antes de que irrumpiera la Primera Guerra Mundial (incluso bajo intensas demandas durante las décadas de 1870 y 1880). Como hemos visto, entre las principales potencias económicas del momento la globalización financiera tan extraordinaria (que no se observaría hasta muy recientemente) era producto de creencias convergentes dentro de un club exclusivo de titulares de bancos centrales que tomaban todas las decisiones importantes. 

\paragraph{Proyección internacional: Imperialismo}
No obstante, entre los países de la periferia de la economía global, Latinoamérica, Oriente Próximo y Asia, la disciplina impuesta por medio de prácticas imperialistas fue una vez más crucial para mantener el libre flujo de capitales. La ortodoxia monetaria también se impuso, aunque no todos hicieran una transición completa al patrón oro. En este sentido, el gran quebradero de cabeza persistente (hasta hoy) de las finanzas internacionales concistía en saber cómo garantizar que los países, y los prestatarios dentro de ellos, fueran a pagar sus deudas. Cuando un prestatario de un determinado país se niega a realizar el pago de una deuda contratada localmente, la entidad crediticia agraviada puede ir a los tribunales, pedir que embarguen los activos del prestatario y esperar que las autoridades hagan cumplir la sentencia. Cuando un deudor extranjero no paga, el acreedor tiene pocas opciones. No existe un tribunal internacional que dicte sentencia ni una policía internacional que haga cumplir la decisión. En esencia, la única forma de control sobre el prestatario es la probable merma de su reputación y el coste potencial de ser excluido de los mercados crediticios internacionales durante un periodo de tiempo. A pesar de este coste para su reputación, la historia demuestra que los morosos acaban volviendo a entrar en los mercados financieros internacionales. Y esto tiene las siguientes consecuencias: (i) el prestatario podría dejar de cumplir sus obligaciones internacionales no solo cuando no pudiera pagar, sino también cuando no quisiera a pagar, un umbral mucho más bajo; y, (ii) si un banco, o un comprador de bonos, prevé esta posibilidad, es poco probable que llegue a hacer muchos préstamos internacionales, o los hará solo a un precio muy alto. 

En definitiva, las finanzas internacionales no pueden crecer a menos que haya mecanismos creíbles que garanticen los pagos. Así que allí estaba el \textit{imperium} para asegurar los potenciales beneficios \textit{para todos}. Algunos ejemplos clásicos son:\footnote{%
    Estados Unidos tuvo también una accidentada historia de pago de deudas, pues muchos de sus estados incumplieron sus obligaciones a lo largo del siglo XIX, de modo que es irónico que los estadounidenses acabaran siendo los garantes de las deudas en el hemisferio occidental.
}
\begin{lnum}
    \item Cuando el Imperio Otomano dejó de cumplir sus obligaciones con los titulares de bonos, en su mayoría británicos y franceses, en 1875, los europeos convencieron al debilitado sultán para que les permitiera establecer una agencia extraterritorial para recaudar los pagos vía impuestos otomanos. La Administración de la Deuda Pública Otomana (que empezó a funcionar en 1881) llegó a ser una vasta burocracia dentro del Estado otomano cuyo fin principal consistía en reembolsar el crédito extranjero;
    \item Cuando las revueltas nacionalistas en Egipto amenazaron los intereses económicos británicos en 1882, los británicos invadieron el país para \textit{restaurar la estabilidad política} y asegurarse de que se seguirían devolviendo las deudas contraídas. En aquel momento, el primer ministro, William Gladstone, tenía una buena parte de su riqueza invertida en obligaciones egipcias, de modo que, en este caso, el vínculo entre globalización financiera y poder militar fue particularmente transparente. Gran Bretaña acabó gobernando Egipto directamente, a pesar de que sus primeras intenciones habían sido mucho más limitadas; y,
    \item Theodore Roosevelt dejó claro en 1904 (en el llamado Corolario de Roosevelt a la Doctrina Monroe) que Estados Unidos garantizaría que los países latinoamericanos cumplieran con su deuda internacional. Mostró que lo decía en serio enviando cañoneras a Santo Domingo en 1905 y asumiendo la recaudación de ingresos aduaneros cuando la República Dominicana dejó de pagar su deuda una acción que indicaba su determinación de proteger los intereses de los acreedores extranjeros, e hizo que subieran los precios de los bonos estatales latinoamericanos⁠. La cuestión antes de que las cañoneras estadounidenses aparecieran no era si se cobraría o no la deuda, sino si serían los europeos o los estadounidenses quienes la cobrarían. Al adelantarse a los europeos, Roosevelt pretendía que no quedara duda de que se trataba de un área de influencia estadounidense.
\end{lnum}

El capitalista británico que invertía en los ferrocarriles indios sabía que el Raj Británico estaba allí para garantizar la seguridad de su inversión: \textit{Mientras se le garantizara el 5\% de sus ingresos en la India}, comentaba un funcionario británico, \textit{le tenía sin cuidado que los fondos que había prestado fueran a Hooghly o se convirtieran en ladrillos y mortero}. 

En definitiva, el patrón oro y la globalización financiera fueron posibles, al igual que en el caso del libre comercio, por una peculiar combinación de política nacional, sistemas de creencias y control por parte de terceros. Y el sistema de creencias que sostenía el patrón oro y la globalización financiera entre 1870 y 1914 no logró sobrevivir a los golpes mortales recibidos con la Gran Depresión y la revolución del pensamiento económico debida a John Maynard Keynes. Cuando estas fuerzas se debilitaron, al ir adquiriendo fuerza una política más dirigida a las masas, también se debilitaron las finanzas internacionales, lo que condujo al colapso último del patrón oro durante la década de 1930.

\subsubsection{Movilidad internacional de trabajadores: La verdadera globalización} 
La función principal de las migraciones internacionales en el siglo XIX fue redistribuir parte de la población de regiones agrícolas de Europa a nuevas regiones productoras de bienes primarios de ultramar, donde las condiciones físicas y sociales permitían obtener una mayor producción por trabajador de la que se habría obtenido en el país de origen. 

Muchos inmigrantes, a pesar de proceder de economías esencialmente agrícolas, contribuyeron al desarrollo de nuevos países ocupando puestos en sectores no agrícolas, liberando una parte adicional de la población nativa para expandir la agricultura.

Una minoría de migrantes poseían habilidades industriales y fueron responsables de laintroducción de nuevas técnicas de producción, pero la principal contribución de la inmigración al desarrollo industrial de los países de acogida fue a través del aumento de la población que provocó tanto directamente en un primer momento como posteriormente a través de las altas tasas de fertilidad que mantuvieron. Este crecimiento de la población proporcionó tanto la fuerza de trabajo como el mercado en expansión tan necesarios para la industrialización y la producción a gran escala. De los aproximadamente 46 millones de migrantes que existieron entre 1821 y 1915, en torno a 44 millones partieron de Europa y el resto principalmente de Asia. El grueso de las migraciones europeas tuvo lugar después de 1880 (32 millones), momento en el que también hubo un cambio significativo en el origen de las salidas de población: antes de 1880, la mayor parte de los emigrantes procedían del norte y oeste de Europa, mientras que después la mayoría procedían del sur y del este. A lo largo de todo el período, las Islas Británicas fueron la principal fuente de migrantes, registrando en torno al 37\% del total de emigrantes de Europa. Otros importantes países de origen fueron, en orden descendiente de importancia, Italia, Alemania, Austria-Hungría, España y Rusia. Del total bruto de inmigrantes registrados entre 1821 y 1915, en torno al 62\% entraron en Estados Unidos, 9\% a Argentina, 8\% a Canadá y aproximadamente el 7\% a Brasil. Del resto, la mayor parte consistió en flujos migratorios a a Australasia y a la zona templada de África.

 
\clearpage
\section{Estancamiento: periodo de entreguerras}
\puntosclave{
    La I Guerra Mundial tuvo consecuencias devastadoras sobre las economías tanto de los países vencedores como de los vencidos, particularmente en el caso de Alemania dadas las reparaciones de guerra exigidas tras la fima del Tratado de Versalles. 
    
    La tendencia hacia el libre comercio de la etapa previa a la guerra dio paso a movimientos de corte marcadamente proteccionistas, y la cooperación internacional en materia económica se volvió más complicada con ocasión del surgimiento de movimientos nacionalistas. 
    
    El intento de volver a una estabilidad en el plano monetario similar a la que existió durante el patrón oro a través del sistema cambios oro no fue fructífero, y su impacto podría decirse que fue incluso netamente desestabilizador.
}

Todo este breve repaso histórico, con sus más y sus menos, nos conduce a un suceso histórico de sobra conocido: la Primera Guerra Mundial (1914-1919). Por tanto, sería interesante preguntarnos si podemos encontrar sus causas en lo descrito hasta ahora.

Como hemos visto, la primera globalización ni era tan global ni estaba fundmentada sobre los argumentos liberales que se esgrimen habitualmente. Se trataba, por el contrario, de una situación frágil motivada por intereses de clase y mantenida mediante usos, más o menos frecuentes de la posición dominante. Se trataba, en esencia, de un equilibrio inestable y frágil. La situación podría haber desencadenado un final diferente, pero sin duda no podía persistr inmutable. La Primera Guerra Mundial generó elevadísimos costes desde un punto de vista político y económico. No obstante, tanto sus casuas como sus efectos más devastadores, que caracterizarán el periodo que vamos a describir son de vertiente social. 

El periodo de entreguerras no es un periodo homogeneo. Esta caracterizado fuertemente por un componente ciclo que se expande desde 1919 hasta 1939. Por ello es importante mantener una linea histórica coherente y ordenada, destacando los hechos más relevantes, sin perder nunca de vista las disparidades que experimentaron los territorios de las distintas potencias.

\textbf{¿Qué queremos decir?} Veremos en primer lugar el contexto geopolítico, explicando brevemente el impacto, político y económico, inmediato de la guerra y cómo se organizó el mundo en este periodo. Esto nos ayudará a entender tanto las etapas posteriores como las causas que condujeron a la II Guerra Mundial (sin entrar en detalle). 

\subsection{¿Cuál es el contexto geopolítico?: Las secuelas de la Guerra}
La I Guerra Mundial tuvo un impacto fortísimo en muchos ámbitos. 

\subsubsection{Reconfiguración política: Fragmentación de mercados}
Empecemos primero con la reconfiguración del orden geopolítico. El conflicto con profunda repercusión se dió entre dos bandos:\footnote{%
    Estas alianzas no solo respondían a intereses militares, sino también a necesidades políticas y a la posición que cada país quería ocupar en el equilibrio europeo.
} (i) los Imperios Centrales (Alemania y el Imperio Austro-Húngaro) y el Imperio Otomano, que pretendían mejorar su posición frente a las potencias coloniales;\footnote{%
    Italia nunca se sintió cómoda del todo en ese bloque. Las tensiones históricas con Austria-Hungría por el control de zonas en los Alpes seguían vivas. Aunque formaba parte de la alianza, en 1902 firmó un compromiso secreto con Francia: si uno era atacado, el otro no atacaría.

El tratado se renovó varias veces, pero sin reforzar la cohesión entre sus miembros. Cuando comenzó la guerra en 1914, Italia decidió mantenerse al margen. Dos años después, en 1915, se pasó al bando contrario tras recibir garantías de que podría anexionar territorios en disputa.
} y, (ii) la Triple Entente (Gran Bretaña, Francia y el Imperio Ruso) y Estados Unidos, que intentaban frenar esta expansión y/o defender sus intereses político-económicos.

\imagenfit{GranGuerra.png}{La Gran Guerra: Alianzas}{\href{https://historicodigital.com/la-gran-guerra.html#google_vignette}{\textit{Histórico Digital}. BLOG (2011)}}

El desenlace de la Guerra es desobra conocido.\footnote{%
    Se puede consultar la cronología de la Guerra en [\authorfont{Ver Anexo \ref{anx:gran_guerra}}]
} Las potencias aliadas se declararón vencedoras y se procedió a represaliar a los Imperios atacantes. No obstante, estas represalias no se limitaron a sanciones económicas, que luego veremos, sino que supusieron una verdadera transformación sin precedentes de las fronteras políticas del mundo hasta entonces conocido. El Imperio Otomano, el Imperio Ruso y, en último término, el Imperio Austro-Hungaro fueron desmembrados por diferentes causas. Solamente en Europa, se constituyeron más de diez nuevas naciones soberanas por causas diversas. El Imperio Otomano, con su basta extensión en Oriente Próximo y Norte de África fue repartido entre las potencias coloniales vencedoras (Francia y Gran Bretaña), estableciendose protectorados en territorios que hoy en día persisten (Irak, Palestina, Líbano, etc.). El Imperio Alemán perdió la posesión de todas sus colonias (más de dos millones de kilómetros cuadrados), que se traspasaron a manos británicas. En conjunto, la reorganización creó o redefinió fronteras sobre varias decenas de miles de kilómetros.\footnote{%
    El proceso estuvo formalizado, sobre todo, en los tratados de: (i) \href{https://es.wikipedia.org/wiki/Tratado_de_Brest-Litovsk}{Tratado de Brest-Livostk, 3 de marzo de 1918)}, entre el Imperio Alemán y el Imperio Ruso; (ii)  \href{https://es.wikipedia.org/wiki/Tratado_de_Versalles_(1919)}{Tratado de Versalles, 28 de Junio de 1919}, entre las Potencias Aliados y la República de Weimar; (iii) \href{https://es.wikipedia.org/wiki/Tratado_de_Saint-Germain-en-Laye}{Tratado de Saint-Germain-en-Laye, 10 de Septiembre de 1919}, entre las Potencias Aliadas y Austria; (iv) \href{https://es.wikipedia.org/wiki/Tratado_de_Trianón}{Tratado de Trianon, 4 de Junio de 1920}, entre las Potencias Aliadas y el Reino de Hungría; y, (v) \href{https://es.wikipedia.org/wiki/Tratado_de_Sèvres}{Tratado de Sèvres, 10 de agosto de 1920}, entre las Potencias Aliadas y el Imperio Otomano), parcialmente revisado luego por el \href{https://es.wikipedia.org/wiki/Tratado_de_Lausana}{Tratado de Lausana, 24 de julio de 1923}, entre las Potencias Aliadas y el recién Estado de Turquía. 

Los nuevos Estados surgidos del desmembramiento de los imperios europeos continentales fueron: (i) del Imperio Ruso, Polonia, Finlandia, Letonia, Estonia, Lituania, Bielorrusia, Ucrania y Moldavia (sin delimitarse tal y como están hoy en día); (ii) del Imperio Austro-Hungaro, Austria, Hungría, Checoslovaquia y Yugoslavia; (iii) del Imperio Otomano, Irak, Palestina, Jordania, Siria, Líbano, Yemen, Egipto y Arabia Saudita. Además, poco después también se independizó Irlanda, debido a la debilidad de la Gran Bretaña. También hubo ampliaciones territoriales importantes de Rumanía, Italia y Grecia.

El Imperio Austrohúngaro tenía una superficie en 1914 de $\sim$ 676.000 km$^2$. Tras la guerra quedó fragmentado en al menos siete unidades políticas relevantes y se trazaron aproximadamente entre 12.000 y 20.000 km de nuevas fronteras internacionales, dependiendo del criterio cartográfico utilizado. El Imperio Alemán perdió unos 70.000 km$^2$ en Europa (Alsacia-Lorena, Posnania, corredor polaco, Eupen-Malmedy, Schleswig septentrional, etc.), además de todas sus colonias, unos $\sim$2,9 millones de km$^2$. El Imperio Ruso perdió aproximadamente 1 millón de km$^2$ en Europa occidental y báltica y surgieron Finlandia y los Estados bálticos, mientras Polonia recuperó extensos territorios. El Imperio Otomano perdió prácticamente todos sus territorios árabes. Antes de la guerra controlaba unos 5,2 millones de km$^2$; tras el proceso quedó reducido esencialmente a Anatolia y Tracia oriental, base de la futura Turquía. Quedaron en manos del mandato británico: Irak, Palestina y Jordania; y bajo mandato Francés: Siria y Líbano. Además, sufrío las pérdidas (previas o simultaneas) de los territorios de Arabia Saudita, Yemen y Egipto.

En total, se estima que el Reino Unido obtuvo control directo o indirecto sobre $\sim$1,8 millones de km$^2$, Francia administró $\sim$790.000 km$^2$ en Siria y Líbano e Italia y Grecia aspiraron a zonas de Anatolia bajo el Tratado de Sèvres, aunque gran parte se perdió tras la Guerra de Independencia Turca. En suma, desaparecieron cuatro imperios multinacionales, se consolidó el \href{https://es.wikipedia.org/wiki/Wilsonianismo}{principio wilsoniano} de autodeterminación (aunque aplicado de manera desigual) y aparecieron numerosos Estados multiétnicos inestables, formando la base geopolítica de muchos conflictos, sucesivos y contemporáneos, en Europa oriental y Oriente Próximo.
}

Esto tuvo un impacto decisivo en lo que seguirá del periodo. Aunque pueda parecer tremendamente costoso para unos y beneficioso para otros, la realidad es que ninguno surgió beneficiado del nuevo reparto colonial. Los Imperios desmembrados sufrieron pérdidas que, salvo Alemania, no eran capaces de estabilizar en los años previos; mientras que las Potencias Alidas, en particular Francia y Gran Bretaña, absorbieron tales territorios en unas condiciones que imposibilitaban su control como venía haciéndose sobre los territorios coloniales precedentes. Desgastados por la Guerra y sufriendo las gravosas deudas económicas que habían contraído, las posesiones coloniales recién adquiridas desgastaron progresivamente sus capacidades, muy mermadas, dificultando el control de territorios previamente sometidos (por ejemplo, los territorios del Subcontinente Indio).

De esta forma, la reconfiguración político tuvo un impacto económico tremendo que se magnificó por dos vías: (i) los territorios recién independizados fragmentaron el mercado europeo en muchos más territorios de los que habían existido previamente, con el subsecuente incremento de los costes de transacción que esto supone; y, (ii) la ampliación de mercados y territorios coloniales de las Potencias Aliadas no compensaron los recursos que fue necesario destinar para su control y sometimiento, lo que acabó generando un efecto desbordameinto sobre otras regiones coloniales que se sumirían en el caos.

El mapa político de Europa y Medio Oriente se redefinió una vez terminada la guerra, y hubo cuatro grandes perdedores en materia territorial: Alemania, el Imperio Austro-Húngaro, el Imperio Otomano y Rusia. En Europa aparecieron 12 nuevos países independientes. Alemania perdió el 13\% de su territorio, perdiendo regiones estratégicas a manos de Francia y Polonia. El viejo imperio ruso, perdió los territorios de las actuales Finlandia, Estonia, Letonia, Lituania y Polonia. El Imperio Austro-Húngaro desapareció, y de sus territorios se formaron Austria (reducida), Hungría (reducida), Checoslovaquia y Yugoslavia; parte del territorio pasó a manos polacas. El Imperio Otomano tuvo una suerte similar: el Estado Turco quedó reducido a su tamaño actual y perdió todos sus dominios en Medio Oriente y Arabia; su fragmentado territorio se convirtió en un polvorín hasta el presente.

Toda esta redefinición de fronteras, importante en sí misma, es clave para entender la inestabilidad política de las regiones mencionadas durante todo el siglo siguiente. Por otra parte en términos económicos, significó un enorme retroceso: en dónde antes hubo mercados nacionales integrados pasó a haber importantes barreras al comercio y a la integración. Se desintegraron mercados en tiempos en que había que reconstruir las economías luego de una costosa guerra con altas perdidas de capital humano y físico. La crisis económica era inminente.

\subsubsection{Socialismo: el (único) vencedor de la Guerra}

En este contexto creció la adhesión al comunismo, que había triunfado en Rusia en 1917, y a los movimientos nacionalistas, que se alimentaban del resentimiento por el trato humillante ofrecido a algunos de los países tras la Primera Guerra Mundial (como fue el caso del fascismo italiano, que llegó al gobierno en 1922 e inspiró al nazismo alemán). Una expresión del conflicto político e ideológico del período fue la guerra civil española (1936-1939), que enfrentó a sectores republicanos, socialistas, comunistas y anarquistas con el bando sublevado o “nacional” que contó con el apoyo del fascismo y el nazismo. Este conflicto concluyó pocos meses antes del inicio de la Segunda Guerra Mundial.

El 11 de noviembre de 1918 el kaiser Guillermo II, que llevó a Alemania a la guerra, huyó hacia el extranjero presionado por la sublevación de marineros en el puerto alemán de Kiev, misma que fue secundada con disturbios en otras ciudades alemanas. Al siguiente mes Rosa Luxemburgo y Karl Liebknecht declararon la República Socialista Alemana (Childs, 1980: 13-14). En Hungría, en marzo de 1919 el régimen comunista de Bela Kun expropió la propiedad privada y tomó el control de los bancos y el comercio (Willmott. 2003: 301). Al concluir la guerra, los trabajadores no solo alemanes sino de toda Europa no estaban dispuestos a regresar a las mismas condiciones laborales y de vida prevalecientes antes de la guerra, se sintieron motivados a transformar el sistema siguiendo el ejemplo de la Revolución rusa de apenas un año antes. Para los trabajadores resultaba muy atractivo terminar en definitiva con las extremadamente difíciles condiciones de vida y de trabajo descritas en el apartado anterior, por lo tanto la burguesía tenía que apresurarse para evitar el avance del socialismo no solo en Alemania, sino en el resto de Europa. Al terminar la guerra, las huelgas proliferaron en Inglaterra incluyendo Irlanda del Norte, el número de huelgas pasó de 730 en 1917, a 1165 en 1918 y todavía en 1919 aumentaron las huelgas a 1 607, posteriormente disminuyeron hasta incrementarse de nuevo en los años cincuenta y sesenta (Mitchell, 1988: 142). Apenas cinco meses después de concluida la guerra, en abril de 1919, como parte del Tratado de Versalles, que dio fin a la guerra, se fundó la Organización Internacional del Trabajo (oit), provocando un sustancial aumento de salarios en los países que ahora forman el grupo de los países más ricos (Douglas, 1966). Tan solo en los Estados Unidos, en 14 años, de 1890 a 1919, el salario aumentó apenas 14%, pero en cuatro años, de 1920 a 1924, el salario se incrementó 24%, posteriormente se estancó en tanto la productividad continuó creciendo, provocando así la crisis de 1929 (gráfica 5.1).




\subsection{Consecuencias económicas de la Guerra: el coste de la guerra}
A esta reconfiguración política del panorama internacional que, como hemos visto, dista mucho de generar unas condiciones óptimas para el crecimiento en los años posteriores, deben sumarse las consecuencias estríctamente económicas que supuso la guerra. STAMP calculó que la guerra destruyó el crecimiento normal de unos tres o cuatro años de las rentas derivadas de la propiedad en Europa o una trigésima parte de su valor original.

Aunque se suele decir que la guerra tuvo unas consecuencias económicas diferentes para los vencedores que para los vencidos, no es del todo correcto. Las consecuencias económicas directamente derivadas de la Guerra (situación económica a 1919) fueron muy similares para todos los países a excepción de los Estados Unidos y algunos países de la periferia (no colonial). 

Siguiendo al historiador económico español TAFUNELL (2005), podemos decir que el conflicto bélico tuvo cinco tipos de efectos económicos, manifestados a través de tres canales claramente diferenciados.

\subsubsection{Capital humano: Pérdida de población}
En primer lugar, la pérdida de capital humano.

Es difícil hacer un cómputo exacto de las pérdidas de población producidas por la guerra. Esto es así, en parte, porque los datos del período distan mucho de ser perfectos, pero también porque las bajas militares tan sólo fueron una pequeña proporción del número total de muertes ocurridas en ese período. Murió mucha más gente de hambre y enfermedades, o como consecuencia de la guerra civil rusa, que en el campo de batalla. Además, debe hacerse alguna estimación del déficit de población causado por la falta de nacimientos como consecuencia de las condiciones del tiempo de guerra. En base a los datos recopilados por (ALDCROFT, 1990):

\begin{center}
\begin{tabular}{|p{0.15\linewidth}|p{0.20\linewidth}|p{0.075\linewidth}|p{0.175\linewidth}|p{0.075\linewidth}|p{0.175\linewidth}|}
\hline
\textbf{Potencia} & \textbf{Acciones militares} & \textbf{Civiles} & \textbf{\% Población Activa} & \textbf{Déficit} & \textbf{\% Población Total} (est.) \\
\hline
\textbf{Francia} & 1,4 M & \multirow{5}{*}{5 M} & 10\% & 3 M & 7\%\\
\cline{1-2}\cline{4-6}
\textbf{República del Weimar} & 2 M & & 10\% & 5 M & 7,5\% \\
\cline{1-2}\cline{4-6}
\textbf{Reino Unido} & 1,4 M &  & 5\% & 1,5 M & 4\%\\
\cline{1-2}\cline{4-6}
\textbf{Imperio Austro-Húngaro} & 1,4 M &  & 5,7\% & 5 M & 6,6\%\\
\cline{1-2}\cline{4-6}
\textbf{Italia} & 1,4 M &  & 6\% & 2,3 M & 6,1\% \\
\hline
\textbf{Europa} & 1,4 M & 5,5 M & 8\% & 22-24 M & 7\% \\
\hline
\textbf{Rusia} & \multicolumn{2}{c|}{16 M} & 3,3\% & 26 M & 20\% \\
\hline
\end{tabular}
\end{center}

Como vemos, las bajas militares fueron muy pequeñas en términos relativos. Durante el período de hostilidades, unos 8,5 millones de hombres (incluyendo estimaciones aproximadas para Rusia) perdieron su vida en servicio activo (un 15\% de los movilizados) equivalente al 8\% trabajadores varones. Además, debemos incluir unos 7 millones de hombres que quedaron incapacitados permanentemente, y otros 15 millones más o menos seriamente heridos.\footnote{%
    La mayoría de las personas muertas estaban en la flor de su vida y por tanto constituían la parte más productiva de la fuerza de trabajo. En el caso de Alemania, el 40\% de las bajas estaba dentro del grupo de edad de 20-24 años y el 63\% entre 20 y 30 años. 

Las cifras para Rusia son mucho menos fiables, aunque es probable que las pérdidas en este país superen la cifra total del resto de Europa. Las bajas militares en la Gran Guerra estricta fueron relativamente pequeñas, pero murieron millones en la revolución y en la guerra civil que vinieron a continuación. El número total de víctimas no estuvo muy lejos de los 16 millones. Añádanse a esta cifra unos 10 millones por déficits de nacimientos y se llega a 26 millones, sin tener en cuenta las pérdidas en los territorios cedidos por Rusia como consecuencia del tratado de paz concluido con Alemania en 1918.
}

En definitiva, se perdió por muertes directas y el déficit de natalidad el equivalente al 7\% de la población europea de antes de la guerra, o al conjunto de su crecimiento natural entre 1914 y 1919. Así, \textbf{a principios de 1920 la población de Europa era aproximadamente la misma que al principio de la guerra}. Además, nótese que se incluyen los Estados ex-post, puesto que debemos tener en cuenta que la desmembración de algunos territorios también generó déficits desde una perspectiva ex-ante han sido contabilizados. Europa sufrió un serio agotamiento y deterioro en la calidad de la población durante el período de la guerra. Las cifras citadas, además, no son estrictamente completas, porque hubo pérdidas adicionales, debidas a causas asociadas con la guerra, que se produjeron en el período posterior al armisticio. La epidemia de gripe de 1918-1919 se cobró muchas víctimas, mientras que un número sustancial de personas murieron en Europa oriental y en los Balcanes a consecuencia del hambre. Conflictos fronterizos posbélicos y matanzas entre 1919 y 1921, especialmente en el sudeste de Europa, añadieron más víctimas al total. 

Una cifra aproximada sería entre 50 y 60 millones, de los que Rusia supondría aproximadamente la mitad. Las muertes militares directas en la guerra estricta representaron sólo una pequeña proporción de las pérdidas totales de población en este período. En términos humanos el desastre puede considerarse poco menos que trágico. Pero es dudoso que la pérdida tuviera un impacto fuerte y permanente en los países afectados. La mayoría de los países, por supuesto, perdieron parte de su mejor mano de obra, a menudo altamente cualificada, pero pocos, aparte de Francia, sufrieron escasez de trabajadores en la década siguiente a la guerra. En efecto, como así sucedió, el período posbélico estuvo marcado por un alto desempleo en muchos países europeos y así puede argumentarse que el freno al crecimiento de la población fue algo así como una especie de bendición.

\subsubsection{Capital físico: Destrucción de infraestructura y doctrina \textit{guerra total}}
La razón que lo explica es doble. 

Por un lado se debe a la destrucción de capital físico que había generado la Guerra. De manera directa o indirecta, supuso la destrucción de la muy diversa infraestructura productiva generada en las décadas precedentes. El valor del stock de capital de Europa se deterioró durante la guerra como consecuencia del daño físico, la venta de activos extranjeros, el freno a la inversión y el descuido en el mantenimiento. Veamos lo por partes.
\begin{la}
    \item \textbf{Impacto destructivo}. Europa perdió aproximadamente una trigésima parte de sus activos fijos como consecuencia de la destrucción y daño físico. Además, por supuesto, algunos países sacrificaron territorio y propiedad por los acuerdos del tratado de paz. No obstante, la incidencia del impacto destructivo varió considerablemente de país a país: el daño físico fue máximo en los principales teatros bélicos, especialmente en Francia y Bélgica, aunque a Italia, Rusia y algunos países de Europa oriental también les fue mal;\footnote{%
        Fue inevitable que Bélgica y Francia soportaran la carga principal, dado que buena parte de la lucha se produjo en sus territorios. La destrucción de granjas, fábricas y casas fue amplia y sustancial en ambos casos, aunque en Francia la mayor parte del daño físico tendió a concentrarse en el norte del país. (Aunque geográficamente concentradas, las pérdidas de Francia fueron severas y se produjeron en la parte más rica y avanzada del país). No obstante, Bélgica fue menos afortunada. Prácticamente todo el país fue invadido y la lista de daños hace lúgubre su lectura: un 6\% de los edificios, la mitad de las acerías y las tres cuartas partes del parque móvil ferroviario fueron destruidos o dañados sin reparación posible, miles de acres de tierra se convirtieron en inservibles para el cultivo, mientras que la población animal fue diezmada. En términos absolutos estos dos países representaron el grueso de las pérdidas de la propiedad en el tiempo de la guerra. No obstante, en términos relativos, es probable que algunos de los países más pequeños situados más al este, salieran de la guerra en condiciones de una mayor devastación.
    } en comparación, Gran Bretaña, Austria y Alemania, aunque beligerantes principales, fueron castigados con bastante poca severidad. Indudablemente, los territorios ocupados corrieron la peor suerte;\footnote{%
        El valor de la propiedad perdido por Polonia fue sólo un poco menor que el de Alemania, pero el impacto fue mucho mayor. Las potencias ocupantes devastaron literalmente el país por la destrucción y el saqueo. Grandes extensiones de tierra agrícola fueron dejadas yermas, el 60\% de la cabana ganadera desapareció, gran parte del parque móvil ferroviario fue requisado, muchas fábricas fueron destruidas o despojadas de sus equipos, y 1,8 millones de edificios se perdieron por el fuego. Lo mismo puede decirse de Serbia, de partes de Austria y también de Rusia, aunque en el último caso gran parte del daño se produjo a consecuencia de la guerra civil. De hecho, en algunas áreas la escala de la destrucción fue tan grande que la cuestión de la reparación difícilmente podía considerarse; más bien era cuestión de limpiar la tierra y reconstruir de nuevo. Bulgaria lo hizo mucho mejor que sus vecinos de la península balcánica, porque el país nunca se convirtió en zona de guerra y así evitó una destrucción fuerte o el despojo de la propiedad
    }
    \item \textbf{Política de guerra total: Deterioro de capital físico}. A esto debe añadirse una cantidad desconocida por el deterioro del stock de capital existente, debido al descuido o falta de mantenimiento. La mayoría de los países beligerantes experimentaron cortes sustanciales en la inversión, con el resultado de que sus stocks de capital eran menores al terminar las hostilidades. Además, la conversión de las economías europeas a \textbf{economías de guerra} desequilibró la producción a escala nacional (y global) e implicó la reasignación de grandes cantidades de recursos hacia sectores estratégicos para sostener el esfuerzo bélico (con el consecuente auge de las importaciones, aumento del déficit comercial y descuido de otra infraestrutura productiva no estratégica), afectando seriamente la capacidad de recuperación y crecimiento económico de los países europeos. Se estima que, en los países beligerantes, el gasto del Estado acaparó entre el 30 y el 50\% del producto bruto (COMÍN, 2014);
    \item \textbf{Descapitalización de activos en el extranjero}. Por último, merce especial atención la descapitalización de activos mantenidos en el extranjero. Alemania perdió pocos activos nacionales, pero la mayor parte de sus activos exteriores fueron vendidos o embargados. La mayor parte de las pérdidas físicas de Gran Bretaña fueron buques, pero también vendió una pequeña parte de sus inversiones ultramarinas para pagar la guerra. Francia perdió dos terceras partes de sus activos exteriores de antes de la guerra, por venta, insolvencia (como en el caso de las inversiones en Rusia) o a causa de la inflación. 
\end{la}

En este contexto, la tarea de reconstrucción era ciertamente sustancial y en algunos países. Sin embargo, el proceso de recuperación europeo se vió definitivamente condenado por dos una última pieza del puzle.

\subsubsection{Secuelas financieras: ¿cómo (no) financiar la guerra?}
La dimensión del programa de gasto bélico (estimado entre 250.000-300.000 millones de dólares, incluyendo todos los beligerantes) no fue particularmente significativo.\footnote{%
    Los mayores gastos correspondieron a Gran Bretaña, Estados Unidos, Alemania, Francia, Austria-Hungría e Italia, en este orden. Puede obtenerse una idea de la magnitud del desembolso total por el hecho de que representó unas seis veces y media la suma de toda la deuda nacional acumulada en el mundo desde el final del siglo XIX hasta la víspera de la primera guerra mundial.
} La I Guerra Mundial proporciona un gran ejemplo de como no financiar un esfuerzo bélico: casi de la noche a la mañana los gobiernos abandonaron precipitadamente la sana ortodoxia financiera del siglo XIX, abandono de la disciplina del patrón oro y el recurso a la financiación con déficit. 

Las operaciones de crédito de una u otra clase, más que los impuestos, fueron la principal fuente de financiamiento de la guerra. Alemania y Francia, por ejemplo, confiaron casi por completo en el préstamo. En Estados Unidos sólo un poco más del 23\% de los gastos de guerra se obtuvo de fuentes de renta. En promedio, un 80\% o más del gasto total de guerra de los beligerantes se financió por medio de préstamos. La deuda pública en los países europeos se disparó. Además, buena parte de esta deuda fue colocada en el interior de cada país, aprovechando el auge del sentimiento patriótico y nacionalista de la población además de un claro convencimiento de brevísima duración (la Guerra que todos iban a ganar/nadie iba a perder). No obstante, estos préstamos no se obtuvieron de auténticos ahorros (como sí sucedió a menudo en la II Guerra Mundial, \href{https://es.wikipedia.org/wiki/Bono_de_guerra}{Bonos de Guerra de EEUU}), sino que procedió principalmente del crédito bancario concedido por: (i) creación de nuevo dinero; y/o, (ii) aumento de la oferta monetaria a cmbio de \textit{promesas de pago} que utilizaron como reservas (irónico). Aunque los detalles del mecanismo variaron de un país a otro, el resultado final fue con mucho el mismo: \textbf{las deudas públicas aumentaron rápidamente, incrementándose la proporción de la deuda a corto plazo a medida que pasaba el tiempo; la oferta monetaria aumentó considerablemente y las reservas metálicas de los bancos, en relación con el pasivo, cayeron notablemente}.\footnote{%
    A fines de 1918 la oferta monetaria alemana había aumentado nueve veces y el déficit presupuestario seis veces, mientras que la relación entre las reservas metálicas y los billetes de banco y depósitos había bajado del 57\% al 10\%. La situación fue incluso peor en el caso del imperio Austro-Húngaro, mientras que Francia y Bélgica también lo pasaron mal. En general, la degradación de las condiciones financieras fue más aguda en los países de Europa central, menos aguda en los países neutrales y de moderado alcance en los demás.
}

El desastre era inminente -y evidente, añadiríamos-. Tales condiciones automáticamente dieron lugar a inflación de precios y depreciación monetaria, dado que casi todos los países congelaron la convertibilidad durante la guerra (muchos países abandonaron directamente el patrón oro durante la Guerra). La mayoría de los países experimentaron aumentos del doble al triple de sus precios, y en algunos casos mucho más, dependiendo del grado de inflación monetaria. Por ejemplo, los precios al por mayor en Alemania al final de las hostilidades eran cinco veces los de antes de la guerra, mientras que el marco había bajado al 50\% de su antiguo valor. Austria y Hungría experimentaron una inflación todavía mayor con valores monetarios cayendo al 30 y 40\% de la paridad original. Otros países cuyas monedas habían empezado a depreciarse significativamente eran Finlandia, Francia, Italia, Bélgica y Portugal.\footnote{%
    La mayoría de los neutrales, por otra parte, intentaron mantener o mejorar el valor de sus monedas, a pesar de grados significativos de inflación.
} Además, durante el primer año de paz continuaron las descuidadas políticas monetarias y fiscales.

Se ha comentado lo anterior porque, a menudo, se atribuyen las consecuencias desastrosas del periodo inmediatamente posterior al grave régimen de sanciones que se impuso sobre los países perdedores. Nada más lejos de la verdad. La situación de todos los países europeos, en términos financieros, era muy, muy problemática. No obstante, los problemas inflacionistas no eran, en su mayor parte, irresolubles (con las posibles excepciones de los de Alemania y Austria-Hungría). 

Sin embargo, el golpe final sí que, de forma más o menos acentuada, fue dado por las sanciones y la falta de \textit{coordinación/amabilidad} de las Políticas orientadas al cumplimiento de objetivos estríctamente nacionales.

En definitiva, las implicaciones financieras de la primera guerra mundial fueron más serias que las de la segunda y ello se debió a que el método de control financiero fue mucho menos exigente en el primer caso, y no a que la escala de las operaciones militares fuese mayor; de hecho fue al revés. 

El cambio de política se produjo en 1920, cuando varios países, en particular Estados Unidos, Gran Bretaña y Suecia, impusieron una fuerte política de economías, que afectó adversamente a sus economías nacionales e hizo mucho más difícil la recuperación en el resto de Europa. La mayoría de países europeos, sin embargo, continuaron con inflación, con consecuencias desastrosas para algunos, en especial Alemania, Austria, Polonia y Hungría. 

\begin{lnum}
    \item \textbf{Sistema Monetario Internacional}. La inestabilidad monetaria dificultó el proceso de recuperación, mientras que el fracaso de los gobiernos en aceptar el declive de los valores monetarios y el abandono del oro como algo sólo temporal, con el tiempo llevó a un intento desorganizado de tratar de restaurar el sistema monetario de antes de la guerra en condiciones completamente diferentes. En tercer lugar, debe hacerse referencia a la compleja serie de deudas intergubernamentales contraídas entre los aliados y las cargas por reparaciones impuestas a los vencidos, que no sólo demostraron ser una fuente de fricciones internacionales a lo largo de los años veinte, sino que también impidieron el proceso de reconstrucción financiera.
\end{lnum}


\subsection{Consecuencias económicas de la paz: el desorden mundial}
Si hasta aquí hemos visto las consecuencias económicas de la Guerra, veamos brevemente las consecuencias económicas de la paz. 

El período de entreguerras fue la franja de tiempo que transcurrió entre el final de la Primera Guerra Mundial (1914-1918) y el comienzo de la Segunda Guerra Mundial (1939-1945). Distinguimos tres etapas en este periodo: (i) una primera etapa de colapse económico e intento de reconstrucción (1919-2024); (ii) una segunda etapa de progresiva recuperación económica y distensión política que incrementó el consumo de bienes y entretenimiento (1924-1929); y, (iii) una tercera etapa caracterizada por la Gran Depresión y su contagio internacional. Veamoslas con detalle.

\imagenfit{Crono_Entreguerras.png}{Cronología del periodo de entreguerras}{\href{https://humanidades.com/periodo-de-entreguerras/}{Enciclopedia humanidades}}

\subsubsection{Primera etapa: Sembrando en invierno}
Las reparaciones de guerra son, sin duda, el aspecto central de las consecuencias económicas de la paz y la primera etapa del periodo de entreguerras.\footnote{%
    Este término fue acuñado por el reconocido economista John Maynard Keynes, quien criticó duramente las sanciones impuestas por los países vencedores de la guerra a las naciones vencidas. Los tratados de paz, firmados por el conjunto de los países vencedores con, separadamente, Alemania, Austria, Hungría, Turquía y Bulgaria, tendrían, como auguró Keynes, nefastas consecuencias a corto y largo plazo. El Tratado de Versalles, reconocía a Alemania como culpable de la guerra y le obligaba a compensar a los países vencedores por los daños causados. Estos daños se cifraron en torno a 31.500 millones de USD (132.000 millones de marcos) más el 25\% de los aranceles a la exportación durante 42 años. El problema fundamental asociado a estos pagos era su transferencia a divisas de pago, ya que Alemania no disponía de reservas suficientes y, dada la inmediatez con que se hacía necesario el pago, no podía depender de la generación de superávits continuos de balanza de pagos para ir acumulando divisas. El recurso utilizado fue el acceso a créditos internacionales, que debilitaron aun más a la economía alemana. Estos pagos y el aumento de medidas proteccionistas (ver más adelante) llevaron a Alemania a una situación insostenible de balanza de pagos, que condujo a la devaluación del marco en 1922 y alimentó la espiral inflación-devaluación. Además, al ser estos pagos parcialmente satisfechos en especie, provocaron convulsiones en numerosos sectores económicos, especialmente los extractivos, en los países aliados.
} Sin embargo, no explican todo el fondo del asunto. 

Al terminar la Primera Guerra, a excepción de los Estados Unidos, los países participantes enfrentaban condiciones económicas extremadamente difíciles: (i) millones de personas regresaron del frente de batalla desempleadas, miles de ellas lisiadas o incapacitadas con apremiantes necesidades sanitarias, asistenciales, alimentarias y de todo tipo; (ii) carestía generalizada, mercado negro y escasez de productos básicos (DOUGLAS, 1966); (iii) los gobiernos enfrentaban deudas sumamente elevadas contraídas para financiar la guerra; (iv) pérdida generalizada del poder adquisitivo. Sin embargo, nada era visto con tanto tempor y animádversión por parte de la clase político como nuestro querido vencedor: el socialismo. 

El final de la Primera Guerra Mundial marcó un parteaguas en la historia del capitalismo que obligó a las empresas a romper las viejas prácticas laborales seguidas hasta entonces. De forma generalizada, Gobiernos y productores, pese a la catástrofe económica en que se encontraban, se entregaron a las concesiones sociales con el fin de evitar la deblace del sistema. Apenas cinco meses después de concluida la guerra, en abril de 1919, como parte del Tratado de Versalles, que dio fin a la guerra, se fundó la Organización Internacional del Trabajo (OIT), provocando un sustancial aumento de salarios en los países que ahora forman el grupo de los países más ricos (DOUGLAS, 1966). Tan solo en los Estados Unidos, en 14 años, de 1890 a 1919, el salario aumentó apenas 14\%, pero en cuatro años, de 1920 a 1924, el salario se incrementó 24\%.

\imagenfit{Fig1.png}{Estados Unidos, Índice de salario real en la industria manufacturera, 1914=100}{Douglas, Paul H. Real Wages in the United States, 1890-1926, Nueva York, Cuadro 24, p. 108.}

Pocos tienen en cuenta las reformas laborales adoptadas en el Tratado de Versalles. En su articulado (art. 427) expuso las nuevas condiciones laborales que deberían de imponerse a nivel universal:\footnote{%
    Con el propósito de asegurar la implementación en todo el mundo de las condiciones universales de trabajo, el Tratado de Versalles creó, art. 387, a la Organización Internacional del Trabajo (OIT), encargada de garantizar que los países cumplieran el compromiso de otorgar condiciones remunerativas de trabajo. Llama la atención el discurso humanitario empleado por los diseñadores del Tratado de Versalles para justificar la creación de la OIT, cuando unos años antes se había perpetrado todo tipo de represión y crímenes contra los trabajadores, que las más de las veces terminaban en masacres.
}
\begin{lnum}
    \item El trabajo no puede ser considerado fácilmente como una mercancía o como un artículo comercializable;
    \item El derecho a la asociación tanto de empleados como de empleadores;
    \item El pago de un salario al trabajador que le asegura condiciones de vida apropiadas de acuerdo con las condiciones de su tiempo y su país;
    \item La adopción de una jornada de trabajo de ocho horas o 48 horas de trabajo a la semana;
    \item La adopción de un día de descanso a la semana que incluya por lo menos 24 horas;
    \item La abolición del trabajo infantil y de la responsabilidad dañina para el trabajo de jóvenes;
    \item Asegurar salario igual para hombres y mujeres para trabajos similares.
\end{lnum}

Lo anterior no tenía el afán de mejorar las condiciones de vida de los trabajadores, pues nunca habían sido importantes en la historia de la humanidad, y el capitalismo no era la excepción, su intención fue más bien frenar el avance del socialismo, eliminar el descontento de los trabajadores y motivarlos a regresar a las fábricas convencidos de las bondades del sistema. En este cotexto, al terminar la I Guerra Mundial se reconoce que la existencia de entornos mínimos de bienestar y de seguridad del trabajo eran condición necesaria para garantizar la paz mundial.

Lo que cabría preguntarse entonces es: ¿Se trata de una situación óptima en la aplicación de dichas Políticas? Aunque su proyección real dista mucho de lo estipulado, la realidad de la época se entiende muy bien con lo expuesto. Aquellos países con el suficiente fondo de armario como para poder hacer frente a las reivindicaciones sociales que exigía el cambio de época, lograron sortear (mejor o peor) las dificultades económicas heredadas. Por el contrario, aquellos países que no contaban con esta flexibilidad fueron sumidos en el caos, con el impacto económico que esto supone. 

Dos ejemplos polares son muy esclarecedores. 

Por un lado, Estados Unidos, un país aliado y acreedor (directa o indirectamente) de toda la deuda europea, enfrentó este conflicto social gracias al margen del que disponía y el rendimiento económico observado fue absolutamente excepcional. Aunque era un hueso duro de roer, fue uno de los países más reacios en adoptar las reformas laborales de 1919 y solo cedió ante un salario mínimo de dos dólares al día. No obstante, frente a la oposición gubernamental de intervenir, tan pronto aparecieran las reformas laborales de la OIT (1919) se declararon en huelga 350.000 trabajadores de la siderurgia, 120.000 trabajadores textiles se sumaron en Nueva Inglaterra y Nueva Jersey fueron a la huelga, y otros 30.000 trabajadores de la seda en Paterson (Nueva Jersey). \textbf{La policía se declaró en huelga en Boston}. Y, lo mismo hicieron en Nueva York los fabricantes de puros, los camiseros, los panaderos, los camioneros y los barberos [...] 5.000 trabajadores de la International Harvester y otros 5.000 trabajadores urbanos se echaron en las calles (ZINN, 2011: 280). Ante esta situación, los empresarios, que veían en toda manifestación popular una amenaza a su poder político, adoptaron rápidamente algunos beneficios para los trabajadores.\footnote{%
    En 1920 el Departamento de Justicia ordenó el arresto y deportación de individuos que pensaban en doctrinas comunistas y anarquistas. En Detroit, el corazón del Sindicato de Trabajadores Automotrices, este tipo de redada resultó en el arresto de 827 personas, de las cuales 234 fueron subsecuentemente deportadas (KEERAN, 1980: 32).
}

\imagenfit{Fig2.png}{Evolución del salario real y de la producción por trabajador en manufacturas en los Estados Unidos, 1899-1925}{\href{http://www.iisg.nl/hpw/data.php}{International Institute of Social History.} (27.10.2016).}

A consecuencia de este nuevo clima laboral, en 1919 los salarios superaron en 5\% el aumento de la productividad. Además, aunque se experimentó cierta inflación inmediata a consecuencia de la escasez provocada por la guerra, los precios se controlaron y descendieron progresivamente en el periodo siguiente. En 1925 los precios disminuyeron 30\% con respecto a 1920 (ELGAR, 2000). Entre 1922 y 1928 los Estados Unidos alcanzaron una completa estabilidad de precios: \textit{Para 1923 los precios al mayoreo habían recuperado solo una sexta parte de su disminución desde 1920-1921. Desde entonces hasta 1929 cayeron en promedio 1\% al año} (FRIEDMAN, 1963: 298).

\imagenfit{Fig3.png}{Índice de crecimiento del empleo, salario real, de producción física, del valor estimado de edificios, maquinaria y equipos y del gasto de equipo por trabajador, 1989-1925}{Paul H. Douglas, Real Wages in the United States, 1890-1925, Nueva York, 1966.}

Los mayores efectos del aumento de los salarios y de la reducción de la jornada laboral se dejaron sentir en el incremento de la producción y de la productividad. Acompañando este progresivo aumento salarial, se experimentó un aumento más que proporcional de la inversión en capital físico. Para 1925 la producción física había aumentado a más del doble. En 1919 y 1920 el desempleo fue de $0.4$ y $1\%$, respectivamente (pese al regreso del frente de batalla de millones de jóvenes desempleados). Entre 1921 y 1922 el desempleo alcanzó $15,3$ y $12,1\%$, respectivamente, pero en los siguientes dos años el desempleo cayó a $5$ y $7\%$, respectivamente (DOUGLAS, 1966: 427). En 1925 el empleo aumentó $30\%$ con respecto al inicio de la guerra en 1914. Observaciones similares de vieron en otros países como Canadá y, en menor medida, Inglaterra (menos preparada).

Sin embargo, estas fuertes demandas sociales también se experimentaron en países con muy, muy poco margen de maniobra. Aunque, el ejemplo más claro es Alemania, las causas de los desequilibrios estructurales que persiguen a los países que quedaron rezagados pueden encontrarse en alguno de los siguientes factores.

\paragraph{Trade-off: Conflicto social y precios}
En los primeros años, fueron varios los países que perdieron totalmente el control sobre la evolución de los precios. No obstante, los efectos que generó dependieron de la postura frente a esta. El Gobierno Británico, donde los precios se habían triplicado durante la guerra y permanecían más altos que en Estados Unidos (su principal competidor), consideró prioritario el mantenimiento de la estabilidad cambiaria, sirviendo a los intereses financieros en detrimento de la economía real. El razonamiento fue sencillo: si la mano de obra y la industria británicas acababan siendo poco competitivas, no habría más remedio que soportar un periodo de deflación de salarios y precios. En este contexto, las industrias orientadas a la exportación, como las del carbón, el acero, la construcción naval y las textiles, recibieron un duro golpe y el desempleo acabó por llegar al 20\%. Los conflictos laborales y las huelgas se dispararon. La situación conexa en los países perdedores se abordó de manera diametrialmente opuesta. Enfrentados a conflictos sociales extremas, Alemania, Austria y otras regiones, abordaron la combinación de endeudamiento previo, los gastos de reconstrucción, las pérdidas por ocupación y la imposición de las reparaciones mediante la monetización del déficit. Entre 1919 y 1923, los precios llegaron a multiplicarse por más de un billón (o sea, un millón de millones); el valor de la moneda cayó a cero. La hiperinflación fue, además de un fenómeno monetario, un fenómeno político: provocó una bancarrota encubierta de la deuda interna alemana (COMÍN, 2014) y agravó el conflicto social (recayendo sobre los sectores de la población que no podían protegerse en activos seguros). 

\paragraph{Deuda heredada: Los desequilibrios exteriores}
Por otro lado, los desequilibrios exteriores heredados actuaban también como factor de lastre común. Gran Bretaña había contraído una enorme deuda con Estados Unidos, que ahora tenía el control de un gran porcentaje de las reservas de oro del mundo. Errando en favor de la economía nominal, la economía Británica perdía competitividad mientras el Gobierno británico mantenía tipos que evitaran la fuga de capitales; elevando los costes de financiación y el repago de la deuda. En los países perdedores emergian los mismos problemas, las reparaciones de Guerra exigidad por los países conexos obligaban al Gobierno a recurrir a la monetización, por ser la imposición una vía políticamente inaceptable, agravando a su vez la espiral inflacionista que venían experimentando.

A estos efectos comunes se deben también añadir los profundos cambios en los flujos comerciales a nivel internacional, de donde emergieron nuevos miniestrellas en el continente americano (Argentina, Venezuela, Chile y Uruguay) y la menor movilidad de factores productivos (principalmente de trabajo) que agravaban el desempleo estructural. Además, en los países perdedores se produjeron \textbf{problemas de subsistencia}, por la destrucción de su sistema productivo y la imposibilidad de comerciar por su deteriorada relación de intercambio.

\subsubsection{Segunda etapa: Buenas noticias}
El desarrollo económico, como vemos, fue absolutamente desigual en las diferentes potencias. La tensión acumulada llego a su cumbre en 1923 cuando, debido al impago sostenido de las reparaciones de guerra, Francia ocupó la región alemana del Ruhr (1923).

Incapaz de hacer frente a la situación, la comunidad internacional, impulsada por Estados Unidos, abrió pasó a una serie de concesiones que condicionarían positivamente la recuperación económica en el viejo continente (especialmente alemania): el Plan DAWES (1923).

En diciembre de 1923, británicos, franceses y alemanes llegaron al acuerdo de nombrar una comisión que sería presidida por CHARLES G. DAWES, el primer director de la Oficina del Presupuesto de Estados Unidos. Dicha comisión debía considerar las formas de equilibrar el presupuesto alemán, estabilizando la moneda y fijando un nivel nuevo y viable de reparaciones anuales.\footnote{%
    El plan estipulaba que entre 1924 y 1929 la deuda sería abonada de una forma asumible para la economía alemana, fijándose de esta forma un pago inicial de mil millones de marcos-oro anuales, de los cuales ochocientos vendrían a ser cubiertos por un préstamo exterior y los pagos se elevarían hasta dos mil quinientos millones de marcos en 1929. Los pagos se complementarían posteriormente con una cifra que se fijaría de acuerdo con la situación de la economía alemana. El Reichsbank sería reorganizado bajo la supervisión de los vencedores, y para evitar la inflación, se le prohibía descontar los bonos del Tesoro y se le obligaba igualmente a mantener unas reservas de oro y moneda extranjera equivalentes, por lo menos, al cuarenta por ciento del papel moneda emitido.

Asimismo se estableció que franceses y belgas se retirarían del Ruhr y que sólo podrían adoptarse medidas como la ocupación del Ruhr, en caso de flagrante incumplimiento por parte de Alemania y únicamente con el acuerdo unánime de todos sus acreedores. El gobierno alemán aceptó el plan el 17 de abril de 1924 y el 24 lo hicieron los gobiernos inglés, italiano y belga. El primer ministro francés Raymond Poincaré firmó el día 25, ante el temor de verse aislado y presionado por el desmoronamiento del franco.
} Para sustentar las reparaciones de guerra, se abandonó la iniciativa de requisar materias primas y bienes muebles fijadas en el Tratado de Versalles (pues ello impediría a la larga el sostenimiento de la economía alemana) y, por el contrario, las reparaciones se financiarían con impuestos al consumo, tributos aduaneros y derivados de la explotación de los ferrocarriles. 

El nuevo plan pudo considerarse, a todos los efectos, un éxito. Por un lado, alivió la carga que implicaban las reparaciones de guerra para la economía alemana, estimuló las inversiones extranjeras en este país y ayudó a que los industriales germanos contasen con financiamiento para volver a participar en los mercados mundiales con productos manufacturados. A partir de 1925, Europa experimentó una fase de estabilización monetaria y recuperación económica, que coincidió con la reconstrucción de circuitos comerciales y un boom industrial tanto en Europa como en Estados Unidos.

Esta época se conoce como años locos, años dorados o felices años veinte; años donde se dejó atrás el pesimismo de la inmediata posguerra y se expandió el consumo, el confort y el entretenimiento. Algunas costumbres y manifestaciones culturales de los felices años veinte, que empezaron en EEUU en 1920, se repitieron en algunas ciudades de Europa (como Berlín o París) durante la segunda mitad de la década de 1920. Incluso en Alemania, pasados los peores años de la posguerra, la República de Weimar se caracterizó por una gran creatividad intelectual y estética.

\paragraph{Talón de Aquiles: EE.UU}
No obstante, el Plan DAWES dependía de un aspecto crucial: se sostenía casi exclusivamente en la marcha de la economía de EE. UU. Como, en realidad, eran plenamente conscientes de que la economía alemana no podría hacer frente a los pagos a partir de su rendimiento económico interno, se le permitió endeudarse a un tipo preferente ofrecido por los estadounidenses. De esta forma, se generó un círculo (vicioso) crucial: (i) bancos estadounidenses otorgaban fondos para que su Gobierno los prestase a Alemania; (ii) el Gobierno alemán usaba gran parte de esos fondos para pagar las indemnizaciones de guerra a otros países europeos (como Gran Bretaña, Francia y Bélgica); y, (iii) estos a su vez destinaban las reparaciones abonadas por Alemania para pagar a EE. UU. su propia deuda de guerra.

\textbf{INSERTAR DIAGRAMA DE FLUJOS}

Y, ¿qué pasó? El suceso es de sobra conocido: este círculo virtuoso devino vicioso tras el estallido bursátil del 29. La crisis de liquidez que siguió al Crack del 29 motivó que los bancos estadounidenses no pudiesen seguir prestando dinero al Gobierno, lo que influyó también en las reparaciones abonadas por los alemanes a Gran Bretaña y Francia y volvía de nuevo impactando en EE. UU., que tampoco recibía pago alguno por parte de sus (antiguos) aliados europeos.

Por tanto, si antes de 1929, los años de prosperidad de EE. UU. influyeron en el desarrollo del plan, pues los estadounidenses no presionaban al Reino Unido ni a Francia (ni a Alemania) para devolver las deudas de guerra, y esto repercutía en las reparaciones que Alemania debía pagar. A partir del crack, tanto la fragilidad del sistema económico internacional se hizó más que evidente y el conflicto socio-económico subyacente (en el que venimos incidiendo) no pudo ser contenido más tiempo. 

\subsubsection{Tercera etapa: La larga agonía}
Llegamos así a la tercera etapa del periodo de entreguerras, cuyos protagonistas principales serán: (i) la Gran Depresión y el colapso de la economía internacional, en el plano económico; y, (ii) la adopción generalizada de nuevas y radicalmente diferentes alternativas en el plano político.

Si bien es cierto que esta marea de cambios llegó de manera desigual a las diferentes naciones, el patrón es relativamente estable.

\paragraph{Gran Depresión: Perspectiva económica}
Desde una perspectiva económica, los efectos del crack se extendieron rápidamente al resto del mundo, tanto por el canal financiero mencionado (deuda) como los canales que emergieron de las Políticas emprendidas por los respectivos Gobiernos, en una reacción al más puro estilo de \textbf{Guerra Total} economica -lo que se desconoce esta vez es el enemigo: ¿extranjeros o nacionales?-. 

En primer lugar, conocemos el canal financiero de crédito a través del que se extendió la crisis a todos los países. El clima de desconfianza se generalizó en Estados Unidos y condujo a una crisis de liquidez (\textit{pánico bancario}).\footnote{%
    Después del pánico de 1929, y durante los primeros 10 meses de 1930, 744 bancos estadounidenses quebraron. (En total, 9,000 bancos quebraron durante los años 30s). Para abril de 1933, alrededor de \$7 billones en depósitos habían sido congelados en bancos quebrados o aquellos a los que les fue retirada la licencia después de la Ley de Emergencia Bancaria. Las quiebras bancarias se multiplicaron a medida que banqueros desesperados retiraron créditos, para los cuales los prestatarios no tenían ni tiempo ni dinero para pagar. Con la expectativa de rendimientos futuros pobres, inversión de capital y la construcción se desaceleraron o cesaron completamente. Frente a los préstamos no productivos y el empeoramiento de futuros prospectos, los bancos sobrevivientes se hicieron aún más conservadores en sus préstamos.
} Mientras que, en la escena internacional, el alto grado de interrelación financiera sumado a la contracción monetaria de la FED, que tuvieron que acomodar el resto, sentenció la partida: (i) quiebra bancos europeos; (ii) congelación de créditos a (y en) Europa, y, (iii) el colapso de la demanda estadounidense, que redujo la producción en Europa -y viceversa-, obligando a compensar de nuevo por cuenta corriente.

En definitiva, la colapso bursátil se canalizó por múltiples vías a la economía real y, muy desde el principio, se contagió por la serie de interdependencias que se habían construido durante el lustro precedente. Esta fue, sin duda, una crisis de intensidad desconocida hasta esos momentos. En EE. UU. se alcanzaron hasta 12 millones de parados, 7 en Alemania. La producción industrial de los países desarrollados se contrajo en más de un 40\% y el comercio internacional en un 20\%. 

\paragraph{Corrientes de cambio: Perspectiva política}
Ante este escenario, todo lo que conocemos de la época de los años treinta no debería resultarnos extraño. Como indica JEFFRY FRIEDEN, el politólogo de Harvard: \textit{Quienes apoyaban el orden clásico habían argumentado que dar prioridad a las relaciones económicas internacionales requería quitar importancia a preocupaciones tales como las reformas sociales, la construcción de la nación y la defensa nacional}. Una vez se quedaron sin argumentos, se abrieron las compuertas. Los comunistas optaron por la reforma social frente a la economía global y se mantuvieron al margen de los mercados mundiales. Los fascistas optaron por dar prioridad a la nación, generando una ola de nacionalismo económico en Europa y en los países en vías de desarrollo.

Una ola de proteccionismo asoló los países.\footnote{%
    Particularmente perjudicial fue la adopción generalizada de límites cuantitativos (contingentes) sobre las importaciones, que se habían abolido en gran medida en favor de unos impuestos más transparentes sobre las importaciones.

Además se vio acompañado por la imposición de controles de cambios: monopolización de los mercados de divisas de cada país por parte del Estado, para limitar las salidas de capital, regular las importaciones y/o manipular el tipo de cambio con fines de política económica.
} Estados Unidos estuvo entre los peores transgresores. La infame Ley Smooth-Hawley (1930) tenía como fin resguardar tras altos muros protectores a todo sector económico que tuviera alguna influencia en el Congreso. Los países europeos, que tenían parecidas razones económicas, aprovecharon el paso dado por Estados Unidos, que sirvió tanto de excusa como de pistoletazo de salida. Países como Francia y Suiza, que siguieron mucho tiempo con el oro y disponían de menos libertad para estimular sus economías, fueron especialmente propensos a las barreras al comercio. Cuando Hitler llegó al poder en 1933, utilizó la Políticas Comercial de forma estratégica para aprovecharse al máximo de sus vecinos del sureste de Europa. Hasta Gran Bretaña se apuntó imponiendo aranceles del 10\% en una amplia gama de importaciones a través de la Import Duties Act de 1932.\footnote{%
    Unos años más tarde, JOAN ROBINSON etiquetó a las políticas económicas aplicadas luego de la I Guerra Mundial como políticas de empobrecer al vecino; el problema para los europeos fue que estas políticas acabaron empobreciendo a los propios países que las implementaron, desintegraron la economía internacional y minaron las relaciones entre naciones. Dos décadas después de terminado el conflicto bélico, el objetivo dejó de ser el empobrecimiento del vecino: pasó a ser su destrucción.
} De hecho, la tendencia proteccionista llegó también a zonas en  desarrollo como la India y Latinoamérica; la armada británica estaba demasiado debilitada y distraída con otros asuntos como para apoyar al libre comercio en la periferia. \textbf{Entre 1929 y 1937, el volumen del comercio mundial se redujo a la mitad}.


Uno podría ver esto como una clásula de escapa ante la adversidad económica. Pero las raíces más profundas del proteccionismo se encuentran en el cambio de papel del gobierno en la sociedad. Una sociedad activa y con más poder político exigía una mayor protección económica del gobierno ante la adversidad extrema. Los gobiernos todavía no proporcionaban las redes de seguridad amplias y la protección social actuales para mitigar la competencia internacional y paliar sus consecuencias negativas sobre los trabajadores. Los sistemas de creencias y los hábitos de cooperación internacional, que sirvieron bien al mundo bajo situaciones económicas razonablemente sanas, se hundieron bajo el impacto conjunto de unas circunstancias económicas distintas y el aumento del número de personas ante las cuales los gobiernos eran responsables y debían rendir cuentas.


\subsection{Implicaciones: ¿Qué debemos aprender?}
Recapitulemos todo lo que hemos presentado para ver si podemos aprender algo. ¿Qué era diferente esta vez? Pensándolo bien, nada era muy diferente a lo ya experimentado. Al fin y al cabo, los países se habían visto en situaciones más o menos iguales en situaciones no muy lejanas de la historia (aunque, estríctamente hablando, el carácter tan internacional sí era un factor novedoso). Primero la  economía, luego la política, y después la economía otra vez.  

Empecemos con los primeros factores económicos. El modelo de libro de texto de un ajuste externo (bajo TC fijo como el patrón oro) supone unos mercados laborales descentralizados e individualizados, con salarios flexibles: si un sector de la economía nacional se hacía poco competitivo a escala global, sus salarios y demás costes caerían, ayudándolo a retener cuota de mercado. La mano de obra más barata también reduciría el desempleo. Desde luego, no era así como funcionaba la economía de verdad, pero con el tiempo esta historia pasó a ser todavía más fantasiosa en la década de los 20. Como hemos visto, los conflictos sociales fueron un cambio estructural del periodo de entreguerras. Concesiones laborales, salario mínimo, jornada de trabajo, etc.\footnote{%
    Hubo un aumento significativo de afiliaciones sindicales en las dos décadas anteriores a los años 1920, y el malestar obrero fue en aumento hasta culminar en la huelga general de 1926.
} La capacidad de los trabajadores para mantener sus salarios significaba que una contracción monetaria daría como resultado un mayor desempleo.

Ahora la política hacía su aparición. Los amos políticos, en especial de los Bancos Centrales, comprendieron que ya no podían permanecer al margen de las consecuencias políticas de la recesión económica y el alto desempleo. Los trabajadores no solo estaban agrupados en sindicatos, ahora también tenían derecho a voto:\footnote{%
    El sufragio se había cuadruplicado en Gran Bretaña en la década posterior a la Primera Guerra Mundial. La prensa y la radio iban en camino de convertirse en medios de masas; los diarios nacionales británicos tenían en conjunto una circulación media de 10 millones de ejemplares en la década de 1930.
} la Política Económica empezaba a democratizarse y había un incipiente movimiento socialista con el que lidiar. 

La economía da el golpe de gracia. Una vez que los mercados financieros empiezan a cuestionar la credibilidad del compromiso de los gobiernos con una paridad fija respecto al oro, surge un factor de inestabilidad. Los gobiernos son ahora una presa fácil para los ataques especulativos. Con la más mínima sospecha de que las cosas van a salir mal, los inversores venden la divisa nacional, compran moneda extranjera y sacan el capital del país. Si la paridad se mantiene, les basta con dar marcha atrás a esas transacciones y no pierden nada. Pero si la moneda se devalúa, ganarán mucho dinero volviendo a comprar la moneda nacional a un precio mucho menor y repatriando su capital. Este síndrome resulta familiar en presencia de tipos de cambio fijos: para los mercados financieros es como lo de si sale cara, gano yo; si sale cruz, pierdes tú. En el proceso de la venta de una moneda nacional, los especuladores, naturalmente, ejercen presión a la baja sobre el valor de la moneda y acaban acelerando el colapso de la paridad. Es fácil que sus expectativas lleven en sí mismas la semilla de su propio cumplimiento. El caso de Gran Bretaña en el periodo de entreguerras demostró que unas reglas rígidas monetarias y financieras, como las del patrón oro, no combinan bien con una economía moderna y una constitución política  moderna. 

La narrativa del patrón oro como sistema financiero global autorregulado y que funciona como la seda dejó de resultar convincente a la vista de la nueva realidad política creada por la democracia. Es una lección que hubo que volver a aprender en la década de 1990.



\clearpage
\section{Bretton-Woods: 1946-1973}
\puntosclave{
    Si bien la II Guerra Mundial tuvo un impacto catastrófico en términos económicos para la mayor parte de los países involucrados, tras el final de la misma comenzó una época de recuperación marcada, en contraposición con el anterior periodo de entreguerras, por una intensificación de la cooperación internacional. El nacimiento del sistema de Bretton Woods supone una clara muestra del espíritu de cooperación que existió en esta etapa. 
    
    En el plano monetario, el sistema, a pesar de poseer ciertos problemas inherentes a su diseño, funcionó con un éxito razonable favorecido por la coyuntura y ciertas medidas tomadas por los principales países del sistema. No obstante, los cambios ocurridos a principios de los años 70 terminaron por precipitar la ruptura del mismo y la instauración de lo que se ha venido a conocer como el “no sistema”. 
    
    En el ámbito comercial, si bien no se logró la creación de una Organización Internacional de Comercio como tal, sí existió un importante impulso del multilateralismo a través de acuerdos como el GATT. 
    
    En términos generales, se trató de una época de elevado crecimiento sostenido y relativa estabilidad en comparación con los shocks que experimentaría la economía mundial a partir de las décadas de 1970 y 1980.
} 

La IIGM tuvo efectos devastadores sobre el aparato productivo de los participantes, mucho mayores de los que la I Guerra Mundial provocó. Si bien es cierto que las pérdidas en capital humano fueron, de alguna manera, equivalente; las pérdidas e capital físico fueron significativamente mayores, con una cobertura geográfica mucho más amplia, una devastación sin precedentes; así como también el nivel de transformación del aparato económico-productivo de las naciones implicadas. A modo de ejemplo, puede observarse la transformación que experimentó la economía estadounidense:

\imagenfit{Fig5.png}{Cuadro Macroeconómico: Economía de Estados Unidos (1929-2009)}{\href{https://www.economicsandpeace.org/wp-content/uploads/2015/06/The-Economic-Consequences-of-War-on-US-Economy_0.pdf}{Economic Consequences of War on the U.S. Economy}. Institute for Economics and Peace (2015)}

\textbf{¿Qué queremos decir?} Sin embargo, la memoría colectiva aún albergaba las consecuencias de la I Guerra Mundial y se la culpaba, en gran medida, por la sucesión de graves perturbaciones económicas que habían experimentado las economías desarrolladas y, que en última instancia, habían conducido inevitablemente a la II Guerra Mundial: \textbf{para evitar desastres económicos y políticos cualquier futuro orden económico internacional tendría que encontrar un mejor equilibrio entre las demandas de la economía internacional y las de los grupos sociales en los distintos países}. El diseño de este compromiso requeriría entender mejor la relación entre el libre comercio y las tensiones sociales.\footnote{%
    Esto es, que la política comercial es conflictiva en términos políticos porque tiene consecuencias importantes sobre la distribución interior de la renta y porque contrapone valores e instituciones que pueden ser muy distintos en los países implicados. Nada de esto importaría mucho si las políticas comerciales pudieran aislarse de las políticas nacionales y fueran solo un asunto que concerniera a los expertos —⁠la fantasía del economista del libre comercio⁠—. Nunca ha sido ese el mundo donde vivimos; ni es probable que nos encontremos en un mundo así en un futuro cercano. Bajo el mercantilismo, la política comercial y la política nacional eran una misma cosa, como hemos visto. Incluso en el apogeo del liberalismo económico durante el siglo XIX, el aislamiento de las políticas comerciales del resto de las políticas de un país ocurrió brevemente y el proteccionismo comercial hizo una rápida reaparición en cuanto cayeron los precios agrícolas. La politización de las políticas comerciales tuvo un momento álgido en el periodo de entreguerras. La incapacidad de los gobiernos para responder a las quejas de sus empresas, trabajadores y campesinos en el contexto de una economía abierta contribuyó a la Gran Depresión.
} Con esto en mente, el nuevo clima de cooperación internacional culminaría con la firma del Acuerdo de Bretton Woods.

\subsubsection{¿Cuál es el contexto geopolítico?: Multilateralismo}



\subsection{Bretton-Woods: Las lecciones aprendidas}
¿Cómo podría restaurarse una economía global abierta en un mundo en el que predominaban las políticas nacionales?\footnote{%
    KEYNES ya era reconocido como el economista más destacado de su generación y como un agudo comentarista sobre la política contemporánea y los políticos. WHITE era un funcionario muy respetado del Tesoro de Estados Unidos del que más tarde se descubrió que había pasado información secreta a los soviéticos antes y durante la Segunda Guerra Mundial.
} El sistema económico internacional que diseñaron los economistas KEYNES y WHITE, y otros representantes de 44 naciones, fue una asombrosa obra de ingeniería institucional:\footnote{%
    Se le conoce como \textbf{Sistema de Bretton Woods} por el complejo hotelero de New Hampshire (EE.UU) donde tuvieron lugar las reuniones.
} demostraron firmemente que estaban dispuestos a evitar los errores del pasado. Veamos rápidamente cómo y qué se acordó.

\subsubsection{¿Qué es Bretton Woods? Dos planes, un hegemón}
Prescribir una apertura económica y esperar después que las políticas nacionales se adaptaran de un modo u otro ya no iba a funcionar, si es que había funcionado alguna vez.\footnote{%
    Como reflejo de la supremacía económica estadounidense, WHITE, en particular, se inclinaba por liberar la economía mundial de las grandes restricciones y controles impuestos durante el periodo de entreguerras y endurecidos todavía más durante la guerra. Por su parte, KEYNES había escrito un artículo importante en el apogeo de la Gran Depresión, en 1933, en el que describía su cambio de actitud respecto del libre comercio y su recién hallada preferencia por un cierto grado de autosuficiencia nacional. Al igual que la mayoría de los ingleses, KEYNES escribió que sentía una adhesión casi moral a la doctrina del libre comercio. \textit{Consideraba las deserciones habituales del }[libre comercio] \textit{como una estupidez y una atrocidad al mismo tiempo}. Pero cuando echó un vistazo a sus argumentos en defensa del libre comercio en la década de 1920, dejaron de parecerle sólidos. El compromiso total con el libre comercio era posible solo cuando las distintas sociedades estaban gobernadas por tecnocracias imbuidas de la fe en un tipo similar de capitalismo. Ese tipo de compromiso dejó de ser posible, e incluso deseable, en un mundo en el que las diferentes naciones experimentaban con políticas económicas distintas.

En esencia, estos dos hombres notables eran también realistas y comprendieron que las reglas del comercio internacional y el dinero internacional necesitaban cambiar.
} La experiencia histórica era clara al respecto, cuando las necesidades nacionales chocaban con los requerimientos de la economía global, las necesidades nacionales acaban saliendo victoriosas. Por tanto el nuevo sistema debía basarse en un delicado compromiso: \textbf{permitir suficiente progreso y disciplina internacional en aras de la liberalización comercial para garantizar un comercio mundial con vitalidad, pero permitiendo que los distintos países respondieran con holgura a sus necesidades sociales y económicas}. O, dicho de otra forma, las Políticas Económicas internacionales debían someterse a los objetivos de las políticas nacionales y no a la inversa.

Bajo esta prermisa, KEYNES (Reino Unido) y WHITE (EE.UU), protagonistas/antagonistas máximos de la negociación, fueron envíados por sus respectivos Gobiernos para formular sus propuestas y alcanzar un Acuerdo. No hay que alabar su altruismo cosmopolita de estos economistas, ni KEYNES ni WHITE estaban únicamente motivados por esto, y así quedaba claro en los planes presentados.\footnote{%
    KEYNES estaba al tanto del declive económico británico y de su dependencia de Estados Unidos e hizo todo lo posible para favorecer los intereses británicos dentro de esos límites. WHITE promovió la causa del comercio y las inversiones estadounidenses y trabajó para fortalecer el poder de Estados Unidos en las nuevas organizaciones internacionales.
} Brevemente, la siguiente tabla resume los principales puntos:

\begin{center}
\begin{tabular}{|p{0.3\linewidth}|p{0.3\linewidth}|p{0.3\linewidth}|p{0.3\linewidth}|}
\hline
 & \textbf{Plan KEYNES} & \textbf{Plan WHITE} \\
\hline
\textbf{Tipo de Cambio} (SMI) & {\scriptsize El Tipo de Cambio sería fijo, con una amplitud de fluctuación puntual permitida de 10 puntos básicos ($\pm5\%$). La paridad fijada sería revisada si fuera necesario.} & {\scriptsize El Tipo de Cambio sería fijo salvo desequilibrios fundamentales que debieran ser corregidos. Esta revisión solo podría adoptarse con 4/5 partes de los votos.} \\
\hline
\textbf{Instituciones} & {\scriptsize Se constituye una International Clearing Union (ICU), de carácter técnico, entre Bancos Centrales. Cualquier saldo de un Banco Central se puede usar para negociar con otro bajo el marco de la ICU.} & {\scriptsize Se constituye un Fondo Internacional de Estabilización, de carácter político (multigubernamental) y con un sistema de votos ponderados.} \\
\hline
\textbf{Cuotas} & {\scriptsize Se aporta un capital inicial (de la ICU) de 30 mil millones de USD, a conseguir mediante un sistema de cuotas según la Balanza Comercial del país (75\% del comercio total: $X+M$ entre 1935-1939 para empezar). La divisa de pagos internacional sería una moneda fiduciaria de nueva circulación, el Bancor.}\footnote{%
    Es decir, el ICU emitiría una moneda internacional (Bancor) que tuviese respaldo en las monedas más fuertes y por medio de la cual se lograría la transferencia de excedentes de los países superavitarios a los países deficitarios. De esta forma, se vitaría lidiar con los mercados cambiarios y los movimientos internacionales de oro y monedas, logrando un crecimiento del comercio mundial.
}  & {\scriptsize Se aporta un capital inicial de 5 mil millones de USD, a distrbuir en función del nivel de renta, reservas y Balanza de Pagos. El capital del Fondo no sería de nueva circulación, sino que estaría basado en el principio de depósito.}  \\
\hline
\textbf{Liquidez} & {\scriptsize La unidad de cuenta del sistema y la liquidez sería una divisa de nueva creación: el Bancor. Esta divisa sería únicamente utilizada en el SMI, estableciéndose su valor a paridad con el oro pero ajustable. Además, los déficits o excesos de liquidez se cubren con un sistema de descubierto automático. Los países pueden servirse de este mecanismo de ajuste (o liquidez) automática hasta el límite de la cuota que le corresponde.} & {\scriptsize La unidad de cuenta del sistema sería la Unitas (1Un = 10\$), pero sería únicamente eso, la unidad de cuenta (por poner alguna). La liquidez del Sistema venía únicamente determinada por la contribución de los miembros y el capital total del Fondo, no podían haber descubiertos. El Fondo, además, no podría tener más del 200\% de la cuota de un país en su moneda, debía haber una distribución adecuada de las monedas en posesión del Fondo (preferiblemente, liquidez de las monedas fuertes).} \\
\hline
\textbf{Ajuste} & {\scriptsize El ajuste, tanto bilateral como multilateral, sería simétrico: deudores y acreedores. El ajuste no está sujeto a condicionalidad. Sin embargo, el sobreendeudamiento de un país podía llevar a la ICU a declarar en bancarrota a un país y cancelar a mutu propia la deuda con su(s) acreedor(es)} & {\scriptsize El ajuste se lleva únicamente a cabo con los países deudores. Además, estaría sujeto a condicionalidad.}\footnote{%
    En la negociación, Estados Unidos admite la posibilidad de que un país con superávits constantes genere dificultades, provocando una crisis de liquidez de su moneda. Ante esta circunstancia, permite la posibilidad de llevar a cabo acciones discriminatorias frente a esta.
} \\
\hline
\end{tabular}
\end{center}


Si analizamos atentamente los puntos esenciales de cada Plan, nos daremos cuenta de que estamos ante una concepción \textbf{muy diferente} de cómo debía organizarse el nuevo Sistema Económico Internacional: KEYNES, abanderaba una verdadera revolución de las relaciones internacionales económicas, centrándose en la economía real y la cooperación internacional en el ámbito económico, entre iguales y de manera colaborativa; WHITE, por el contrario, abanderaba la causa continuista (¿quizás por estrechez de miras? recordemos que era funcionario...), con suficientes ajustes para \textit{garantizar} estabilidad y un peso significativo sobre la economía nominal. Ahora bien, aunque sobre el papel había dos planes, sólo había un hegemón. 

\subsubsection{Acuerdo de Bretton-Woods: Plan WHITE}
Llegado a un punto, WHITE, de forma unilateral, decidió los porcentajes del voto de las principales potencias (Estados Unidos, Gran Bretaña, Unión Soviética y China) y encargó a un economista de su equipo que calculara durante la noche la fórmula que diera lugar a esos porcentajes e improvisara su justificación económica. Por tanto, entender el Sistema de Bretton-Woods es entender el Plan WHITE, y en ello nos centraremos. ¿Qué y cómo se implementó? Igual que en lo político, la aportación estadounidense más destacada al sistema económico internacional de la posguerra fue ese mismo multilateralismo.\footnote{%
    Dicho multilateralismo significaba que las reglas del comercio y las políticas económicas se canalizarían a partir de entonces a través de instituciones internacionales —⁠el Fondo Monetario Internacional, el Banco Mundial y el GATT (Acuerdo General sobre Aranceles Aduaneros y Comercio)⁠— en vez de ser impuestas por los países más poderosos. Se trataba de una innovación muy importante. Aunque la influencia de Estados Unidos era innegable, el multilateralismo dotaba a estas instituciones de un cierto grado de legitimidad independiente del poder estadounidense que las sustentaba. Nunca llegaron a ser verdaderamente autónomas de Estados Unidos u otras principales potencias económicas, pero tampoco fueron puramente una extensión de estas potencias. Desempeñaron un papel importante en el establecimiento de reglas, en hacerlas cumplir y en darles legitimidad. El multilateralismo dio voz a las naciones más pequeñas y más pobres y protegió sus intereses de un modo sin precedentes. En consecuencia, los estadounidenses, a diferencia de los británicos antes que ellos, acabaron creando una infraestructura institucional para el gobierno de la economía internacional que duraría más que su hegemonía indiscutida.
}

\paragraph{Plano Monetario: Fondo Monetario Internacional (1946)}
En el plano monetario se crearía un Fondo Monetario Internacional (1946) cuyo objetivo principal sería garantizar la estabilidad cambiaria. Para ello se adoptaron las siguientes normas:
\begin{lnum}
    \item \textbf{Patrón Cambios-Oro: Sistema fiduciario}. Se adoptó un patrón cambios-oro y, en consecuencia, un SMI fiduciario construido sobre el dólar, cuyo valor se fijo en 35\$ la onza de oro. Con respecto a esta moneda de reserva internacional, dólar, se establecería el resto de paridades. Las paridades serían fijas con fluctuación permitida de $\pm1\%$ y sólo se permiten ajustes ante desequilibrios fundamentales, pudiendo realizar: (i) un ajuste automático siempre que la variación fuera menor del 10\%; y, (ii) con autorización si superaba ese nivel;
    \item \textbf{Cuotas: Principio de depósito}. El capital del Fondo se constituiría mediante un sistema de cuotas: a cada país se le asignó una cuota de participación que debía cumplimentar en un 25\% con oro (o moneda convertible en oro) y en un 75\% con moneda nacional;
    \item \textbf{Convertibilidad: Economía nominal}. Los países debían garantizar la convertibilidad corriente de sus monedas, se prohibían las restricciones y discriminaciones en los movimientos de capital por operaciones corrientes. No obstante, el artículo XI permitía posponer \textit{sine die} dicha liberalización y fue la vía seguida por muchos; 
    \item \textbf{Ajustes: Préstamos por desequilibrios}. Se facilitaba financiación ante desequilibrios temporales que estuviesen generando una escasez de reservas. Estas facilidades se obtendrían mediante préstamos a corto plazo y sin condicionalidad, siempre y cuando fueran temporales y la cuantía máxima respetara el límite determinado atendiendo a su cuota (un tanto por ciento sobre la cuota: p.e. $250\%$). Si el desequilibrio fuera permanente, también se podría disponer de financiación pero sujeto a condiciones;
    \item \textbf{Supervisión}. El FMI supervisaría el funcionamiento del sistema, pudiendo regular algunos aspectos, y debía servir como órgano de consulta. Esto, por tanto, se sumaba a la función financiera.
\end{lnum}

Estas, por tanto, son las bases del sistema monetario internacional construidas bajo el paraguas de Bretton-Woods en 1944, e institucionalizadas en 1946 cuando el Fondo Monetario Internacional devino operativo. Sin embargo, el sistema monetario internacional estuvo operativa realmente (en su máxima expresión) a partir de 1958, cuando los últimas restricciones impuestas sobre la convertibilidad fueron derogadas por los países miembros.

\paragraph{Plano Financiero: Banco Mundial}
En el plano financiero se constituyó el Banco Mundial (más precisamente, el Banco Internacional de Reconstrucción y Fomento), la génesis del ahora Grupo del Banco Mundial. El objetivo establecido para esta institución era claro: financiar el desarrollo de los países facilitando recursos financieros cuando fueran necesarios; y, para este fin actuaba, precisamente, como eso: un verdadero Banco al estilo anglosajón. 

Para dotar de capital al Banco Mundial se establecieron dos formas: (i) capital suscrito por los Estados miembros; y, (ii) capital exigible, que permitía al Banco Mundial emitir obligaciones (al estilo de bonos corporativos) en los mercados financieros para financiar sus operaciones.\footnote{%
    Cada país miembro, por tanto, se comprometía a suscribir/aportar capital según una cuota ligada a su peso económico, que se descomponía en: (i) capital desembolsado, pequeña fracción efectivamente ingresada; y, (ii) capital exigible, un compromiso contingente que sólo se activaría si el Banco no pudiera cumplir sus obligaciones frenta a sus acreedores.

La innovación fue precisamente el capital exigible, ya que permitía al Banco emitir bonos con alta credibilidad sin necesitar enormes aportaciones de capital inicial. Además, constituye esto también la principal diferencia con el FMI: mientras el Grupo Banco Mundial se puede apalancar (endeudarse en los mercados), el FMI no.
} Iniciaría con esto sus operaciones financieras, siendo sus primeros préstamos para Francia e Irán (el primero, sustancialmente más elevado: 250M€).\footnote{%
    Los préstamos del Banco Mundial supusieron un verdadero alivio para la reconstrucción de Europa, pero no surtieron el efecto catalizador a la velocidad que se esperaba. Es por ello que surgió el \textit{European Recovery Plan}, comúnmente conocido como Plan Marshall (aprobado en 1948). Este \textit{relief program} fue canalizado a través de la Organización Económica de Cooperación Europea (OECE antecesora de la OCDE), que repartió finalmente 12.400 millones de dólares en toda Europa, salvo en España. 

En la fase previa a su aprobación, 1947, se planteó incluir desde un primer momento a la URSS, aceptando Stalin en la creencia de que demostraba la debilidad capitalista. Sin embargo, la inquietud en el bando capitalistas, que creyó que esto conduciría directamente al rechazo por parte del Congreso de los EE.UU., hizo que se presentara al ministro de Stalin, MOLOTOV, unas condiciones muy difíciles de aceptar. En consecuencia, la URSS rechazó cualquier opción de cooperación y creó el COMECON (CAME en español, Consejo de Ayuda Mutua Económica), un área de libre comercio planificado que dividía el territorio en áreas de producción según la especialización regional y que llegó a dominar el 10\% del comercio mundial.
}

El Banco Mundial es, a fecha de hoy, un grupo de organizaciones que busca erradicar la pobreza y promover el desarrollo económico. 

\paragraph{Plano comercial: Organización Internacional del Comercio (GATT, 1947)}
Por último, en el plano comercial se proyectó incialmente la constitución de una Organización Internacional de Comercio.\footnote{%
    La OIC, en la propuesta inicial, abarcaba acuerdos sobre estabilización de precios de productos básicos, medidas antitrust internacionales y normas laborales justas, pero se fue a pique debido a la política estadounidense. Al Congreso de Estados Unidos le preocupaba su invasión de las prerrogativas nacionales.
} Sin embargo, al no llegar  a constituirse, la encarnación institucional del multilateralismo en el comercio fue el Acuerdo General sobre Comercio y Aranceles (\textit{General Agreement on Trade and Tariffs} de 1947, GATT). Y así permanecería durante cincuenta años.\footnote{%
    Si bien el GATT no se constituyó formalmente como una organización plena como el FMI o el Banco Mundial, se gestionaba en un pequeño secretariado en Ginebra. Esto le permitió convertirse, de hecho, en el foro multilateral que supervisaba la liberalización del comercio global. ¡Y qué gran éxito!
}

\begin{specialblock}{La OIC - ¿Qué era? ¿Por qué era importante?}
    Para los participantes de Bretton Woods, el fracaso de la OIC fue un duro golpe. El objetivo de la Organización Internacional del Comercio era establecer un sistema de comercio que incrementase el bienestar pero limitase o corrigiese los déficit o superávit comerciales crónicos y pudiesen poner en brete los otros pilares del Sistema.
    
    Más aún, si se observa desde la óptica de KEYNES, donde el sistema de pagos internacionales hubiese sido centralizado y mediado en la ICU (International Clearing Union), hubiese servido de cierre perfecto del círculo, permitiendo a los países utilizar un foro multilateral para mediar sus problemas comerciales mientras los ajustes realizados en otro les daban \textit{respiro} para hacerlo de manera progresivo. Su acta fundacional, la Carta de la Habana, fue adoptada pero nunca llegó a constituirse debido a la reticencia del Congreso de los Estados Unidos, que ostenta la potestad última en materia de Política Comercial. 
    
    Todos los aliados salvo los soviéticos participaron en sus negociaciones. 
\end{specialblock}

A pesar de unos comienzos lentos, las sucesivas \textit{rondas} de negociaciones comerciales multilaterales (ocho en total entre 1947 y 1995) lograron eliminar una parte sustancial de las restricciones sobre las importaciones existentes desde la década de 1930. Además, se introdujo el principio de nación más favorecida (NMF), por el cual se garantizaba que todos los firmantes del GATT se beneficiarían de cualquier bajada de aranceles, al margen de lo activa que fuera su participación en las negociaciones. 

No obstante, es probable que el multilateralismo comercial de la era Bretton Woods se entienda mejor por sus límites que por sus logros. Como es normal, las políticas comerciales nacionales siguieron en conflicto. Sin embargo, esto tuvo muy poca presencia en la política dese una perspectiva nacional. ¿Por qué? La razón última son sus límites y numerosas excepciones. En primer lugar, la liberalización apenas tocó sectores críticos:
\begin{lnum}
    \item \textbf{Agricultura}. La agricultura se mantuvo fuera de las negociaciones del GATT y permaneció con aranceles y otras barreras. De hecho, para mayor infamia, en forma de contingentes de importación variables cuya finalidad era estabilizar los precios nacionales en niveles mucho más altos que en los países exportadores⁠; y,
    \item \textbf{Servicios}. La mayoría de los servicios (seguros, banca, construcción, suministros, etc.) se excluyeron también de cualquier liberalización;
    \item \textbf{Países en vías de desarrollo}. Además a los países en vías de desarrollo se les proporcionó bastante libertad para hacer lo que quisieran con sus Políticas Comerciales. De hecho, normalmente, ni se les pedía que ofrecieran concesiones arancelarias durante las negociaciones del GATT, incluso cuando se beneficiaban de las reducciones arancelarias de otros por el MFA. Habían recurrido a varias cláusulas del GATT que les permitían limitar sus importaciones prácticamente de forma permanente. 
\end{lnum}

Además, si la liberalización había llegado lo suficientemente lejos como para poder generar tensiones en la política local, ¡no pasaba nada! Había una larga lista de herramientas que permitían reducir la presión exterior del sector nacional que se quisiera proteger. De hecho, las reglas contenían resquicios tan amplios que habría podido pasar un elefante.\footnote{%
    Por ejemplo, cualquier empresa podía comprarse protección por medio de las cláusulas Antidumping (AD) o de Salvaguardia del GATT. El acuerdo AD, concretamente, era una abominación desde el punto de vista del libre comercio. El país que importaba podía imponer aranceles siempre que determinara que un exportador había vendido sus productos a un \textit{valor inferior al normal} y había causado \textit{perjuicio} a la industria local. Las autoridades nacionales podían manipular fácilmente el significado de \textit{valor inferior al normal}. Y los aranceles punitivos podían imponerse incluso si el comportamiento en cuestión constituía una práctica comercial normal, tal como vender por debajo del coste en una fase bajista del ciclo económico, o cuando el infractor no tiene la capacidad de monopolizar el mercado local. Las empresas del país explotaban ampliamente estas reglas con el fin de obtener una protección a medida.
}

Y esta posibilidad fue explotada, en abundancia dirían algunos. Por ejemplo, los sectores industriales que habían sido liberalizados pero empezaban a enfrentarse a la amenaza de una competencia significativa de otros exportadores con bajos costes y alta productividad, pronto recibieron protección para evitar enfrentarse a su destino. Así, las industrias textil y de la confección de los países desarrollados estuvieron protegidas desde 1974 por el Acuerdo relativo al Comercio Internacional de los Textiles (MFA), un conjunto de contingentes negociados bilateralmente sobre las exportaciones de los países en vías de desarrollo. Más aún, la década de 1980 vio la difusión de las limitaciones voluntarias a la exportación (VER).\footnote{%
    Acuerdos por los que los exportadores japoneses (sobre todo) de automóviles, acero y otros productos industriales se comprometían a mantener sus exportaciones dentro de unas cuotas específicas.
} 

Por otro lado, y para más inri, la capacidad para hacer cumplir los acuerdos del GATT era un chiste. Si un gobierno pensaba que otro se había saltado las reglas, podía pedir a un panel del GATT que arbitrara. Si el panel decidía a favor del demandante y el informe del panel era aprobado por los miembros del GATT, la parte culpable tenía que cambiar la política infractora o de lo contrario el demandante tenía derecho a una compensación. El único inconveniente era que la aprobación del informe del panel requería una decisión unánime: cada uno de los miembros del GATT, incluyendo el país que había incumplido las reglas, tenía que firmar la decisión. Si un jurado incluye al acusado, es seguro que no se pronunciará unánimemente en su contra. Así que las reglas del GATT dejaban segmentos completos del comercio mundial al descubierto; eran débiles donde existían; y estaba claro que no se podía obligar su cumplimiento.\footnote{%
    Estas características hicieron que la institución tuviera deficiencias evidentes y que la nueva Organización Mundial del Comercio, que nació en 1995, resultara mucho más atractiva desde el ángulo del libre comercio. Pero criticar el GATT porque se quedaba considerablemente corto en cuanto al libre comercio sería juzgarlo desde una perspectiva inadecuada. Puede que el GATT no hubiera tenido como finalidad minimizar las complicaciones económicas entre las naciones, el objetivo para el cual el KEYNES de entreguerras había expresado cierta simpatía. Pero fue diseñado para dar espacio a cada nación participante para que persiguiera sus objetivos sociales y económicos de forma relativamente libre de restricciones externas, aunque dentro de un marco flexible de cooperación internacional.
}

En definitiva, podemos repetir el mantra más recurrente de la exposición: \textbf{cuando el comercio internacional amenazara los pactos o sensibilidades nacionales sobre la distribución de la renta, el comercio sería el que cedería}. 

\subsubsection{Implicaciones de Política Económica: ¿Qué podemos aprender?} 
Ahora que conocemos el célebre Sistema de Bretton Woods y sabemos cuáles fueron sus pilares funcionales, podemos preguntarnos ¿cuál fue el resultado? ¿cómo funcionó?

Como suele pasar, los datos describen mejor la situación: el volumen del comercio mundial creció a una tasa media anual de casi el 7\% entre 1948 y 1990, considerablemente más rápido que hasta la fecha; también lo hicieron los niveles de producción de los países ricos y pobres (en ambos casos a la vez, causa y efecto del rápido auge del comercio⁠); en términos de amplitud y profundidad del progreso económico, \textbf{el régimen de Bretton Woods eclipsó todos los periodos anteriores, incluido el del patrón oro y la época del libre comercio del siglo XIX}. 

Entre 1950 y 1970 los países industrializados experimentaron un crecimiento sostenido a tasas anuales promedio de más del 5\%.\footnote{%
    No obstante, el ritmo de crecimiento económico no fue uniforme entre este grupo de países. Durante las décadas de 1950 y 1960, Japón creció mucho más rápido que otros países, a tasas anuales promedio del 9,8\%, seguido por Alemania Occidental (6,3\%), Italia (5,6\%) y Finlandia y Francia (5\%). Las tasas de crecimiento más bajas fueron las experimentadas por EE. UU. (3,5\%) y Reino Unido.
} También la inversión creció de manera increíble, estimulando la actividad económica incrementando la demanda en numerosos sectores, aumentando la cantidad de capital por trabajador y permitiendo un mayor grado de innovación.\footnote{%
    ucho se ha especulado acerca de las causas de este elevado grado de inversión: elevados niveles de demanda doméstica y externa; una oferta sustancial de ahorro doméstico y extranjero; importantes mejoras tecnológicas; un rol progresivamente más enfocado al fomento del crecimiento por parte del sector público; amplia oferta de trabajadores a salarios bajos (constituida por desempleados y migrantes de zonas rurales a nivel doméstico e inmigrantes llegados del exterior) explica también por qué el output y las exportaciones crecieron más rápidamente en países como Alemania, Italia y Suiza que en otros como Reino Unido o Bélgica, donde el trabajo era relativamente más escaso.

Pero se suele perder de vista el punto esencial. La inversión no aumentó más ni menos que en otros periodos, por ejemplo el actual, sino que se dirigió a sectores reales y no a los mercados financieros, como nunca antes se había experimentado, dada la magnitud del crecimiento y la relativamente frágil internacionalización del capital.
}

\paragraph{Perspectiva política: embedded liberalism}
Desde una perspectiva política, JOHN RUGGIE es probablemente el que mejor definió este periodo, al que llamó \textbf{embedded liberalism}: \textit{A diferencia del nacionalismo económico de los años treinta, [El régimen] tenía un carácter multilateral; a diferencia del liberalismo del patrón oro y el libre comercio, su multilateralismo estaba basado en el intervencionismo nacional}. 

Este Sistema permitió a los países desarrollados construir a medida sus versiones del capitalismo en torno a criterios diferentes de institucionalización del mundo empresarial, mercados de trabajo, regímenes impositivos, relaciones gobierno-empresa y formas de organizar el Estado de bienestar: surgieron \textbf{variedades de capitalismo} (HALL y SOSKICE).\footnote{%
    Estados Unidos, Gran Bretaña, Francia, Alemania y Suecia eran economías basadas en el mercado, pero las instituciones que sostenían sus mercados diferían sustancialmente e incorporaban características inequívocamente nacionales. Solo en la Europa continental había al menos tres tipos diferentes de capitalismo: el modelo alemán de economía de mercado social; el Estado de bienestar escandinavo, y el sistema francés basado en la \textbf{planificación indicativa} y una sustantiva reglamentación de los mercados. Japón también emprendió su propio camino y creó un sector exportador hipercompetitivo junto a una economía tradicional regulada y protegida. Estados Unidos se mantuvo como el ejemplo de referencia de la economía de mercado liberal, aunque a su liberalismo económico le faltaba la ambición que adquirió en la década de 1980.
} 

La situación era muy parecida en el mundo en vías de desarrollo, donde los esfuerzos nacionales se dirigían a impulsar la industrialización y el crecimiento económico. A falta de disciplina exterior, las naciones en vías de desarrollo eran libres para desplegar una ingente gama de políticas industriales para transformar sus economías y reducir su dependencia de los recursos naturales y mercancías. Así, muchos de ellos lograron conseguir un alto crecimiento basado en la producción de manufacturados.\footnote{%
    También se citan a menudo los efectos \textit{spillover} generados por el fuerte crecimiento de los países desarrollados y el altruismo como promotores de este crecimiento. 

En el segundo caso, se sugiere que el incremento de la riqueza de los países desarrollados surgió un deseo por mejorar el bienestar económico de las regiones más pobres del planeta. Para racionalizar esto, como gusta a los economistas, se alude a: (i) el atractivo del comunismo como ideología política y modelo económico, que suponía una amenaza para la mayor parte de países occidentales; (ii) que los beneficios económicos de un aumento generalizado de los estándares de vida y las mayores oportunidades para el comercio derivadas del mismo compensarían con creces las posibles pérdidas experimentadas por ciertos países avanzados con posiciones privilegiadas (que podrían verse afectadas por los cambios asociados al desarrollo económico de los países en desarrollo); o, (iii) que la concesión de independencia a las antiguas colonias europeas en otros continentes vino aparejada del otorgamiento de un trato favorable a las nuevas naciones emergentes para suavizar el ajuste a su nueva situación. 

Todo esto es posible que jugase un papel importante, sin embargo, para el primer factor ya estaban las armas (como se demostró amplíamente en el mismo periodo: Guerra de Corea y Guerra de Vietnam) y los segundos, aunque importantes, proporcionan una visión parcial de la situación. Desde mi punto de vista, el crecimiento económico fue tal que es normal que emerjan motivaciones altruistas en interés de un bien común. Al fin y al cabo, cuando hay tanto para repartir, los principios de suma cero que rigen el intercambio son perfectamente reinterpretados por la máxima comunista: \textit{de cada cual según sus capacidades, a cada cual sus necesidades}.
} 


\paragraph{Comercio internacional}
Desde una perspectiva comercial, el comercio, como vemos, se hizo (y permaneció) libre solamente donde planteó escasos retos a las instituciones, preferencias de distribución o valores nacionales. Y, sin embargo, experimentó un crecimiento anual como ninguno hasta la fecha. 

De hecho, ¡el volumen de comercio creció, de media, a una tasa superior al de la producción mundial! El objetivo era un comercio más libre en algunas áreas, no el libre comercio en todas ellas. De hecho, aunque los gastos de transporte siguieron bajando, los responsables políticos mostraron una firme y decidida falta de ambición por liberalizar más el comercio. Partes enormes del comercio mundial permanecieron completamente fuera de cualquier acuerdo multilateral o protegidos por generosas excepciones a los acuerdos existentes. 

¿Qué nos permite ver esto? Debemos entender que lo que impulsó la globalización, lo que hizo explotar el crecimiento del comercio internacional, no fueron las políticas de liberalización. Sino el propio crecimiento económico que la equidad, seguridad y estabilidad de los acuerdos de Bretton Woods ayudaron a respaldar. El crecimiento de una amplia base facilitaba la globalización porque ayudaba a amortiguar los impactos distributivos del comercio internacional. Resulta menos notoria la agitación de las aguas cuando la marea de las oportunidades económicas sube repentinamente y levanta todas las barcas. Así, las \textbf{políticas nacionales favorecieron la globalización como subproducto de un crecimiento económico ampliamente compartido junto con alguna apertura modesta}. El éxito de la época de Bretton Woods parece indicar que unas economías nacionales sanas son la base de una economía mundial dinámica, incluso en presencia de restricciones al comercio.

\begin{specialblock}{Mierda de CECO que debe ser depurada}
    Sin embargo, no todas las regiones del planeta compartieron los mismos ritmos de crecimiento de su actividad comercial, dado que las tasas de crecimiento de los países industriales de Europa, América del Norte y Japón fueron superiores a la media mundial en los años 50, por lo que vieron incrementada su participación en el total de exportaciones del 61\% al 70\%. En lo que a composición del comercio internacional se refiere, cabe destacar hasta 1973 un aumento en la importancia de las manufacturas y una correspondiente disminución en la de materias primas y alimentos. Los cambios en la composición del comercio de manufacturas siguieron el patrón establecido a finales del siglo XIX, intensificándose los flujos de productos de ingeniería (especialmente de vehículos a motor) en detrimento de los textiles. 
    
    Desde 1963, se observa una importante caída en el comercio de metales, mientras que los productos químicos han continuado ganando importancia. En conjunto, en esta etapa continuó el cambio desde los bienes de consumo y semimanufacturados hacia bienes de capital que comenzó en el período de entreguerras, por lo que los países industrializados fueron los más beneficiados por el crecimiento del comercio internacional después de 1950. Este sesgo creciente del comercio de mercancías a favor de los bienes de capital tuvo como consecuencia un cambio en la dirección del comercio total en favor del comercio entre países industriales frente a las exportaciones de los países en desarrollo hacia las regiones más avanzadas. 
    
    Por tanto, dicha tendencia parecería contradecir las predicciones hechas por algunos economistas en los años 30 de que, a medida que las naciones se desarrollan y sus dotaciones de factores se vuelven relativamente más uniformes, la necesidad del comercio basado en la distribución desigual de los factores productivos se reduciría. Esta aparente divergencia entre teoría económica y realidad empírica ha sido explicada por diferentes motivos. Según algunos autores, gran parte de este incremento del comercio entre países industrializados fue una mera vuelta a relaciones económicas más normalizadas siguiendo las reducciones de las restricciones al comercio y los pagos internacionales introducidas en los años 30 y durante el principio de la posguerra. 
    
    Además, parte de la expansión del comercio entre estos países podría estar simplemente reflejando las mayores tasas de crecimiento económico alcanzadas por los mismos y el hecho de que países con niveles similares de renta son más proclives a comerciar más intensivamente entre ellos (dado que los patrones de demanda doméstica están influenciados hasta cierto punto por la renta). Finalmente, la mayor renta per cápita fomentó una mayor diversificación de la demanda de los consumidores, lo que, unido a avances tecnológicos rápidos y desiguales (que generan un proceso continuo de ajuste en las ventajas comparativas de los países), posibilitó un patrón creciente y más complejo de comercio de mercancías entre los países industrializados.
    
    Si alguna vez hubo una época dorada de la globalización, fue esta. Salvo por una cosa rara: las políticas del GATT no apuntaban precisamente a la globalización. Como hemos visto, la globalización requiere una reducción significativa de los costes de transacción en el comercio y las finanzas internacionales. Y así ocurrió en determinadas áreas. 
\end{specialblock}

\paragraph{Movimiento internacional de factores: atando a la bestia en corto}
Por último, ¿qué podemos decir de los flujos internacionales de factores de producción?

Los flujos de inversión directa extranjera (IDE) se incrementaron rápidamente hasta aproximadamente los 20.000 millones de dólares estadounidenses por año a principios de los 70. 

Gran parte de este notable aumento se debió a las operaciones de inversión realizadas por empresas multinacionales estadounidenses en territorio europeo para evitar el arancel exterior común de la CEE. En Europa, el Reino Unido se mantuvo como el principal emisor y receptor de IED, y, mientras que en los años 50 sus flujos de inversión tuvieron como principal destino países de la Commonwealth, en los 60 el foco de los mismos se situó en el continente europeo. 

En cuanto a flujos de trabajadores, cabe destacar que desde 1945 a 1960 la mayor fuente de migrantes fue Europa, de la que algo menos de 5 millones de personas emigraron durante los años 50 (cifra significativamente menor que los totales de cada década desde 1880 a 1910). Los principales países de destino fueron Estados Unidos, Canadá, Australia (cada uno de estos países absorbiendo más de un millón de personas), Argentina y Brasil, a pesar de que existieron durante este período importantes restricciones a la entrada de migrantes en Estados Unidos y otros países receptores. A diferencia de los emigrantes europeos del siglo XIX, que contribuyeron a expandir la oferta de alimentos y materias primas para el mercado europeo, el grueso de los inmigrantes de los años 50 entró en industrias manufactureras y terciarias en los países de destino (algunos de los cuales eran competidores de las industrias exportadoras europeas). La llegada de trabajadores extranjeros también contribuyó a mantener niveles salariales más estables y tasas de inflación más reducidas en los países de destino. Tras alcanzar unos niveles razonablemente elevados y estables de empleo en la mayor parte de países desarrollados, la atención se dirigió hacia las posibles formas de incrementar la capacidad productiva de las economías, es decir, aumentar el nivel de output alcanzable cuando todos los factores de producción son plenamente utilizados a través de cambios en la cantidad y calidad del stock de capital físico y humano.

A partir de 1960, las migraciones intercontinentales se redujeron en cierta medida, y las altas tasas de crecimiento económico en los países de Europa Occidental fomentaron la demanda de trabajadores de forma tan rápida que millones de trabajadores extranjeros de países cercanos fueron absorbidos por esta región. La relajación de restricciones al movimiento de trabajadores dentro de la Comunidad Europea también contribuyó al desplazamiento general de trabajadores existente dentro de Europa.

En consecuencia, la canalización de ayuda de los países desarrollados hacia las economías en desarrollo representó un cambio notable en el ámbito de los flujos internacionales de capital en comparación con el período anterior a 1939, y el conjunto de las economías en desarrollo experimentaron elevadas tasas de crecimiento (alrededor del 6\% anual promedio entre 1960 y 1973).


\subsection{El fin de Bretton Woods: morir de éxito}
Y llegamos al fin del periodo que, por constumbre o convención, suele marcarse/definirse con una fecha concreta. La realidad no es esa, pues el Sistema de Bretton Woods permanece, hoy en día, perfectamente visible. Tanto desde el plano comercial como desde plano financiero, tanto las instituciones que se constituyeron como los objetivos que se les atribuyeron se mantienen intactas (con ligeras modificaciones) hasta la actualidad. También se mantiene el Fondo Monetario Internacional, aunque con unos objetivos distintos a los que inicialmente se previeron.

Entonces, ¿qué colapso realmente? ¿por qué hablamos del colapso de Bretton Woods? Lo que verdaderamente colapso fue el acervo del Sistema de Bretton Woods. Las intenciones, objetivos y deseos de aquellos que lo formularon. La herencia que, durante casi tres décadas, dejó al mundo la mayor época de crecimiento económico experimentada hasta la fecha. El Sistema de Bretton Woods murió de éxito. 

Sin embargo, tradicionalmente se ha dicho que el sistema nació muerto. ¿Por qué? La razón es que el Sistema de Bretton Woods es visto, por muchos, únicamente como un Sistema Monetario Internacional. Aunque sabemos que esto no es así, veamos rápidamente por qué se dice que el colapso del Sistema se dió en 1971 (aunque, más formalmente, enre diciembre de 1971 y junio de 1976: sí, fue una muerte agonizante). 

\subsubsection{El colapso del Sistema Monetario Internacional: 1971-1976}
El problema que planteaba el pilar monetario de Bretton Woods era evidente. Tan evidente, de hecho, que no hubiese ocurrido si se hubiese implementado la visión keynesiana del Sistema, pues era, probablemente, el que mejor entendía lo que iba a ser realmente \textbf{Bretton Woods}. 

\paragraph{Paradoja de TRIFFIN: ¿un final evitable?}
Imaginemos la economía internacional a 1944. Podemos suponer, sin mucha imaginación, que estaba completamente devastada y el dinamismo económico era, cuanto menos, reducido. De hecho, podemos simplificar al máximo esta situación y considerar una foto fija de la economía mundial en 1944 (31 de diciembre, por ejemplo) plasmada en la ecuación cuantitativa del dinero:

\eqblock{
\begin{aligned}
    &\text{Problema}:&&
    \left\{\begin{aligned}
        &1)\,\,\text{Crecimiento Real}: &&\left(\frac{M}{P}\right)_0=\bar{v}\,Y_0&&\overset{\text{20 años}}{\Longrightarrow}&& \text{si},\,Y_1=2\,Y_0\,\,\Rightarrow\left(\frac{M}{P}\right)_0=2\left(\frac{M}{P}\right)_0\,\,\\[1em]
        &2)\,\,\text{Ancla nominal}: &&35\$_0=1\text{oz.}&&\Longrightarrow&&35\$_1=1\text{oz.}
    \end{aligned}\right.\\[2em]
    &\text{Solución}:&&35\$_0\equiv 1\text{oz.}\quad\overset{\frac{P_1}{P_0}=2.33}{\Longrightarrow}\quad 82\$_1=1\text{oz.}\\[1em]
\end{aligned}
}{Paradoja de TRIFFINF: Intuición económica. Colapso del Sistema Monetario Internacional}

Si, como hemos dicho, durante los siguentes veinte años el crecimiento económico real fue absolutamente desorbitante, una cosa está clara, las necesidades de liquidez para satisfacer dicho dinamismo solo podían aumentar. Pensémoslo para hacerlo más intuitivo. Imaginemos que todo el aparato industrial de Alemania se recupera y empieza un ciclo exportador sin precedentes. Los estadounidenses compararán marcos alemanes y venderán dólares, que serán transferidos a Alemania. Ante estasituación, cada vez más importante, surge un dilema inexorable: o repatriamos los dólares a la economía estadounidense a cambio de oro, o las autoridades monetarias amplían la oferta de dinero.\footnote{%
    Si una persona nos dice: \textit{Espera un momento, ¡no estás teniendo en cuenta la oferta de oro!}; lo más gentil que podemos llamarle es estúpido. Pero no pocos de nuestros compañeros economistas dirían algo semejante. 

En esta linea argumentativa subyace la idea de que el dinero es, necesariamente, fiduciario. Que hay algo real que tiene que respaldar su valor puesto que sino no cumpliría las funciones que la ciencia económica le atribuyen (medio de intercambio, depósito de valor y unidad de cuenta). Aunque comprensible, está completamente al margen de los resultados observados en inumerables estudios antropológicos que demuestran que, el dinero tal y como lo concebimos hoy en día, es en realidad dinero crédito. No es una mercancia como tal sino un facilitador. No obstante, para evitar la discusión podríamos decir, como resulta evidente, las impresoras de la FED van más rápido que los pobres mineros sudafricanos.
} Sin llevar a cabo ninguna de las dos acciones, Estados Unidos se quedaría sin dólares. O, como se conoce formalmente, la Paradoja de TRIFFIN: la liquidez necesaria para sostener el crecimiento observado ponía en riesgo el mismísimo Sistema. En otras palabras: morir de éxito.

Para respetar la regla de convertibilidad en oro, EEUU podía emitir dólares solamente en la medida en que acumulase reservas en ese metal, lo cual se traducía en una liquidez adicional limitada. Para proporcionar al resto del mundo el aumento deseado de reservas internacionales de dólares, EE. UU. debía generar continuos déficits por cuenta corriente, deteriorando su solvencia. Además, al incrementar sus pasivos en dólares en cuantía no correspondida por sus activos en oro introducía elementos de desconfianza en el sistema9. El dilema de Triffin se traducía por tanto en la imposibilidad del dólar de cumplir simultáneamente las funciones de moneda de circulación y activo de reserva con un tipo fijo ligado al oro.


\paragraph{Flujos internacionales de capital: un final ¿inevitable?}
Ahora bien, sería muy simplista considerar que este Sistema estaba abocado al fracaso. Que el final era inevitable. Lo cierto, y curioso curioso es que, durante 20 años, el Sistema funcionó perfectamente a pesar de que no uno, sino las dos acciones fueron realizándose paralelamente: la Reserva Federal llevaba a cabo Políticas Monetarias expansivas y los Bancos Centrales de países superavitarios recompraban oro a cambio de dólares americanos. Por tanto, falta un ingrediente más para entender el colapso del Sistema Monetario de Bretton Woods.\footnote{%
    factores que facilitaron su funcionamiento: - la firma de un Acuerdo General de Préstamos en 1962, entre los Bancos Centrales de 10 países, poniendo a disposición del FMI un mayor volumen de recursos financieros. Este acuerdo se renueva cada cuatro años. Si cuando un país solicita divisas al FMI, este no dispone de ellas, consulta al G-10 que decide si concederlas o no. Si las conceden, se reparten entre ellos. A cambio el FMI les entrega títulos no negociables que se reembolsan, junto con los intereses, a los cinco años; - la utilización activa de un mayor número de monedas como monedas de reserva (como el marco alemán o el franco suizo); - acuerdos bilaterales de crédito entre banco centrales (swap de divisas) para frenar movimientos especulativos, y - la creación de una nueva fuente de liquidez internacional, los Derechos Especiales de Giro (DEG) en 1967, en la primera enmienda al Convenio Constitutivo del FMI, con lo que se facultó al FMI a crear liquidez.
}

Para encontrar este ingrediente, debemos saber primero cómo se puede mantener la estabilidad del Sistema. Bajo un Sistema de este tipo, la única estabilidad a largo plazo viene dada por la coordinación internacional estricta. ¿Por qué? Si nuestra moneda está sobrevalorada respecto a su ancla nominal y, progresivamente, también respecto a otras divisas, el déficit no dejará de crecer: nuestra Balanza Comercial empeorará (inflación, sobrevaluación) progresivamente porque nuestros bienes será relativamente más caros; por otro lado, nuestra Posición de Inversión Internacional cada vez tendrá más obligacionespara con el exterior (influjos de capital, Estados Unidos se financia con entrada). Esto cuadra el Sector Exterior pero menoscaba la credibilidad del Sistema.

Entonces, ¿qué soluciones hay? La primera solución, y más obvia, hubiese sido una depreciación respecto al ancla: el oro. No obstante, esto hubiese sido difilmente justificable hacia nuestros socios, que ostentaban dichos derechos. Otra opción, más saludable para el sistema, sería el reajuste de paridades. ¿Reconocería este desequilibrio fundamental Estados Unidos? La verdad es que no, durantelargo tiempo, las autoridades supusieron que se trataba de un desajuste temporal, transitorio. Además, pensemos que los reajustes conllevaban pérdidas de credibilidad para con la comunidad internacional. 


Sin embargo, lo sentenció el Sistema fue mucho más básico: la especulación contra monedas, ya que un sistema de ajustes, de manera discontinua y por cuantías relativamente importantes, hacía evidente el sentido del ajuste y facilitaba que los especuladores pudieran anticipar esos movimientos corriendo riesgos mínimos. Por ello las devaluaciones se hicieron sin notificar al FMI y por sorpresa, resultando superiores que lo que habría sido necesario si se hubieran hecho en el momento apropiado. En concreto, el problema del ajuste generaba 3 tipos de asimetrías: - Asimetría del ajuste: al recaer sobre los países deudores, que no pueden mantener esa situación indefinidamente (pérdida de reservas).\footnote{%
    En previsión de esta asimetría se había incluido en los Acuerdos de BW una cláusula de recursos escasos, con medidas contra países permanentemente superavitarios, que nunca llegó a utilizarse.
} 



En primer lugar, porque en un Sistema con múltiples divisas, no pueden fijarse todos los Tipos de Cambio. De hecho, lo óptimo sería fijar únicamente un ancla nominal, el dólar. Si se opera de esta forma, el valor de la moneda reserva puede ser sostenida por el resto de Bancos Centrales resto del mundo. Así EEUU podía despreocuparse de sus problemas de balanza de pagos, interviniendo en los mercados los bancos centrales del resto de países para mantener la paridad de sus monedas respectivas. Aunque hubiera presiones depreciatorias contra el dólar, eran los demás países los que compraban dólares y vendían sus monedas, ya que no les interesaba que su activo de reserva (el \$) se devaluara. Esto proporcionaba a EE. UU. más margen de libertad en la conducción de su política económica. 





Los países se habían comprometido a hacer lo necesario para mantener su paridad con el dólar, lo cual implicaba la necesidad de adoptar políticas de modificación y desviación del gasto, con importantes costes para el país que las aplicaba. 





Asimetría del señoreaje 8: el único medio para conseguir medios de pago para el comercio internacional era generar superávit con el banquero (EE. UU.). Tener una balanza superavitaria frente a EE. UU. implica que se transfieren recursos reales y derechos sobre activos reales a EE. UU., y se reciben a cambio dólares, que suponemos que tienen un menor coste de obtención para EE. UU. (en principio, valdría con imprimir dólares). La emisión de dinero internacional genera una fuerte cantidad de ganancias por señoreaje.









\clearpage
\section{Conclusión} 


\subsection{Recapitulación}


\subsection{Extensiones y relación con otras partes del Temario}



\subsection{Opinión}

\subsubsection{Idea final (Salida o cierra)}

\clearpage
\pagenumbering{gobble}
\MyBibliography


\href{https://ocw.ehu.eus/pluginfile.php/41139/mod_resource/content/1/Tema_5/Aldcroft-Consecuencias_economicas_de_la_Gran_Guerra.pdf}{\textit{Consecuencias económicas de la primera guerra mundial}. UPV}

\href{https://www.funcas.es/wp-content/uploads/Migracion/Articulos/FUNCAS_PEE/017art14.pdf}{\textit{Keynes y la reforma del monetaria internacional} VARELA y MOLINA}

\href{https://www.funcas.es/wp-content/uploads/Migracion/Articulos/FUNCAS_PEE/017art14.pdf}{\textit{Las ideas de KEYNES para el ordén económico mundial} DÍEZ (2007)}


% ANEXOS
\clearpage
\anexo{\textbf{Preguntas Test}}
%\input{\TemaCode/questions.tex} 


\clearpage
\anexo{La Gran Guerra: Cronología de un desastre}\label{anx:gran_guerra}
\puntosclave{
Toda la información que aquí se expone ha sido extraída de las siguientes fuentes:
\begin{lnum}
    \item Blog: \href{https://historicodigital.com/la-gran-guerra.html#google_vignette}{\textit{La Gran Guerra}. Histórico Digital (2011)}
\end{lnum}
}

El inicio de la Gran Guerra en 1914 y los cuatros años de enfrentamientos bélicos entre los Estados europeos se convirtieron en la transición entre el mundo decimonónico y el siglo XX. La apacible vida de los europeos cambió tras el conflicto, así como \textbf{la forma de ver la vida}. Cabría destacar, para ejemplificar este cambio, algunos fragmentos de la insigne obra del austriaco ZWEIG, \textit{El mundo de ayer. Memorias de un europeo}:

\textit{Si busco una fórmula práctica para definir la época de antes de la Primera Guerra Mundial, la época en que crecía y me crié, confío en haber encontrado la más concisa al decir que fue la edad de oro de la seguridad. Todo en nuestra monarquía austríaca casi milenaria parecía asentarse sobre el fundamento de la duración, y el propio Estado parecía la garantía suprema de esta estabilidad. Los derechos que otorgaba a sus ciudadanos estaban garantizados por el Parlamento, representación del pueblo libremente elegida, y todos los deberes estaban exactamente delimitados. Nuestra moneda, la corona austríaca, circulaba en relucientes piezas de oro y garantizaba así su invariabilidad. Todo el mundo sabía cuánto tenía o cuánto le correspondía, qué le estaba permitido y qué prohibido. Todo tenía su norma, su medida y su peso determinados. Quien poseía una fortuna podía calcular exactamente el interés que le produciría al año; el funcionario el militar, por su lado, con toda seguridad podía encontrar en el calendario el año que se ascendería o se jubilaría. Cada familia tenía un presupuesto fijo [...] En aquel vasto imperio todo ocupaba su lugar, firme e inmutable y en el más alto de todos estaba el anciano emperador; y si éste se moría, se sabía que vendría otro y que nada cambiaría en el bien calculado orden.

Dicho sentimiento de seguridad era la posesión más deseable de millones de personas, el ideal común de vida. [...]

El siglo XIX, con su idealismo liberal, estaba convencido de ir por el camino recto e infalible hacia el mejor de los mundos. Se miraba con desprecio a las épocas anteriores, con sus guerras, hambrunas y revueltas, como a un tiempo en que la humanidad aún era menor de edad y no lo bastante ilustrada. Ahora, en cambio, superar definitivamente los últimos restos de maldad y violencia sólo era cuestión de unas décadas [...]
}

Todo esto, seguirá diciendo el autor, acabó con la Gran Guerra, y la que le continuará.

$$\text{\textbf{1. ¿Qué motivó la Guerra?}}$$
Hacía tiempo, desde 1870, que las potencias europeas esperaban una guerra, aunque tampoco podemos decir que no las hubiera habido a lo largo del periodo conocido como la paz armada (entre 1870 y 1914). Eran breves enfrentamientos, a veces en las zonas coloniales (lo que permitía ocultar que se tratara de una guerra entre Estados), o en su caso habían sido entre potencias de segundo orden, como en los Balcanes, en las que habían estado las grandes potencias detrás de todo ello, pero siempre en la sombra. Pese a todo, la diplomacia siempre estuvo por encima, y los conflictos entre Estados, que fueron muchos, se resolvían de forma diplomática, y las resoluciones aceptadas en teoría.

Pero ¿por qué estalló entonces la Gran Guerra? Precisamente porque todos esos conflictos fueron resueltos de forma puntual, aceptados a duras penas y sin que convencieran a casi nadie. Dieron lugar a un cúmulo de rencores que acabarían por aflorar en 1914. A este respecto, los rencorres fueron selectivos, puesto que entre los propios miembros de las alianzas existían motivos suficientes como para que la guerra hubiera sido un \textit{todos contra todos}. Austria-Hungría parece que había olvidado ya la derrota, a manos prusianas, de Sadowa (1866), al igual que las derrotas sufridas a lo largo de la unificación italiana. Por su parte, Francia pareció olvidar la derrota de Fachoda (1898) a manos inglesas, a la vez que los ingleses no parecieron tener en cuenta que Francia era su principal rival en la carrera colonial, sin mencionar que eran enemigas desde siempre. A Rusia no parece que le importara que sus principales aliados mantuvieran regímenes parlamentarios, mientras que el zarismo mantenía prácticamente un férreo absolutismo.

\imagenfit{Europa1914.png}{Distribución territorial europea: Preludio}{}

En cambio, otras rivalidades no estuvieron dispuestas a ser perdonadas. La principal de ellas era entre Alemania y Francia, puesto que hay que recordar que tras la derrota francesa en la guerra Franco-prusiana de 1870, Francia perdió los territorios de Alsacia y Lorena. Francia nunca lo olvidó, y la obsesión francesa, a partir de entonces, será la de recuperarlos. Mientras que Austria llevaba tiempo rivalizando con Rusia por el dominio de los Balcanes, que serán los que finalmente lleven al desencadenamiento de la guerra. Por su parte, Inglaterra, con el fin de mantener su amplio dominio colonial, debía intentar dejar a un lado a una Alemania que molestaba demasiado por su constante reivindicación de colonias, y por su creciente industrialización. Y es que al final, ante todo, se encontraba una causa económica, la de conseguir mercados y la preponderancia político-económica.

Pero eran cuestiones ante todo de prestigio. Gran Bretaña y Austria-Hungría querían conservar sus gloriosos imperios; Francia, Turquía y Rusia restaurar honores mancillados; y Alemania, Serbia e Italia buscaban victorias que reafirmasen la identidad de sus nuevos Estados. Se podría decir que, en general, el espíritu nacional de los países había crecido considerablemente, sobre todos en aquellos nuevos Estados que habían surgido no hacía mucho, como Alemania, o acababan de conseguir la independencia. A ello se le unió los nacionalismos que todavía esperaban su independencia, y que usaron la guerra para intentar conseguirla, como sucederá cuando ésta terminó.

$$\text{\textbf{1.1.} El desendadenante}$$
¿Por qué todo se desencadeno en 1914? Desde 1913, tras la Segunda Guerra Balcánica, Austria-Hungría había aceptado el tratado de Bucarest a regañadientes, y las malas relaciones con Serbia eran ahora ya totalmente patentes (hacía décadas que Austria intentaba su predominio en la zona, al igual que Rusia), que empeoraron, aún más, cuando Austria-Hungría inició en Bosnia maniobras militares, lo que hizo que el nacionalismo serbio se recrudeciera, apareciendo organizaciones serbias que no dudaron en usar la violencia contra todos los intereses austriacos. En ese contexto, el 28 de junio de 1914, el Archiduque Francisco Fernando, heredero al trono de los Habsburgo, se encontraba de visita oficial en la capital bosnia, Sarajevo, en donde un nacionalista serbio, Gavrilo Princip, le arrojó una bomba a su paso que le ocasionó la muerte.

Las diplomacias europeas se pusieron en marcha. En Viena, tras conocer el atentado, se sospechó que detrás de todo se encontraba el gobierno de Belgrado, así que se envió un ultimátum a Serbia en donde se les solicitaba que, en 48 horas, el gobierno condenara el atentando, iniciase una investigación para identificar a los terroristas, y se prohibiera toda actividad, por los grupos o asociaciones que radicaran en Serbia, en contra del Imperio Austro-Húngaro. Alguno de los términos del tratado eran los siguientes:

\textit{La historia de estos últimos años, y especialmente los acontecimientos de 28 de junio, han demostrado la existencia en Serbia de un movimiento subversivo cuyo fin es separar de la monarquía austrohúngara algunas partes de su territorio [...]. El gobierno real serbio debe comprometerse a:
\begin{lnum}
    \item Suprimir toda publicación que incite al odio y al menosprecio de la Monarquía; [...]
    \item Eliminar sin demora de la instrucción pública en Serbia [...] todo lo que sirva o pueda servir a fomentar la propaganda contra Austria-Hungría;
    \item Separar del servicio militar y de la administración a todos los oficios y funcionarios culpables de la propaganda contra la monarquía austrohúngara; [...]
    \item Abrir una encuesta judicial contra los participantes en el complot del 28 de junio que se encuentran en territorio serbio. 
\end{lnum}

El gobierno imperial y real espera la respuesta del Gobierno real lo más tarde hasta el sábado 25 de este mes, a las cinco de la tarde.

Del Gobierno de Austria-Hungría al Gobierno de Serbia, 23 de julio de 1914.
}

Evidentemente se sabía (y ese fue el objetivo austriaco) que el ultimátum ocasionaría la justificación para una guerra, puesto que, difícilmente, Serbia (apoyada por Rusia) aceptaría tales exigencias. Era una intromisión extranjera en el funcionamiento político serbio. De igual manera, se ha de suponer, el ultimátum fue enviado con el beneplácito alemán, que estaba dispuesto a toda costar a apoyar a Austria-Hungría.

Pese a ello, Austria envió a Rusia una notificación en la que se aseguraba que de ninguna manera existía una voluntad de anexionar territorio serbio (se trataba únicamente de un castigo), mientras que Rusia contestaba que se ampliara el plazo del ultimátum (quizás para ganar tiempo en la movilización de su ejército en una posible guerra), algo que Austria no aceptó, al tiempo que Rusia, sin tapujo alguno, comunicaba el apoyo a Serbia. Ésta, por su parte, intentó que mediara la diplomática europea, aceptando parte de los términos del ultimátum, mientras que el resto de las peticiones serían arbitradas por una conferencia integrada por los principales Estados europeos. Pronto, el canciller inglés, Sir Edward Grey, inició dicha misión, solicitando a ambos contendiente (Serbia y Austria-Hungría) que el problema fuera resulto por Gran Bretaña, Francia, Italia y Alemania, al considerarse que ninguna de estas potencias tenía intereses en los Balcanes. Rusia lo aceptó, pero no así Alemania, que propuso que Rusia se debía mantener neutral, y que se permitiera que Austria y Serbia resolvieran el conflicto entre ellas, lo que habría supuesto que Austria se impusiera a Serbia. Ni Francia, ni Inglaterra, ni Rusia lo validaron.

El 28 de julio de 1914, Austria-Hungría declaraba la guerra a Serbia, mientras que Guillermo II solicitaba al zar (Nicolás II) y a Inglaterra que no entraran en el conflicto. El Zar le volvió a proponer el arbitraje de naciones, pero de poco sirvió. Ni tampoco la nueva proposición de Inglaterra, por la cual Austria penetraría en territorio serbio a modo de castigo punitivo, pero se detendría en un lugar acordado, para después someterse al arbitrio internacional. Ninguna fórmula agradó a Alemania (dispuesta para la guerra), que envió un ultimátum a Rusia para que parara la movilización de su ejército iniciada días atrás.

A partir de aquí, fracasada la diplomacia, se pusieron en marcha las dos alianzas militares: la Entente y la Triple Alianza. El día antes de que las potencias se declararan la guerra las unas a los otras, \textit{La Tribuna} (periódico de la época), publicaba lo que pensaba cada Estado:
\textit{
\begin{lnum}
    \item Rusia dice: si no se respeta la integridad territorial de Serbia, intervendré contra Austria;
    \item Alemania dice: si Rusia pone un pie en Austria, apoyaré militarmente a esta;
    \item Inglaterra y Francia dicen: secundaremos a Rusia en su acción si interviene Alemania;
    \item Japón dice: enviaré dos escuadras al Mediterráneo y al Atlántico para apoyar a Inglaterra si se ve envuelta en una guerra;
    \item Rumania, Grecia y Montenegro dicen: apoyaremos a Serbia si se atenta contra su independencia;
    \item Bulgaria dice: apoyaré a Austria si intervienen Rumania y Grecia;
    \item Austria ha declarado que respetará la nacionalidad de Serbia, a la que solo quiere castigar;
    \item Italia secundará a sus aliados, Austria y Alemania, en caso de estallar el conflicto europeo;
    \item España permanecerá en todo caso neutral.
    \item Y mientras las naciones preparan sus ejércitos, se siguen celebrando en Viena, en San Petersburgo y en Berlín conferencias para que no se rompa la paz europea. 
\end{lnum}

La Tribuna, 31 de julio de 1914.
}

El 1 de agosto, Alemania declaró la guerra a Rusia. A su vez, al día siguiente, lo hizo Francia a Alemania. Y ésta lanzó un ultimátum a Bélgica para pasar por su territorio. El día 4, Gran Bretaña declaraba la guerra a Alemania, justificando la defensa de la neutralidad de Bélgica y Luxemburgo. Gran Bretaña no tenía los suficientes motivos, o al menos no tan fuertes, para entrar en guerra como Francia y Rusia, pero sabían que de ser derrotada Francia, Alemania obtendría el predominio sobre Europa, lo que no beneficiaría a Gran Bretaña. Especialmente porque los alemanes presionarían para obtener mayor número de colonias, y dominarían el plano internacional (tanto político, como económico). Por otra parte, no entrar en la guerra, y permitir que fuera ganada por Francia, supondría el mismo peligro que una Alemania victoriosa. Era, por tanto, mejor entrar en el juego bélico, y que los posibles beneficios de la guerra fueran repartidos.\footnote{%
    De todas formas, los ingleses, que -como todos- esperaban una guerra corta, consideraron que no sufrirían ningún tipo de desastre, puesto que, en principio, su apoyo sería básicamente económico (su aportación militar fue muy reducida en los primeros días), dando crédito a los franceses por valor de 3.000 millones de dólares. Acabaría, sin embargo, en ser también deudora de Estados Unidos, con un valor de 4.700 millones de dólares.
}

Unos y otros fueron declarándose la guerra mutuamente. Lo que había comenzado hacía apenas un mes, como un conflicto localizado en los Balcanes, afloró en una guerra europea (civil, como muchos autores sostienen). Era la primera gran guerra desde la época napoleónica, que implicaba a la gran mayoría de Europa divididos en dos bandos: la la Triple Entente, que pasarán a denominarse como los Aliados; y, la Triple Alianza, que se conocerá como los Imperios Centrales. 

A lo largo de la guerra entrarán en el conflicto el Imperio Otomano (por su rivalidad con Rusia) y Bulgaria (por su rivalidad con Serbia) del lado de los Imperios Centrales, mientras que Japón (que comerciaba con Inglaterra) y Estados Unidos lo harán del lado Aliado.\footnote{%
    También una serie de Estados independientes latinoamericanos, al menos de forma teórica (anecdótica básicamente), así como China y Siam lo hicieron del lado de estos últimos.
} Mientras, una Italia indecisa acabo saliendo de la Triple Alianza (sin que hubiera declarado la guerra todavía a nadie) y entrando en los Aliados. Se trataba, por tanto, de una verdadera guerra mundial, aunque el escenario será principalmente europeo, pese a que hay que recordar que se darán batallas en las colonias, así como en el mar.

Pocos fueron los Estados europeos que se mantuvieron neutrales: España, los Estados nórdicos, Holanda y Suiza. Aunque desde ellos, como en el caso de España, se podía operar para vender armas a los dos contendientes indistintamente.\footnote{%
    Se podría mencionar al español Juan March, que amasó una mayor riqueza en estos años de guerra. No sería el único, y, a lo largo del conflicto, unos cuantos consiguieron crear grandes fortunas a costa de ésta, lo que ocasiono, a posteriori, amplias críticas.
}

$$\text{\textbf{2. El inicio de la Guerra}}$$
¿Con qué entusiasmo recibieron la guerra los ciudadanos? Con el mejor posible. En una reacción de patriotismo, la gran mayoría de los ciudadanos de todos los Estados consideraron que la decisión de la guerra era lo mejor. Apoyaron, con vítores y aplausos, a sus respectivos gobiernos. Al fin y al cabo ¿Qué suponía esa guerra? Unas semanas de luchas (en la que se enviaría a los ejércitos profesionales, así como voluntarios), tras las cuales todos los Estados estaban seguros de que se alzarían con la victoria. Tras ello, todos podrían volver a casa, y la vida seguiría igual.\footnote{%
    No era extraño ese pensamiento. Europa estaba acostumbrada a unas guerras (las libradas a lo largo del siglo XIX) que no habían superada más de un par de meses, o apenas unas semanas. La última guerra, de larga duración, había sido en la época napoleónica, es decir, ningún europeo del momento había vivido una larga guerra. Pero este conflicto iba a ser diferente, iba a suponer el comienzo de la guerra moderna, por mucho que los Estados Mayores, formados en estrategias decimonónicas, intentaran continuar tácticas desfasadas para los nuevos armamentos. Pero el caso es que la guerra entusiasmo a todos, y le dieron tan poca importancia como la que refleja en su diario KAFKA: \textit{Alemania le ha declarado la guerra a Rusia; por la tarde me he ido a nadar}.
} Tan solo algunos intelectuales avisaron sobre lo que la guerra pudría suponer, posicionándose en contra de ella, al igual que lo hizo la Segunda Internacional, sin éxito alguno.\footnote{%
    El entusiasmo fue decayendo conforme se comprobó la crueldad de la guerra -no solo en los frentes, sino también en las propias calles de las capitales- y la longitud que ésta adquiría. Y respecto a esto último nos podríamos preguntar por qué se prolongó durante cuatro años. Por una parte, los contendientes eran países altamente industrializados, que habían innovado en su armamento a lo largo de casi un siglo (sobre todo a partir de 1870), y, especialmente, desde comienzos del siglo XX habían comenzado a destinar gran parte de los presupuestos nacionales al desarrollo armamentístico. Ello hizo, como ya se ha comentado, que la guerra tradicional de batallas decisivas se convirtiera en una guerra de desgaste. Sin embargo, la propia bonanza económica de la industrialización que permitió financiar la guerra, logró que ese desgaste se fuera alargando, ante unos recursos muchos más amplios para destinar a la guerra durante años. Y entre esos recursos se encuentra el factor humano. Europa podía contar ahora con miles de efectivos militares, en una población que se había multiplicado en el último siglo: de los 98,8 millones de habitantes a principios del XIX, en 1910 había nada menos que 355,5 millones de habitantes que podían ser enviados a los frentes.
}

$$\text{\textbf{2.1.} La Guerra rápida de los alemanes}$$
Con ese entusiasmo inicial, el Estado Mayor Alemán planeo que podrían celebrar la victoria en seis semanas. Diseñaron el plan Schlieffen, por el nombre del general que fue la cabeza pensante de éste, por el cual Alemania movilizaría sus tropas contra Francia, pasando por la neutral Bélgica, y acabarían en París, en donde Francia se rendiría. 

El paso por Bélgica (y Luxemburgo) era esencial para hacer más amplio el frente para que no se acumulara el potencial francés en un único sitio, es decir, en la frontera que unía Francia y Alemania. Ello daría tiempo para que Austria pudiera movilizar sus tropas contra los rusos, y, que además, cuando estos últimos hubieran movilizado las suyas -considerándose que un enorme imperio sin apenas infraestructuras, como el ferrocarril, tardaría semanas en traer sus ejércitos al oeste–, los alemanes, con Francia derrotada, moverían sus ejércitos hacía la frontera rusa, en donde derrotarían al Zar, y la guerra acabaría.

\imagenfit{GranGuerra.png}{Gran Guerra: Alianzas}{}

Pero todo salió mal para los alemanes. Tras penetrar en Bélgica, y encontrarse a pocos quilómetros de París, las tropas francesas y británicas -creyendo los alemanes que Inglaterra no entraría en la guerra– frenaron el avance alemán, dirigido por el general Helmuth von Molke. El general francés Joffre consiguió recuperar, en la batalla de Marne, buena parte del suelo francés ocupado. Fue una agotadora batalla, en la que se puso en marcha un armamento que ya no permitía un enfrentamiento cuerpo a cuerpo. Ambos bandos empezaron a escavar trincheras, y allí permanecieron prácticamente el resto de la guerra, desgastándose unos a otros, sin que se pudiera avanzar. A todo ello, se le unió que Rusia movilizó sus ejércitos antes de lo que se esperaba y Alemania tuvo que trasladar apresuradamente parte de su ejército del frente francés al frente ruso. 

La pesadilla alemana de enfrentarse en dos frentes se había cumplido.

$$\text{\textbf{2.2.} La Guerra de trincheras}$$

A comienzo de 1915, era evidente que la guerra había cambiado. Largos sistemas de trincheras a lo largo de los frentes llevó a que el conflicto se convirtiera en una guerra de posiciones, puesto que los frentes prácticamente no se movían. No había forma posible de lanzar un ataque que pudiera penetrar en las líneas enemigas, lo que conllevó la creación de nuevas armas que permitiera que el enemigo saliera de sus trincheras, como el lanzallamas, gases tóxicos, morteros, granadas, piezas de artillería, los primeros tanques, ametralladoras, la nitroglicerina (NOBEL, 1866) así como el avión que, de ser usado únicamente para la observación, pasó a convertirse en un elemento de ataque, bombardeando las trincheras. 

Pese a todo, no faltaron los intentos de lanzar una gran ofensiva a lo largo de los años de la guerra. Pero estos siempre acabaron con el desgaste de los contrincantes, que tras un gran esfuerzo, y miles de pérdidas humanas, no lograban apenas avanzar.\footnote{%
    Y es que, en el fondo, los bandos eran reacios a prescindir de la guerra tradicional. Las antiguas técnicas de combate se mezclaron con las nuevas, al igual que el nuevo armamento se usó junto con burros, mulas y caballos. E incluso algunos ejércitos fueron reticentes a cambiar los uniformes, de vivos colores, por otros que ocultaran al soldado en el territorio en donde se debía de esconder. Los franceses siguieron llevando guerreras azules, y pantalón rojo. Como se decía, \textit{era contrario al gusto francés} cambiarlos, y suponían la identidad del ejército, como alegaba el propio ministro de guerra, negándose a cambiarlos. Aunque es cierto que los alemanes sustituyeron el azul prusiano por el gris, y los ingleses tomaron el caqui.
} Pero no solo los enfrentamientos comenzaron a hacer estragos en la vida de los soldados. Las trincheras –estrechas, húmedas, frías¬– se convirtieran en auténticos nidos de enfermedades en donde los soldados mal vivían en el mejor de los casos, o directamente morían (especialmente los heridos).\footnote{%
    Una carta de un soldado francés, que recibían el nombre de Poilú por su aspecto, permite ver la situación:

\textit{Esos tres días pasados encogidos en la tierra, sin beber ni comer: los quejidos de los heridos, luego el ataque entre los boches alemanes y nosotros. Después, al fin, paran las quejas; y los obuses, que nos destrozan los nervios y nos apestan, no nos dan tregua alguna, y las terribles horas que se pasan con la máscara y las gafas en el rostro, ¡los ojos lloran y se escupe sangre!, después los oficiales que se van para siempre; noticias fúnebres que se transmiten de boca en boca en el agujero; y las órdenes boches dadas en voz alta a 50 metros de nosotros; todos de pie; luego el trabajo con el pico bajo las terribles balas y el horrible ta-ta-ta de las ametralladoras.

Verdún, marzo de 1916
}
}

Mientras, en el frente oriental (donde Paul von Hindenburg y Erich von Ludendorff labraron su reputación), los alemanes lanzaron una fuerte ofensiva contra los rusos, penetrando en su territorio 500 quilómetros (ya los rusos habían sufrido una gran derrota a finales de 1914 en Tannenberg), perdiéndose dos millones de vidas rusas. Es de entender que los franceses, temerosos de que Rusia cayera, lanzaran nuevas ofensivas en Artois y Campagne, con el fin de aliviar el frente ruso, aunque costaron unas 350.000 vidas francesas y 700.000 heridos. En sí, la penetración en el territorio ruso tampoco suponía una gran victoria alemana. Rusia, con el ejército más atrasado y con poco material bélico, siguió utilizando el sistema de guerra tradicional (pese a la capacidad del general Alexei Brusilov), y ayudados por el clima, prefirieron siempre que el invierno ruso les ayudara, así como sus malas comunicaciones que impedían que por mucho que los alemanes quisieran llegar más lejos, al final acababan por retroceder ante la imposibilidad de que los suministros llegaran hasta el frente, lo que hizo que se acabara por usar también el sistema de trincheras, pese a que existió mucha más movilidad entre 1915 y 1916 que en el frente occidental.

A lo largo de 1915, se abrieron nuevos frentes, al entrar nuevos contendientes. Italia, hasta entonces neutral, se pasó al bando Aliado, pues, ante el dominio británico del Mediterráneo, sabía que incluso neutral acabaría pagándolo más caro que entrando en la guerra.\footnote{%
    No obstante, la operabilidad y eficacia del ejército italiano fue puesto en duda en muchas ocasiones a lo largo de la Guerra. Queda para la historia la burla dirigida hacia ese ejército: \textit{El ejército italiano nunca se retira, solo se da la vuelta y sigue avanzando}. Aunque se desconoce si esa burla se refiere al comportamiento durante la I o II Guerra Mundial, la frase se atribuye al célebre comandante español Gonzalo Fernández de Córdoba (Gran Capitán), a las órdenes de los Reyes Católicos.
} Entraba también Bulgaria, del lado de los Imperios Centrales, así como el Imperio Otomano, que prontamente fue atacado por los Aliados con el fin dominar los Dardanelos y comunicarse con Rusia (\href{https://es.wikipedia.org/wiki/Batalla_de_Galípoli}{Batalla de Galipolí}), aunque se fracasó. El Imperio Otomano acabaría siendo uno de los países que más perdieron. Sin embargo, su entrada puso en peligro los interés británicos del sur asiático -por su ceracnía al Canal de Suez y su presencia (formal/informal) en la Península Arábiga, donde surgió la insurgencia arabe y una gran película: \href{https://es.wikipedia.org/wiki/Thomas_Edward_Lawrence}{Lawrence de Arabia})-, especialmente el suministro de petróleo que comenzaba a ser esencial.

$$\text{\textbf{2.3.} La Guerra de desgaste}$$
A partir de 1916, Alemania observaba impotente como su potencial militar no podía imponerse. Así que el Mando Alemán consideró que la única forma de romper el frente enemigo, sobre todo el occidental, era lanzando ofensivas a cualquier coste humano, con el objetivo de crear aún más bajas en el ejército enemigo, hasta que se rindieran. Era la guerra de desgate.\footnote{%
    De esta manera, en la batalla de Verdún (quizás la más famosa de esta guerra), a lo largo de cuatro meses, medio millón de soldados, entre alemanes y franceses, perdieron la vida en un duro enfrentamiento que no logró prácticamente cambiar el frente. Por su parte, los aliados lanzaron un ataque similar, la batalla de Somme, en donde, tras tres meses, se perdieron más vidas que en Verdún.
}

Por por aquel año, se dio el único combate naval de la Gran Guerra, pese a que los británicos esperaban un nuevo Trafalgar en donde demostrar que solo ellos eran dueños del mar: la batalla de Jutlandia. La victoria puede considerarse británica, aunque más bien deberíamos creer que quedó en tablas. Los alemanes, que desde hacía tiempo llevaban construyendo una gran armada, consiguieron hundir más barcos que los ingleses, aunque acabaron por retirarse con la flota a sus puertos, de donde no volvieron a sacarla. Quizás por ello se puede considerar que los ingleses vencieron, al no tener que volver a preocuparse de los mares, al menos en la superficie, puesto que los alemanes usaron el submarino para causar daños a los barcos aliados. También, en este año, los ingleses se enfrentaron al Imperio Otomano en Mesopotamia y Palestina. Y lo hicieron, de igual manera, los rusos, que lograron vencerlos en las batallas de Erzurum y Trebisonda, en el Cáucaso. Rusia también fue capaz de lanzar ofensivas a Austria-Hungría. Pese a todo, fueron victorias pírricas que agotaron a los rusos.

La guerra no parecía que fuera a acabar, y lo que es peor, en la retaguardia el desánimo cundía. La población civil empezaba a ver cómo en los mercados desaparecían los alimentos y los precios se encarecían, al tiempo que se necesitaba reclutar cada vez mayor número de hombres de forma forzosa. También subieron los impuestos, especialmente en Alemania y Austria-Hungría, puesto que les fue más difícil conseguir créditos en el exterior. Mientras que en los lugares ocupados, como Bélgica, muchas propiedades fueron destruidas y confiscadas. Era la llamada \textbf{guerra total} en la que todos los recursos del Estados iban encaminados a mantener el enfrentamiento bélico.\footnote{%
    No solo los hombres eran movilizados para el frente, sino que las mujeres, ahora, debían ocupar los puestos de trabajo de los combatientes, si se quería mantener la producción. Se trataba de que toda la población aportara algo a la guerra. El gobierno francés, alentaba de esta forma a la participación de las mujeres:

\textit{En nombre del gobierno de la República [...] os pido que mantengáis la actividad agrícola, recojáis las cosechas del último año y preparéis la del próximo [...]. Hay que atender a vuestra subsistencia, al abastecimiento de las poblaciones urbanas y, sobre todo, al de aquellos que defienden las fronteras, la independencia del País, la Civilización y el Derecho.} Llamamiento de Viviani, presidente del Consejo, agosto de 1914.

Es difícil determinar hasta qué punto esto despertó la conciencia a las mujeres para solicitar la igualdad con los hombres, pero es evidente que reforzó el movimiento feminista, y que muchos Estados dieron el derecho al voto durante los años veinte y treinta.
}

Por otra parte, los gobiernos no solo debían atender los asuntos militares, sino que tuvieron que adquirir nuevas funciones, que en el siglo XIX habrían sido impensables. Principalmente, tuvieron que iniciar políticas sociales ante el malestar y empobrecimiento de sus ciudadanos. Lo hicieron, ante todo, Francia e Inglaterra, en donde se puede apreciar una mejora de las condiciones de vida de sus habitantes, especialmente tras el conflicto. No lo hicieron así en los países con regímenes muchos más conservadores, como Alemania, Austria y Rusia, cuyos ciudadanos se empobrecieron. También los gobiernos tuvieron que encargarse de organizar la mano de obra, la producción y la racionalización de los alimentos, funciones que habían recaído, en el pasado, en manos privadas. Y, de igual manera, la propaganda se convirtió en el medio más eficaz para mantener la moral alta.

A la altura de 1917, los gobiernos eran consciente del desánimo civil y es que todos, en aquel momento, se preguntaron: ¿Por qué seguir luchando? ¿Por qué había empezado la guerra?, ¿Qué objetivos se perseguían? Nadie lo sabía, ni los propios gobiernos. Muchos fueron los que se plantearon que se debía llegar a un acuerdo entre todos los contendientes, y que la diplomacia volviera a reinar, como algunos estadistas norteamericanos recomendaron, donde destaca el coronel E. M. House. También el Vaticano, a la altura de 1917, intentó mediar en el conflicto. El papa Benedicto XV proponía un acuerdo de paz, alegando que se estaban enfrentando católicos contra católicos. La guerra afectaba, ante todo, a la autoridad del Papado en los asuntos religiosos, y muchos, preguntándose dónde estaba la mano de Dios, recelaban de la religión.

Pero ¿cómo se podía llegar a un acuerdo de paz si nadie sabía los objetivos? Si Alemania partió, en origen, con la idea poco definida de asegurar su propia protección, era evidente que controlaban un amplio territorio en 1916. ¿Objetivo cumplido? No, muchos dentro del ejército alemán querían más, y no querían escuchar ningún tipo de acuerdo de paz. (Incluso príncipes, intelectuales y políticos conservadores presionaban por la humillación total de Inglaterra y de Francia). Tampoco entre los sectores más conservadores de los aliados eran favorables a un acuerdo de paz, y sus objetivos no solo eran los de ganar la guerra, sino la de conseguir la mejor situación para sus propios Estados, algo que tampoco podían decir en alto. Solo los socialistas parecieron presionar para intentar una negociación de paz, quizás por ello luego ascendieron vertiginosamente. 

Pese a todo, Francia e Inglaterra lograron fijar una serie de objetivos generales de guerra, que fueron siendo modificados a lo largo de los acontecimientos.

$$\text{\textbf{2.4.} Rusia sale, Estados Unidos entra}$$
En 1917, el rumbo de la guerra dió un importante cambio, con dos trascendentales acontecimientos. 

El primero de ellos fue que Rusia salió de la guerra. Era el país más atrasado política y económicamente, y era el que más estragos había sufrido a lo largo de los casi ya tres años de enfrentamiento. Ante este clima, no es de extrañar que los movimientos internos de Rusia para acabar con el absolutismo zarista dieran sus frutos. La Revolución de febrero de 1917 acabó con la abdicación del Zar, y el establecimiento de un gobierno liberal, que consideró que debía seguir con la guerra. Ello no apaciguó el descontento de la población, que quería salir de ésta cuanto antes. 

No es de extrañar que un movimiento de corte marxista se hiciera poco después con el poder bajo la dirección de LENIN, que volvió a Rusia cruzando la propia Alemania con el apoyo del gobierno del Kaiser, al considerar que el triunfo de la revolución supondría la salida de la guerra de Rusia, lo que permitiría a los Imperios Centrales mover sus tropas, tranquilamente, hacia el frente occidental en donde podrían llevar a cabo una gran ofensiva. Y así fue, una vez conformado el Gobierno bolchevique ruso, LENIN firmó un armisticio con Alemania cesando los combates.\footnote{
Aunque la salida oficial no sería hasta pocos meses después, el 3 de marzo de 1918, cuando Alemania y el gobierno soviético firmaron la paz de Brest-Litovsk.

No todo el partido bolchevique estuvo de acuerdo en tal decisión, ya que, como contrapartida a la paz, Rusia debía abandonar los territorios de Finlandia, Polonia y Ucrania, que pasarían a ser independientes, aunque bajo la batuta alemana. Pero la decisión era inevitable si se quería utilizar el ejército para sofocar la propia guerra civil interna que se abrió entre la resistencia zarista y el gobierno soviético.
}

La salida de Rusia fue un alivio para Alemania. Y lo hubiera sido mucho más si meses antes no hubiera entrado otro contendiente en la guerra, que no había sufrido desgaste alguno, reforzando la posición de los Aliados. El 6 de abril de 1917, el Congreso de los Estados Unidos declaró la guerra a Alemania, al considerarse agraviados por estos. ¿Por qué? Desde 1915, utilizando una nueva arma, el submarino –mientras que sus barcos estaban atracados en los puertos–, los alemanes anunciaron que hundirían cualquier barco enemigo, así como barcos neutrales que comerciaran con el Reino Unido. Si estos intentaban bloquear el comercio marítimo alemán, también lo haría Alemania. Pero un importante comerciante con el Reino Unido era Estados Unidos, y cuando dos de sus barcos mercantes fueron hundidos, el Lusitania (1915) y el Sussex (1916), el pueblo y gobierno estadounidense, con la misma felicidad con la que lo habían hecho los europeos, entraron en la guerra del lado Aliado aportando dos millones de soldados.\footnote{%
    En la decisión, quizás, también contribuyó el telegrama enviado por Zimmermann (ministro de Asuntos Exteriores alemán) al gobierno mexicano, por el que se proponía una alianza para que México recuperara los territorios texanos si Estados Unidos entraba en la guerra. Pero ante todo era una decisión económica, Estados Unidos no podía permitir que Francia e Inglaterra fueran derrotadas, pues ambas tenían una deuda con los EE.UU de 8.700 millones de dólares.

De no haber sido por estos hechos (en conjunto), probablemente Estados Unidos nunca habrían entrado en el conflicto europeo, puesto que como alegaron la gran mayoría de los políticos estadounidenses en 1914, ellos no tenían ningún vínculo con Europa, ellos eran algo distinto, con su propia historia, y no tenían por qué intervenir en asuntos que no iban con el país.
}

$$\text{\textbf{2.5.} El final}$$

La entrada de Estados Unidos había fomentado que los alemanes forzaran la salida de Rusia, con la intención de reforzar el frente francés antes de que los Estados Unidos desembarcaran con sus fuerzas. Creían, una vez más, que podrían lanzar la ofensiva definitiva, que estuvo en manos de Ludendorff, quien creyó que en poco tiempo podría ocupar París. Fracasaron en el intento, y ante el gran número de bajas, y el reforzamiento aliado, el general francés Foch avanzó en el territorio ocupado por los alemanes.

En este momento, los Aliados contaban con tres importantes potencias. Sin embargo, el bando de los Imperios Centrales había acabado por convertirse prácticamente en Alemania. Ésta era la que llevaba todo el peso de la guerra. El resto de Estados, agotados, comenzaron a negociar la rendición por su propia cuenta. Bulgaria firmó el armisticio el 29 de septiembre, el 30 de octubre lo hizo el Imperio Otomano, y el 3 de noviembre de 1918 lo hacía la principal aliada de Alemania, Austria-Hungría (que se había encontrado con la muerte de Francisco José en 1916 y la presión de los nacionalistas).

Alemania estaba sola a finales de 1918 y, aunque el ejército aún era optimista (se consideraban un ejército ocupante) y los Aliados aún no habían puesto los pies en suelo alemán, la clase política comenzó las negociaciones presionando a Guillermo II para que acabara con la guerra. A ello se le unió motines revolucionarios, que acabaron con la abdicación de Guillermo II el 9 de noviembre, el mismo día en que en la localidad de Weimar se establecía la República, bajo el nombre de República de Weimar. La nueva República alemana firmaba dos días después el armisticio, y la consecuente finalización de la guerra, cuya negociación definitiva tuvo lugar en toda una serie de tratados a lo largo de 1919 que se convirtió en un año decisivo para Europa.

Sin tener en cuenta los tratados de Paz y la negociación política, que acabaría por reestructurar el mapa de Europa (caída de los imperios y traslado a Estados Unidos del liderazgo mundial); la guerra supuso una cuantiosa pérdida humana. 

La Gran Guerra se convirtió en la primera de las grandes masacres que azotarían Europa, algo no visto hasta el momento. Los muertos se contaron por millones. Nada menos que ocho millones de fallecidos -Alemania, 1,8 millones; Rusia, 1,7 millones; Francia, 1,385 millones; Austria-Hungría, 1,2 millones; Gran Bretaña, 947.000; y Estados unidos apenes 48.000 muertos–. A ello, siete millones de mutilados, que llenaron las calles de las capitales europeas, y que se convirtieron en el reflejo de la crueldad de la guerra y de la miseria. Deberíamos sumar también los muertos que murieron de hambre en las retaguardias debido a las carestías. Materialmente, las infraestructuras en los frentes quedaron desechas y las potencias europeas se habían endeudado considerablemente. Incluso Gran Bretaña, que empezó la guerra como el prestamista, acabó por estar fuertemente endeudada con los Estados Unidos.

Pese a todo, la población se recuperó a lo largo de los años veinte, y se produjo un nuevo resurgimiento del crecimiento económico. Fueron los felices años veinte, que acabarían en una Segunda Guerra Mundial, mucho más cruel que la primera.


\clearpage
\anexo{Crack del 29 y Gran Depresión: Una visión más profunda}

Las causas de la Gran Depresión en Estados Unidos en los años previos a la Segunda Guerra Mundial son un tema de debate activo entre los economistas, y forma parte de un debate de mayor magnitud acerca de la crisis económica, a pesar de que la creencia popular es que la Gran Depresión fue causada por el Crac del 29. Los eventos específicos en materia económica que tuvieron lugar durante la Gran Depresión han sido estudiados a fondo: una deflación en activo, y precios de mercancía, caídas dramáticas en la demanda y crédito, y desorganización del comercio, resultando finalmente en el crecimiento del desempleo y por lo tanto de la pobreza. Sin embargo, los historiadores carecen de consenso para determinar la relación causal entre diversos eventos y la política económica del gobierno como causa de la Depresión.
Se señalan tres motivos fundamentales que contribuyeron a generar la crisis.


$$\text{\textbf{1.- Teorías Generales:} causas y canales de la Gran Depresión}$$

\textbf{Subconsumo: Keynesianismo}. En primer lugar, existen teorías guiadas por la demanda, como el Keynesianismo y la economía institucional, estas, argumentan que la recesión fue causada por el subconsumo y la sobre-inversión causando así una burbuja económica que, al estallar, se vió agravada por la reacción de las autoridades. La lógica es simple: (i) niegan la estricta relación lineal que existe entre ahorro, tipos de interés e inversión; (ii) argumentan que las empresas invierten sobre la base de las expectativas y, si la caída en el consumo parece ser a largo plazo, los negocios reducirán las expectativas de ventas futuras -lo último que les interesa es invertir para incrementar la producción futura aunque las bajas tasas de interés hagan que el capital sea barato-; (iii) esta dinámica de autorefuerzo, donde las bancarrotas eran comunes y la inversión, se presentaba con muy poca frecuencia, solo podía solventarse con Políticas activas de demanda, la Reserva federal debía emitir moneda para que el gobierno federal tomara prestado y gastara, recortando a su vez los impuestos para estimular el consumo y reducir la carga.\footnote{%
    En Estados, el presidente Hoover decidió hacer lo opuesto a lo que KEYNES consideraba la solución. Hizo que el gobierno federal subiera excesivamente los impuestos para reducir el déficit presupuestario generado por la depresión y restringió el gasto del Gobierno. Desde la Política Monetaria, los tipos de referencia aumentaron para mantener el equilibrio de saldos reales en la economía.
} 

\textbf{Gran contracción: Monetarismo}. Por otro lado tendríamos la visión monetarista. La Gran Depresión comenzó como un fenómeno natural de ajuste ordinario, pero que errores significativos por parte de las autoridades en Política Monetaria provocaron una reducción en la oferta de dinero que agravó fuertemente la situación económica, haciendo que una recesión común se convirtiera en la Gran Depresión.\footnote{%
    Relacionados con esta explicación, se encuentran aquellos que señalan que la deflación de la deuda hizo que los prestatarios debieran más que nunca en términos reales.
} La lógica es simple, durante el periodo anterior, Estados Unidos atravesó varios ciclos de abundancia y escasez. Estas depresiones parece que fueron desencadenadas por pánicos bancarios, los más significativos ocurrieron en 1873, 1893, 1902, 1907 y 1920. Hasta la fecha, en la mayoría de estos pánicos, prestamistas de última instancia habían flexibilizado las condiciones de crédito y evitado el colapso del sistema, lo que provocó la institucionalización de dicho mecanismo a través de la creación de la Reserva Federal en 1913. ¿Qué fue diferente en la Gran Depresión? El pánico bancario, crisis de liquidez y las sucesivas quiebras bancarias, se acompañaron con una política activa de cuentas reales: la oferta de dinero debe ser respaldada activos reales, rehusando prestar dinero a los bancos y permitiendo un desempeño catastrófico que llevará a la quiebra absoluta. Para los monetaristas, los pánicos bancarios de 1931, 1932 y 1933 no habrían ocurrido si la FED hubiese actuado para proveer liquidez a los bancos que estaban quebrando.\footnote{%
    Las explicaciones monetaristas han sido rechazadas en la obra de SAMUELSON \textit{Economía}, donde escribió: \textit{Hoy en día pocos economistas observan a la Política Monetaria de la Reserva Federal como un remedio para controlar el ciclo de negocios. Los factores puramente monetarios son considerados tanto síntomas como causas}.
}

\imagenfit{Fig4.png}{La Gran Depresión desde un punto de vista monetarista}{\href{https://es.wikipedia.org/wiki/Causas_de_la_Gran_Depresión}{\textit{Gran Depresión}. Wikipedia.org}}

\textbf{Productividad: Economía marxista}. La visión austríaca está completamente llena de contradicciones, lo que acabó conduciendo al mismísimo HAYEK a aceptar la visión Monetarista de la Gran Depresión, lo que le supuso las críticas de parte de sus seguidores. Por tanto, presentaremos una versión alternativa de la Economía marxista. Desde este punto de vista, la Gran Depresión fue un aumento sostenido de la inversión en activos nominales, descuidando la inversión real y el aumento de la rpductividad del lustro precedente. \textcolor{red}{\textbf{Seguir aquí}}.


Otros autores, como LABINI, consideran que el gap entre la producción y el consumo se explica mas por un consumo insuficiente que por un exceso de producción. Labini considera que la causa de este consumo reducido está en el escaso crecimiento de los salarios en el periodo, muy por debajo del crecimiento de la productividad, con lo que no se generaba suficiente absorción. Contra esta explicación cabe argüir que el moderado incremento de los salarios permitió una fuerte reducción del desempleo, con consecuencias positivas sobre el consumo. El argumento del subconsumo, no obstante, puede claramente utilizarse en el caso de los productos agrarios. De acuerdo con la Ley de Engel, la demanda de estos productos debía aumentar menos que proporcionalmente a los aumentos de renta. Junto a esto, la fuerte demanda de bienes de consumo duradero, muy influida por la evolución de los tipos de interés, se va ver muy afectada por la subida de tipos a partir de 1928. 


$$\text{\textbf{2.- Teorías Específicas:} causas y canales específicos de la Gran Depresión}$$
Adoptando un enfoque más específico, algunos autores han querido enfatizar el papel que jugaría algún objeto económico concreto como catalizador y acelerador la situación observada.

\textbf{Deflacción de deuda: el Cíulo vicioso de la deuda}. La deuda total como ratio de PIB en Estados Unidos alcanzó una cima justo por debajo del 300\% en la época de la Gran Depresión Depresión, cifra que no sería alcanzada de nuevo hasta finales del siglo XX.\footnote{%
    JEROME (1934) proporciona una cita no atribuida acerca de las condiciones financieras que permitieron la gran expansión industrial del periodo posterior a la Primera Guerra Mundial: \textit{Probablemente nunca antes en este país se tuvo tal volumen de fondos disponibles a tasas tan bajas por un periodo tan largo}.

Los requerimientos marginales eran de solo el 10\%: las firmas de corretaje podían prestar \$90 por cada \$10 que un inversionista depositara.
} Con esto, es razonable suponer que determinados mercados estuvieran desviándose profundamente de su rendimiento natural. Por ejemplo, HANSEN analiza que la construcción de viviendas durante los años 20s excedió el crecimiento poblacional por 25\%. Le evolución del índice busrástil proporciona ilustra también la posible burbuja que se estaba generando en activos nominales. Independientemente del origen concreto de la burbuja (pues no el núcleo de lo que nos interesa), este nivel de sobreendeudamiento podría haber tenido consecuencias devastadoras, de acuerdo con esta teoría. La lógica es sencilla. Cuando el mercado cayó, destruyéndose las expetativas de los agentes, apresurándose a retirar sus depositos, los Bencos comerciales se vieron obligados a solicitar la devolución de préstamos inmediatos, que en su mayoría no podían ser liquidados con tanta brevedad, quebrando casi un millar de bancos en tan solo 10 meses desde el crack. ¿Pero estabamos tan mal? La cuestión, desde este punto de vista, es que no. Pero la desconfianza y temor era tal, que aunque las deudas podrían haber sido pagados, el sobredimensionamiento de estas debido a la destructiva espiral deflación de precios/disminución de la oferta monetaria, hizo imposible su cumplimiento. FISHER proporciona una buena visión de la dinámica: (i) al verse obligados a deshacer sus posiciones para cubrir la retirada de depósitos, se experimentó liquidaciones de deuda y ventas forzosas; (ii) esto, que no se veía acompañado por la FED, provocó una reducción de la oferta de dinero vía depósitos ($M_2$, base amplia) y vía balance (deterioro de las deudas incobrables); (iii) esta reducción de la oferta de dinero provocó una progresiva y sostenida caída en los niveles de precios, tanto de consumo como de los activos nominales; (iv) lo que afecta doblemente al valor de las empresas, caída de las ganancias y pérdida de valor de su colateral, multiplicando el impacto en el valor total de las empresas y precipitando bancarrotas; (v) lo que generaría una reducción en la producción, el comercio y el empleo (echando más leña al fuego); (vii) exacerbaría el pesimismo y la pérdida de confianza; (viii) aumentando la acumulación física de dinero, disminuyendo más aún la oferta de dinero $M_2$; y, (ix) provocando una caída en las tasas de interés nominal, pero en realidad una subida en las tasas de interés real, ajustadas por la deflación. Y, lo paradójico de todo esto es que, aunque las deudas pendientes se podían afrontar (al principio), a pesar de hacerse más fuertes en términos reales, ¡cuántas más saldabas tu deuda, más debías realmente! En esencia, se desarrolló un círculo vicioso y el espiral descendente se aceleró. La liquidación de la deuda no podía mantenerse con la caída de precios que esta causó, los bancos construyeron sus reservas de capital y realizaron menos préstamos, lo cual intensificó las presiones deflacionarias, el efecto masivo de la estampida a liquidar incremento el valor de cada dólar que se debía y el gran esfuerzo de los individuos para disminuir su carga de deuda efectivamente la incrementó.\footnote{%
    Algunos macroeconomistas incluyendo a BERNANKE, el expresidente del Sistema de la Reserva Federal, han revivido el punto de vista de deflación de deuda de la Gran Depresión original por FISHER. El economista KEEN revivió la Teoría del Restablecimiento de la Deuda después de predecir acertadamente la recesión del 2008 basándose en su análisis de la gran depresión, y recientemente aconsejó al Congreso participar en la condonación de deuda o pagos directos a los ciudadanos con el fin de evitar futuros acontecimientos financieros.
}


\textbf{Patrón oro: Un compromiso poco creíble}. Otra rama de la literatura ha justificado el origen de la Gran Depresión y, principalmente, su proyección internacional, en la inadecuación de las reglas antiguas del patrón oro dentro del contexto posterior a la I Guerra Mundial. De acuerdo con esta linea, la Gran Depresión fue causada por la decisión de la mayoría de las naciones de regresar al patrón oro a la paridad anterior. Hay que ser muy cautos al analizar este argumento puesto que, como sabemos, los países \textbf{no volivieron} al patrón oro. Nada más lejos de la verdad ya que, durante el primer lustro de los años 20 se adoptaron una serie de erráticas políticas que generon persistencia en la dificultad del ajuste interno. De esta forma, el Sistema Monetario Internacional inmediatamente posterior se caracteriza por un \textbf{no-sistema}, en el que las devaluación y las suspensiones a la convertibilidad fueron la norma. ¿Cuándo, entonces, se restablece la normalidad e identifican el surgimiento del problema? Todo tiene sus raíces en el Plan DAWE y el buen rendimiento económico que trajo consigo. Dicha estabilidad real y fortalecimiento de las expetativas nacionales, permitió el restablecimiento del Sistema Monetario Internacional. Ahora bien, ¿bajo qué regimen? A menudo se dice que el régimen adoptado fue el \textbf{Patrón Cambios Oro}, pero la realidad fue mucho más heterógenea que la observada en el periodo precedente.\footnote{%
    Este régimen fue inicialmente tratado en la Conferencia de Génova, en 1922, donde se recomendó que se consideraran reservas cualquier divisa que fuese convertible en oro, para evitar que la escasez de oro mundial limitara la oferta monetaria. (Piénsese que, en realidad, se estaba introduciendo de esta forma un nuevo multiplicador de la oferta monetaria: Gran Bretaña oferta libras respaldas por oro, Francia oferta francos respaldados por libras respaldadas por oro). De hecho, esta variedad del patrón oro no suponía un cambio sustancial en relación con la situación prebélica, ya que muchos países habían utilizado la libra esterlinade la misma forma en el periodo preguerra. 

No obstante, el funcionamiento teórico puro del sistema era aún más sencillo: se nombraría un banquero al que se transfiere todo el oro y este extiende en contrapartida pasivos monetarios, que el resto de países acceden a mantener como medios de pago. Por tanto, se crean dos roles claramente diferenciados: (i) países ancla (Estados Unidos y Gran Bretaña), que se comprometen a mantener como única reserva el oro; y, (ii) resto de países, que pueden mantener como reserva moneda convertible y emitir sus pasivos repaldados por esta. ¿Coherente y sencillo, verdad? Este Sistema, tal y como se formula, no fue adoptado por ningún país. Al final y al cabo, ¿quién renunciaría a sus reservas en favor de un adversario? No obstante, triunfaron algunas variantes que permitiron la estabilidad.
} 

Estados Unidos adoptó el Patrón Cambios Oro pues era el que más beneficioso le resultaba. El Reino Unido optó por la estabilidad nominal de la economía y adoptó el patrón lingotes oro: restableció la paridad en 1925 pero sólo se permitiría la conversión a partir de una cierta cantidad de moneda. Por lo que respecta al resto de países, podemos decir, \textit{grosso modo}, que aceptaron mantener sus pasivos respaldados por una combinación entre divisa (reserva de valor) y oro (el poco que les quedaba), lo que permitiría: (i) utilizar tanto el dólar como la libra como monedad de reserva; y, (ii) restaurar, con sus más y menos, la estabilidad del Sistema Monetario Internacional. Y así fue, el sistema funcionó.\footnote{%
    Muchos autores sostienen que este sistema fue disfuncional. No es cierto. Recordemos que los \textbf{objetivos declarados} del Sistema Monetario Internacional eran, únicamente, la \textbf{estabilidad cambiaria}. Bajo este pretexto, no hay duda de que el Sistema funcionó, pues aumentó notablemente los movimientos internacionales de capital y construyó una sólida base de expectativas sobre los términos de la paridad.

Ahora bien: eso no quiere decir ni que fuera óptima ni que fuera fácil. Aunque los focos de inestabilidad acuciante parecían disiparse o reducir su importancia -inflación, deflación, reparaciones, etc.-, se habían producido cambios estructurales díficiles de sortear. En particular, las economías europeas mantuvieron una lucha frenético de tipos de interés para mantener su equilibrio exterior. El problema fue significativamente intenso entre Francia y Gran Bretaña. La economía Inglesa fijó definitivamente la paridad (1925) en los niveles de preguerra, mientras que Francia volvió con una paridad significativamente inferior. Las consecuencias directas para los Ingleses fueron perfectamente predecibles: una pérdida brutal de competitividad (moneda muy apreciada) que drenaba sus reservas (obligando a mantener altos tipos de interés) y agrababa, directa (competitividad) e indirectamente (tipos e inversión), el desempleo estructural y los conflictos sociales que esto generaba. Por lo que respecta a Alemania, su elevado grado de destrucción canalizó ingentes recursos financieros, por lo que la estabilidad monetaria sirvió en gran medida para la reconstrucción del país.

En estas condiciones, el patrón-oro, sirvió para la estabilidad monetaria exterior, estimulando los flujos internacionales de capital, pero no ayudó mucho a la recuperación real de su economía.
}

Entonces, ¿qué pasó para que este efecto positivo pueda ser considerada una posible causa? El Sistema fucionó porque fue sutilmente sometido a los objetivos de política interna (a excepción del Reino Unido).\footnote{%
    El férreo compromiso de las instituciones inglesas era una anomalía. Debe quedar muy claro. El comportamiento de los Estados Unidos fue muy distinto (evidentemente respaldado por el buen rendimiento económico). Estos no consideraban, ni mucho menos, como prioridad el rol que ostentada como garantes del Sistema Monetario Internacional, y en consecuencia no hiceron cargar sobre su capacidad de crédito interno las exigencias de paridad.
} El caso paradigmático lo encontramos en Estados Unidos: bajo el mandato de STRONG, las condiciones de crédito interno y la Política Monetaria mantenida por la FED perseguía como único objetivo la estabilidad de precios, la posición en balance... bueno, era una asunto secundario. Ahora bien, tras la muerte de STRONG y el cambio de rumbo de la Política de la FED, que ahora priorizaba una oferta monetaria respaldada por saldos reales, la cosa era muy diferente. La Oferta Monetaria superaba lo aceptable y, desde prácticamente sus inicios, la nueva Junta de Gobierno aumentó los tipos de interés. Aquí tenemos el primer detonante, si hasta ahora Estados Unidos había sido permisivo con la Política Monetaria de sus compañeros europeos (que mantenían una estrecha guerra entre ellos), la cosa iba a cambiar. 

Cuando la FED subió sus tipos de interés de referencia a principios de 1928 para detener lo que pensaba que era un exceso especulativo en Wall Street, ejerció más presión sobre países con déficit exterior, como Gran Bretaña. Los tipos de interés más altos en Estados Unidos obligaron a esos países bien a seguir el juego y subir sus propios tipos de interés, bien a sufrir más hemorragia de oro y capital. Y este juego podría haber funcionado, pero en 1929 recibió un golpe y contragolpe que sería incapaz de sostener. Para evitar repetición, simplemente situese aquí el argumento monetarista sobre la Política de Tipos y apliquese la lógica del canal exterior para entender lo que sucedió: Tipos, tIpos, tipoS y más tipos...\footnote{%
    TIMBERLAKE, economista de la escuela de banca libre y protegido de FRIEDMAN, específicamente abordó esta instancia en su publicación \textit{El patrón oro y la doctrina de cuentas reales en la política monetaria de Estados Unidos}. En esta destaca libertad de acción que había reinado en la FED durante la década, pese a su adhesión al patrón oro. Bajo la década de los años 20, argumenta, el patrón oro permaneció nada más que como una fachada formal esperando el momento oportuno para reaparecer. Sin embargo, cuando la facción que tomó el dominio de la FED y promovió una doctrina de Cuentas Reales, se dañó inevitablemente la economía al forzar una deflación del 30\% en el dólar.
} Aunque aguantó, en septiembre de 1931 Gran Bretaña se salió del oro una vez más. En cuanto este abanderado del oro se vio obligado a salir, los días de este régimen estaban contados.

Otras investigaciones más reciente pero en esta misma linea (TEMIN, BERNANKE y EICHENGREEN) también contrastan esta idea, enfocadose en cómo las limitaciones de la guerra interina del patrón oro magnificaron la crisis económica y fueron un obstáculo significativo a cualquier acción que habría corregido la espiral negativa. De acuerdo con ellos, la inicial crisis desestabilizadora se debió de haber originado con el Crack en Estados Unidos, pero fue el patrón oro lo que transmitió el problema al resto del mundo. Los gobiernos tenían las manos atadas a medida que la economía se colapsaba, a menos que abandonaran el vínculo de su moneda con el oro.


$$\text{\textbf{3.- Conclusión}}$$

Pese a parecer teorías estancas, en realidad todas las causas, efectos y consecuencias pueden, y deben, considerarse como complementarias. Los economistas, particularmente diría yo, tiene la extraña obsesión de buscar la trascendencia aportando visiones innovadoras y evitando recurrir a observaciones ya realizadas. No obstante, esto lleva a que nos alejemos de una visión del conjunto del fenómeno tratado. 

Todo lo expuesto puede contrastarse empíricamente y parece una linea de argumentación razonable, al menos para mi. Por tanto, intentar entender cada una de ellas sin adherirse plenamente a una descripción del suceso es una condición necesaria para entender las dinámicas que condujeron a este. Es perfectamente posible que algunas, más que otras, motivara o agravaran la situación, pero en Economía el producto no puede descomponerse por partes. El todo es resultado de un comportamiento extremadamente no lineal al que todos aportan algo y nada, en concreto, puede considerarse como principal.




La búsqueda de la estabilidad monetaria internacional una vez terminada la guerra no fue sencilla para los beligerantes. Debían, al mismo tiempo, dar estabilidad interna a sus monedas luchando contra la inflación y, por otro lado, alcanzar la estabilidad cambiaria, para lo cual se intentó volver al sistema de tipos de cambio fijos de preguerra: el patrón oro. Podrían analizarse muchos aspectos vinculados a los problemas de la estabilización monetaria (luchas contra la inflación, fallida vuelta al patrón oro, políticas monetarias contradictorias, movilidad de capitales), pero vamos a centrarnos en reseñar lo que nos parece el proceso más interesante para este post: la hiperinflación alemana.

En los primeros años de la posguerra se produjo un boom al intentar reponerse los stocks a los niveles previos a la contienda.\footnote{%
    Como señala COMÍN (2014), la guerra consolidó la tendencia a la concentración industrial, y la crisis económica de la inmediata posguerra fomentó los acuerdos monopolísticos entre empresas para controlar la oferta. Una vez que Europa alcanzó los niveles de producción de preguerra, la sobreproducción de algunas ramas de actividad comenzó a ser un problema para la economía mundial. Para 1919, los mercados globales estaban saturados.
} Esto permitió inicialmente absorber gran parte de los trabajadores desempleados y condujo a un aumento de la oferta mundial de materias primas, atraídas por el aumento de precios asociado al tirón de la demanda. Sin embargo, pasados los primeros años y cuando los stocks alcanzaron sus niveles habituales, se produjo el primer aumento generalizado del desempleo y un desplome de los precios mundiales de alimentos y productos básicos. 





















\end{document}